% Paper_2_Foundational_Gauge_Program_Supplement v1.0 (2026-07-13): NEW DOCUMENT -- TECHNICAL SUPPLEMENT II, THE FOUNDATIONAL GAUGE PROGRAM.
%   Created at the dual-supplement split (ruling of record, E. Brooke 2026-07-13). Carries the APF-to-gauge chain at its banked grades,
%   split out of Paper_2_Structure_of_Admissible_Physics_Supplement_v4.1 (whose classification core continues as Technical Supplement I v5.0):
%   superadditivity -> noncommutative algebra (standard C* semisimplicity route, review item 4.1.01.06; rotated 3-dim construction re-labeled
%   a worked example), gauge identification (Gau_0 subset-of Inn unconditional + equality under the named full-invariance hypothesis, item .03),
%   matter action defined on the lifted group with the charge representation a named input (item .04), blocks=generations passage DELETED
%   (item .05), C(G)=dim(G) as lower-bound theorem + composed operational statement on the named premise pair (item .07), CM stated at
%   channel-configuration level (item .08), R-EC-inv/I-typing material with the two banked EC checks (v24.3.423: check_L_F6_not_from_EC [P],
%   check_L_EC_inventory_reading [P_structural_reading | R-EC-inv + I-typing]), grade-ledger table, open-lane register headed by the
%   fixed-point program. Register: graded program with named opens.
\documentclass[11pt]{article}
\usepackage[letterpaper,margin=1in]{geometry}
\usepackage{amsmath,amssymb,amsthm,mathtools}
\usepackage{enumitem}
\usepackage{booktabs}
\usepackage{array}
\usepackage[colorlinks=true,linkcolor=black,citecolor=black,urlcolor=black]{hyperref}
\usepackage[most]{tcolorbox}
\tcbuselibrary{skins,breakable}
\usepackage{fancyhdr}
\usepackage{caption}

% ── Theorem environments ──
\newtheorem{theorem}{Theorem}[section]
\newtheorem{lemma}[theorem]{Lemma}
\newtheorem{proposition}[theorem]{Proposition}
\newtheorem{corollary}[theorem]{Corollary}
\newtheorem{theorem_app}{Theorem}[section]
\newtheorem{lemma_app}[theorem_app]{Lemma}
\theoremstyle{definition}
\newtheorem{definition}[theorem]{Definition}
\newtheorem{remark}[theorem]{Remark}
\newtheorem{assumption}[theorem]{Assumption}

% ── Box style (matches Paper 1 / Paper 2) ──
\tcbset{
    apfbox/.style={
        enhanced,
        colback=black!3,
        frame hidden,
        borderline={0.5pt}{0pt}{black},
        arc=0pt,
        left=8pt, right=8pt, top=8pt, bottom=8pt,
        fonttitle=\bfseries\color{black},
        coltitle=black,
        detach title,
        before upper={\tcbtitle\par\smallskip},
        before skip=14pt,
        after skip=14pt,
        grow to left by=-8pt,
        grow to right by=-8pt,
        sharp corners,
        breakable,
        pad after break=4pt,
        pad before break=4pt,
        lines before break=3,
    }
}

% ── Shortcuts ──
\newcommand{\eps}{\varepsilon}
\newcommand{\End}{\mathrm{End}}
\newcommand{\SU}{\mathrm{SU}}
\newcommand{\U}{\mathrm{U}}
\newcommand{\id}{\mathrm{id}}
\DeclareMathOperator{\spn}{span}
\DeclareMathOperator{\conv}{conv}
\DeclareMathOperator{\tr}{tr}
\DeclareMathOperator{\diag}{diag}
\newcommand{\coderef}[2]{\footnote{Code reference: \texttt{#1} in \texttt{#2}.}}
\newenvironment{varlist}{\begin{itemize}[nosep,leftmargin=2em]}{\end{itemize}}

\pagestyle{fancy}
\fancyhf{}
\fancyhead[L]{\small\itshape Paper 2 Technical Supplement II --- The Foundational Gauge Program}
\fancyhead[R]{\small\itshape E.\,S.\ Brooke}
\fancyfoot[C]{\thepage}
\renewcommand{\headrulewidth}{0.4pt}


% ═══════════════════════════════════════════════════════════════
\setlength{\emergencystretch}{3em}%CORPUSGUARD
\usepackage{seqsplit}%CORPUSGUARD
\usepackage{graphicx}%CORPUSGUARD
\newcommand{\codeid}[1]{\texttt{\seqsplit{#1}}}%CORPUSGUARD

%CORPUSFMT 2026-07-04: preamble normalized to the Paper 0 corpus standard
\usepackage[T1]{fontenc}
\usepackage[utf8]{inputenc}
\usepackage{lmodern}
\usepackage{parskip}
\usepackage{titlesec}
\usepackage{microtype}
\titleformat{\part}[display]{\centering\Huge\bfseries}{\partname~\thepart}{1em}{}
\titleformat{\section}{\Large\bfseries}{\thesection.}{0.5em}{}
\titleformat{\subsection}{\large\bfseries}{\thesubsection}{0.5em}{}
\titleformat{\subsubsection}{\normalsize\bfseries}{\thesubsubsection}{0.5em}{}
\setlength{\emergencystretch}{1.5em}

\newcommand{\TSII}{Technical Supplement~II (\emph{The Foundational Gauge Program})}
\newcommand{\TSI}{Technical Supplement~I (\emph{The Classification Core})}

\date{Version 1.0 --- July 2026}
\begin{document}

\begin{center}
{\LARGE\bfseries Paper 2 --- Technical Supplement II:\\[4pt]
The Foundational Gauge Program}

\bigskip
{\large E.\,S.\ Brooke}\\[2pt]
{\small ORCID: \href{https://orcid.org/0009-0001-2261-4682}{0009-0001-2261-4682}}\\[4pt]
{\small Companion to Paper~2: ``The Structure of Admissible Physics''\\
Technical Supplement II of two; the classification core is Technical Supplement~I}

\bigskip
{\small Version 1.0 --- July 2026}
\end{center}

\bigskip

\begin{abstract}
\noindent
This supplement carries the foundational program of Paper~2: the chain
from the APF axioms toward the inputs that the classification core
(\TSI) consumes, each link stated at its banked grade with its
premises named.  The chain runs: cost uniqueness
($L_{\mathrm{Cauchy}}$: $F(n)=n$ on the discrete channel domain); the
Sep/IJC-routed bridge from superadditivity to a finite noncommutative
admissibility algebra, semisimple by the standard finite-dimensional
$C^*$ route, derived given the IJC occupancy (which is read off the
empirical record); the gauge identification --- unconditionally, every
continuous physical gauge redundancy is inner,
$\mathrm{Gau}_0 \subseteq \mathrm{Inn}(\mathcal{A})$, with equality
under a named full-invariance hypothesis; the exact sequence
$1 \to Z(\mathrm{U}(\mathcal{A})) \to \mathrm{U}(\mathcal{A}) \to
\mathrm{Inn}(\mathcal{A}) \to 1$ with the relative-phase torus as the
sole abelian source; the matter action, defined on the lifted group
$\widetilde{G}$ with the charge representation a named input
(I-typing), the algebra side determining it only up to the central
torus; the lift and the global $\mathbb{Z}_6$ quotient; and the
$C(G)=\dim(G)$ cost program --- the lower bound a theorem on the
operational-encoding premise, the upper bound on the disjoint-anchor
realizability premise, the equality resting on that operational
premise pair pending a cost functor.  Enforcement completeness is
carried at its honest grade,
[P\textsubscript{structural\_reading} $\mid$ R-EC-inv $+$ I-typing],
with the counter-construction owned in the bank; capacity minimality
is A2 over the A1$+$MD admissible set.  The program's newest closures
are the two banked EC results: F6 and EC-label-injectivity are
logically independent over the declared scan space, all four quadrants
inhabited and the F6-free argmin at 42 Weyl degrees of
freedom\coderef{check\_L\_F6\_not\_from\_EC}{ec\_inventory\_reading.py},
and the one-U(1) necessity leg holds exactly under the kinematic
type-inventory reading with the entry-typed
inventory\coderef{check\_L\_EC\_inventory\_reading}{ec\_inventory\_reading.py}.
A grade-ledger table near the front states every link, its grade
string, its named premises, and its coderef.  The open-lane register is
headed by the fixed-point program: whether the (template, matter
content, single-abelian-grading) triple is the unique simultaneous
solution of the full constraint set.  The register throughout is a
graded program: theorems at their grades, named premises at their
consumption sites, opens named as opens.
\end{abstract}

\medskip
\begin{tcolorbox}[apfbox, title={Grade ledger: every chain link, its grade, its named premises, its coderef}]
\footnotesize
One row per link of the foundational chain.  Grade strings are the
bank's; a premise named in the third column is consumed at the cited
statement, not derived there.
\smallskip
\begin{center}
\renewcommand{\arraystretch}{1.18}
\resizebox{\textwidth}{!}{\begin{tabular}{@{}p{4.2cm}p{2.6cm}p{5.2cm}p{4.2cm}@{}}
\toprule
\textbf{Link (statement)} & \textbf{Grade} & \textbf{Named premises / readings} & \textbf{Coderef} \\
\midrule
$L_{\mathrm{Cauchy}}$: $F(n)=n$ on $\mathbb{N}$ (\S\ref{sec:cauchy_supp}) & [P] & C1 additivity on separated interfaces ($L_{\mathrm{loc}}$); unit normalization & \codeid{check\_L\_Cauchy\_uniqueness} (\codeid{L\_Cauchy\_uniqueness.py}) \\
Sep/IJC bridge: (IJC) $\Rightarrow$ noncommutative record algebra (\S\ref{sec:nc_summary}) & [P\textsubscript{math}+P\textsubscript{APF}] & IJC \emph{occupancy} empirical (read off the record); bridge derived & \codeid{check\_T\_IJC\_dichotomy}, \codeid{check\_T\_inseparable\_IJC} (\codeid{core.py}) \\
Superadditivity $\to$ finite semisimple carrier (Prop.~\ref{prop:superadd_to_semisimple_v35}) & [P] & A1, MD, $\Delta>0$ (IJC), positive $B$, adjoint closure, finite resolution; standard $C^*$ route & \codeid{check\_T\_alg} (\codeid{core.py}) \\
Gauge identification: $\mathrm{Gau}_0 \subseteq \mathrm{Inn}(\mathcal{A})$ (Thm.~\ref{thm:continuous_gauge_inner_v34}) & [P] & --- (unconditional inclusion) & \codeid{check\_T3} (\codeid{core.py}) \\
Gauge identification, equality leg & [P $\mid$ full-invariance] & full-invariance hypothesis: anchor set, cost kernel, continuation data $\mathrm{Inn}$-invariant as a family & \codeid{check\_T3} (\codeid{core.py}) \\
Exact sequence $1 \to Z \to \mathrm{U}(\mathcal{A}) \to \mathrm{Inn}(\mathcal{A}) \to 1$; $T_{\mathrm{rel}}$ sole abelian source (Prop.~\ref{prop:exact_seq}) & [P] & --- & \codeid{check\_T3} (\codeid{core.py}) \\
Matter action on $\widetilde{G}$; algebra determines it up to $T_{\mathrm{rel}}$ (Prop.~\ref{prop:gauge_acts_on_fields}) & [P $\mid$ I-typing] & charge representation = named input of the matter category (I-typing) & \codeid{check\_L\_EC\_inventory\_reading} (\codeid{ec\_inventory\_reading.py}) \\
Lift and global $\mathbb{Z}_6$ quotient (Props.~\ref{prop:PU_to_SU}, \ref{prop:global_structure}) & [P] & maximal-quotient convention (named); $T_{\mathrm{rel}}$ observability via charged matter & \codeid{check\_T\_gauge} (\codeid{gauge.py}) \\
EC (enforcement completeness) (Lem.~\ref{lem:EC_CM_from_A1}(i)) & [P\textsubscript{structural\_reading} $\mid$ R-EC-inv $+$ I-typing] & R-EC-inv (kinematic type-inventory reading; A1-motivated, not A1-derived); I-typing; counter-construction owned in-check & \codeid{check\_L\_EC\_inventory\_reading} (\codeid{ec\_inventory\_reading.py}) \\
CM (capacity minimality) (Lem.~\ref{lem:EC_CM_from_A1}(ii)) & [P] & A2 over the A1$+$MD admissible set; channel-configuration semantics & \codeid{check\_L\_cost\_gauge} (\codeid{core.py}) \\
$C(G) \ge \dim(G)\cdot\varepsilon^*$ (Lem.~\ref{lem:gen_channel}) & [P $\mid$ op.\ encoding] & operational-encoding premise (invariance of domain applied to the encoding) & \codeid{check\_L\_cost\_gauge} (\codeid{core.py}) \\
$C(G) = \dim(G)$, composed operational statement (\S\ref{sec:nc_chain}) & [P $\mid$ op.\ encoding $+$ disjoint-anchor realizability] & the operational cost premise pair; cost functor = named open lane & \codeid{check\_L\_cost\_gauge} (\codeid{core.py}) \\
$T_{7B}$ positive-definite cost kernel (\S\ref{sec:T7B}) & [P] & K2, K3, SP, $L_\omega$, $T_{\mathrm{adj}}$ & \codeid{check\_T7B} (\codeid{gravity.py}) \\
$L_{\mathrm{irr,uniform}}$ (\S\ref{sec:Lirr_uniform}) & [P] (sufficient-condition) & saturation bound~(\ref{eq:saturation_bound}) & \codeid{check\_L\_irr\_uniform} (\codeid{core.py}) \\
$B1'$ complex carrier (\S\ref{sec:B1prime}) & [P] & oriented-composite robustness & \codeid{check\_B1\_prime} (\codeid{gauge.py}) \\
Theorem~R prerequisites closed (\S\ref{sec:TR_closed}) & [P] & as listed per carrier requirement & \codeid{check\_Theorem\_R} (\codeid{gauge.py}) \\
F6 $\perp$ EC independence (all four quadrants inhabited; F6-free argmin $=$ 42) & [P] & pure computation over the declared scan space & \codeid{check\_L\_F6\_not\_from\_EC} (\codeid{ec\_inventory\_reading.py}) \\
One-U(1) necessity leg ($\ge 1$ abelian factor) & [P\textsubscript{structural\_reading} $\mid$ R-EC-inv $+$ I-typing] & consumed by the closure theorems of \TSI & \codeid{check\_L\_EC\_inventory\_reading} (\codeid{ec\_inventory\_reading.py}) \\
\bottomrule
\end{tabular}}
\end{center}
\end{tcolorbox}


\medskip
\begin{tcolorbox}[apfbox, title={What Is Assumed, What Is Proved, What Is Deferred}]

\textbf{Imported from Paper~1 Supplement (proved there; re-proved in the appendix for self-containment).}  The first four items are the structurally necessary components of PLEC (none derivable from the others; essentiality proofs in Paper~1 Supplement, \S1).
\begin{enumerate}[nosep,label=\arabic*.,leftmargin=2em]
  \item \textbf{A1} (PLEC upper bound): admissibility capacity is finite and positive.
  \item \textbf{MD} (PLEC positive floor): uniform positive cost per unit of admissibility complexity.
  \item \textbf{A2} (PLEC argmin selection): the realized configuration is the $\arg\min$ of total admissibility cost over the admissible set.
  \item \textbf{BW} (PLEC non-degeneracy): the cost spectrum witnesses a budget-window triple.
  \item \textbf{SP, K3, FD1--FD6}: constitutive admissibility semantics.
  \item $T_{\mathrm{form}}$, $T_{\mathrm{embed}}$, $T_{\mathrm{sep}}$, $L_{\eps^*}$, $L_\Delta$, $L_{\mathrm{blk}}$, $L_\Pi$, $T_{\mathrm{alg}}$, $T_{\mathrm{GNS}}$, T2a--T2c, $T_{\mathrm{Born}}$, $T_M$, $L_{\mathrm{loc}}$, $L_{\mathrm{irr}}$, T3, $L_\omega$, $T_{\mathrm{adj}}$: the finite-carrier chain (each re-proved in Appendix~\ref{app:imported}).
\end{enumerate}

\smallskip\noindent\textbf{Chain results proved in this supplement.}
\begin{enumerate}[nosep,label=\arabic*.,leftmargin=2em]
  \item $L_{\mathrm{Cauchy}}$: cost function uniqueness (\S\ref{sec:cauchy_supp}).
  \item Local support and commuting algebras from disjoint substrates (Definition~\ref{def:local_support}, Lemma~\ref{lem:loc_commut}); the generated-algebra block structure (Proposition~\ref{prop:generated_block_structure}).
  \item Superadditivity to a finite semisimple carrier by the standard $C^*$ route (Proposition~\ref{prop:superadd_to_semisimple_v35}).
  \item Gauge identification: $\mathrm{Gau}_0 \subseteq \mathrm{Inn}(\mathcal{A}) = \prod_k \mathrm{PU}(n_k)$ unconditionally; equality under the named full-invariance hypothesis (\S\ref{sec:gauge_ident}).
  \item Exact sequence $1 \to Z \to \mathrm{U}(\mathcal{A}) \to \mathrm{Inn}(\mathcal{A}) \to 1$ and the relative-phase torus (Proposition~\ref{prop:exact_seq}).
  \item The matter category and the action defined on the lifted group $\widetilde{G}$, determined by the algebra only up to central data, closed by the charge representation as named input (Definition~\ref{def:matter_category}, Proposition~\ref{prop:gauge_acts_on_fields}).
  \item Lift from $\mathrm{PU}(n)$ to $\SU(n)$ and global structure $[\SU(3) \times \SU(2) \times \U(1)_Y]/\mathbb{Z}_6$ (Propositions~\ref{prop:PU_to_SU}, \ref{prop:global_structure}).
  \item EC and CM at their grades (Lemma~\ref{lem:EC_CM_from_A1}); the generator cost bound and the composed operational statement $C(G) = \dim(G)$ (\S\ref{sec:nc_chain}).
  \item $T_{7B}$ (\S\ref{sec:T7B}); $L_{\mathrm{irr,uniform}}$ (\S\ref{sec:Lirr_uniform}); $B1'$ (\S\ref{sec:B1prime}); Theorem~R consolidation (\S\ref{sec:TR_closed}).
\end{enumerate}

\smallskip\noindent\textbf{Consumed downstream by \TSI.}  The
classification core consumes the outputs above as cited premises; the
scan, the waterfall, the one-U(1) closure theorems, hypercharge, and
the capacity count are proved there, not here.

\smallskip\noindent\textbf{Not proved here (honest flags).}
\begin{enumerate}[nosep,label=\arabic*.,leftmargin=2em]
  \item $L_{\mathrm{irr,uniform}}$ is a \emph{sufficient-condition} theorem: it guarantees gauge-sector irreversibility at interfaces satisfying the saturation bound~(\ref{eq:saturation_bound}), not at all interfaces.
  \item The equality leg of the gauge identification consumes the full-invariance hypothesis; only the inclusion is unconditional.
  \item $C(G) = \dim(G)$ rests on the operational cost premise pair pending a cost functor (open lane 2, \S\ref{sec:open_lanes}).
  \item EC rests on R-EC-inv $+$ I-typing; the counter-construction showing non-derivability at shape level is owned in the bank.
  \item The IJC occupancy is empirical by design (named conditioning, \S\ref{sec:open_lanes}).
\end{enumerate}
\end{tcolorbox}

\begin{tcolorbox}[apfbox, title={Imported-ingredient hypothesis ledger (verified vs.\ assumed in the finite-cost environment)}]
\footnotesize
Every off-framework mathematical import, with its hypothesis status
\emph{in this document's finite setting}.  ``Verified'' = the theorem's
hypotheses are checked here against the finite structures it is applied
to; ``assumed'' = the hypothesis is taken as a premise and flagged where
used.

\smallskip
\begin{tabular}{@{}p{0.30\linewidth}p{0.30\linewidth}p{0.32\linewidth}@{}}
\toprule
\textbf{Import} & \textbf{Applied to} & \textbf{Hypothesis status} \\
\midrule
Wedderburn--Artin & finite $*$-algebra $\to \bigoplus_k M_{n_k}$ & verified (finite-dimensional $*$-algebra with faithful positive state; semisimple by the standard $C^*$ route; T2a) \\
Field selection $\mathbb{F}=\mathbb{C}$ & scalar field of the carrier & derived at T2c from $B1'$ orientation; real/quaternionic branches close at the field gate, not assumed away \\
Skolem--Noether & per-block automorphisms & verified (applied blockwise to $M_{n_k}(\mathbb{C})$ only; v4.0 removed the central-simple misstatement) \\
Kadison/Wigner (finite) & symmetry implementation & verified in finite dimension; antiunitary branch excluded by $B1'$, not by assumption \\
Doplicher--Roberts & framing/comparison only & \emph{not consumed as a theorem}: finite-regime departure from type-III noted in \S\ref{sec:finite_vs_typeIII} \\
Cartan closed-subgroup & classification class & verified (the class is $\prod_k \mathrm{PU}(n_k)$ and closed subgroups; scope rider travels with \codeid{check\_L\_cost\_gauge}) \\
Brouwer invariance of domain & resolution-rank bound & verified (finite-dimensional, local form; Lemma~\ref{lem:gen_channel}) \\
\bottomrule
\end{tabular}
\end{tcolorbox}
\begin{tcolorbox}[apfbox, title={Terminology: four words with one fixed meaning each}]
\footnotesize
\begin{description}[nosep,leftmargin=1.6em]
  \item[faithful] of a representation or state: trivial kernel.  For
  states: $\omega(a^*a) = 0 \Rightarrow a = 0$ (GNS null ideal zero).
  Never used to mean ``physically realized.''
  \item[irreducible] of a representation: no proper nonzero invariant
  subspace.  Used only rep-theoretically; ``primitive'' is used for
  invariant tensors, never ``irreducible.''
  \item[effective] of a group action: only the identity acts trivially
  (synonym of faithful for actions); used when the group, rather than
  its representation, is the subject.
  \item[physically distinct (multiplet types)] distinct isomorphism
  classes of $\widetilde{G}$-representations
  (Definition~\ref{def:admissible_gauge}); copies within a class
  (generations) are the same type at different multiplicity.
\end{description}
\end{tcolorbox}

\medskip
\begin{tcolorbox}[apfbox, title={Reader's map: one paper, two supplements}]
\footnotesize
Paper~2's technical material lives in two supplements.  This document
--- Technical Supplement~II, \emph{The Foundational Gauge Program} ---
carries the chain from the APF axioms toward the classification
inputs: the admissibility algebra and its gauge identification, the
matter category and the lift, the EC/CM grounds at their graded
readings, the $C(G)=\dim(G)$ cost program, the type-III/continuum
boundary, and the positioning against reconstruction programs, with a
grade ledger and an open-lane register.  \TSI\ carries the finished
classification mathematics on declared inputs --- the conditional
classification theorem with its C1--C10 ledger, the scan, anomaly and
hypercharge closure, the one-U(1) closure theorems, and the capacity
count.  A reader auditing the classification needs Technical
Supplement~I; a reader auditing where its premises come from, and at
what grade, needs this document.
\end{tcolorbox}

\begin{tcolorbox}[apfbox, title={Axioms-to-SM checklist: minimal load-bearing inputs of the whole program}]
\footnotesize
The map below spans the whole program: the chain rows are proved in this supplement; the template-closure, matter-scan, and hypercharge rows are proved in Technical Supplement~I (\emph{The Classification Core}) and appear here so the reader can see the program end to end.
\footnotesize
\begin{center}
\renewcommand{\arraystretch}{1.12}
\begin{tabular}{@{}p{2.3cm}p{5.2cm}p{5.0cm}@{}}
\toprule
\textbf{Step} & \textbf{Irreducible input} & \textbf{Output used downstream} \\
\midrule
PLEC regime & A1 finite capacity; MD positive floor; A2 minimum-cost realization; BW non-degeneracy; constitutive semantics SP/K3/FD1--FD6 & Finite admissible search space with a positive admissibility ledger. \\
Cost uniqueness & Local additivity on independent interfaces; unit normalization & $F(n)=n$ on the discrete channel domain; all cost comparisons use channel counts rather than tunable weights. \\
Sep/IJC routing & Sep/IJC representation theorem plus the non-separable substrate condition for physical interfaces & Nonzero pool coupling, $\Delta>0$, and hence a noncommutative admissibility algebra. \\
Algebraic carrier & Finite-dimensional $*$-algebra; field selection $\mathbb{F}=\mathbb{C}$; faithful GNS state & $\mathcal{A}\cong\bigoplus_k M_{n_k}(\mathbb{C})$ and linear matter modules. \\
Gauge identification & Internal/kinematic split; Skolem--Noether; cost-equipped block structure & $G_{\mathrm{int}}=\mathrm{Inn}(\mathcal{A})=\prod_k\mathrm{PU}(n_k)$; outer block permutations are discrete relabelings, not continuous gauge. \\
Physical action on fields & Matter-category representation (independently specified) and exact sequence $1\to Z\to U(\mathcal A)\to\mathrm{Inn}(\mathcal A)\to 1$ & Lift from projective algebra action to linear matter action; $\mathrm{PU}(n)$ is represented by its $\SU(n)$ cover where fundamental matter exists. \\
Abelian/global structure & Relative-phase torus; no-redundancy/minimum-cost condition; anomaly-compatible matter modules & Exactly one $\U(1)_Y$ and global quotient $[\SU(3)\times\SU(2)\times\U(1)_Y]/\mathbb Z_6$. \\
Template closure & Lie classification; oriented composite robustness; Witten anomaly; capacity minimality & $\SU(3)\times\SU(2)\times\U(1)_Y$ as the unique minimum-cost gauge template. \\
Matter scan & Bounded 1{,}680-template enumeration; F1--F7 exact filters; three-generation input from the capacity ledger & Unique 45-Weyl-fermion SM survivor up to CPT. \\
Hypercharge & Standard anomaly equations applied to the derived template & Unique rational hypercharge ray and charge quantization. \\
\bottomrule
\end{tabular}
\end{center}
\smallskip
\noindent
Necessity is local to each row: relaxing a row does not refute APF, but it
opens the indicated alternatives.  For example, dropping the non-separable
substrate condition leaves the Sep countermodel with $\Delta=0$; dropping
the relative-phase no-redundancy condition allows extra sterile $\U(1)$
directions; dropping the finite-regime assumption restores type-III AQFT
as an effective continuum limit rather than a finite proof environment.
\end{tcolorbox}
\begin{tcolorbox}[apfbox, title={Framing convention: ``forces'' and ``selects'' denote uniqueness, not processes}]
\footnotesize
The supplement's theorems and proofs use the verbs ``forces,'' ``selects,'' and ``rejects'' as standard mathematical shorthand for uniqueness results, $\arg\min$ properties, and inadmissibility results respectively.  Under the APF ontology, these verbs do not denote processes: nothing is forced by an agent, nothing is selected by a mechanism, nothing is rejected by a filter.  Each theorem asserts a property of the set of admissible configurations---that it has a unique member, a minimum member, or no member satisfying some condition.  For example, ``capacity minimality selects $N_c = 3$'' is shorthand for ``$N_c = 3$ is the $\arg\min$ of dim(G) over R1--R3-admissible color groups,'' and the alternative values (SU(4), SU(5), E$_6$, \ldots) are \emph{admissible} in the sense that each has cost below the capacity bound, but have higher cost than SU(3) and are therefore not what is.  The companion note \emph{What Is, Is Minimum-Cost Admissible} spells this out in full.  Readers may substitute the descriptive reading throughout the supplement without any change to the proofs.

\begin{tcolorbox}[apfbox, title={Status of A2: argmin realization is a lawlike selection rule, not a derived dynamics}]
\footnotesize
A2 is intentionally not derived in this supplement from a deeper dynamical variational principle.  It is one of the four structural inputs of the finite-regime APF package: among admissible continuations satisfying A1, MD, and BW, the realized carrier is the minimum admissibility-cost representative.  All selection claims in the supplement have the conditional form
\[
  \text{A1+MD+BW+A2+listed interface hypotheses} \;\Longrightarrow\; \text{unique minimum-cost admissible structure}.
\]
Thus the paper does not claim to prove that nature \emph{must} obey global argmin realization from conventional QFT assumptions.  It claims that, if A2 is accepted as the APF lawlike selection rule, the Standard-Model template is the unique admissible minimum under the stated finite ledger.  This is why counterfactuals that change A2, make new generators cost-free, or subtract structural gauge generators from the ledger are not small variants of the theorem; they define a different selection problem.
\end{tcolorbox}
\end{tcolorbox}
\begin{tcolorbox}[apfbox, title={Operational/structural convention}]
\footnotesize
Throughout the supplement, \emph{operations} are tests, preparations,
updates, comparisons, and continuations performed at an interface.  A
\emph{structure} is the equivalence class, ledger, projection, module, or
algebra induced by those operations after quotienting observationally
indistinguishable continuations.  Thus when a structural object is said to
``act,'' the formal meaning is always its represented action on operational
records or on the field module generated by those records.  This convention
prevents an interpretive slippage that would otherwise make properties
sound like mechanisms: the algebra is not a physical agent; it is the
minimal representation of admissible operational continuations.
\end{tcolorbox}
\begin{tcolorbox}[apfbox, title={Translation table for nonstandard vocabulary}]
\footnotesize
\begin{center}
\renewcommand{\arraystretch}{1.08}
\begin{tabular}{@{}p{3.2cm}p{9.0cm}@{}}
\toprule
\textbf{APF term} & \textbf{Standard-language approximation} \\
\midrule
Pool $\Pi$ & Residual/shared substrate direction not resolved by the current coarse record algebra; comparable to a gluing degree of freedom between contexts. \\
Anchor $M_d$ & Context-fixed support subspace for a tested distinction; analogous to a local sector or record support. \\
IJC & Failure of a faithful global Boolean refinement preserving all tested marginals and the defined partial products; finite analogue of a contextual global-section obstruction. \\
Continuation calculus & Typed partial algebra of operational histories, quotiented by observational equivalence; comparable to a finite process/category-of-contexts presentation. \\
Admissibility algebra & Finite $*$-algebra generated by record projections and continuation-preserving updates. \\
Capacity cost & Resource monotone for holding distinctions available; in QFT language, a finite-regime bookkeeping surrogate for the cost of supporting admissible distinctions through independent channels. \\
\bottomrule
\end{tabular}
\end{center}
\end{tcolorbox}

\medskip
\medskip\noindent\textbf{Glossary of imported objects.}  Each entry gives the object's name, one-sentence definition, source in the Paper~1 Supplement (P1S), and where it is first used in this supplement.

\smallskip
\begin{center}\small
\renewcommand{\arraystretch}{1.15}
\begin{tabular}{@{}lp{5.8cm}ll@{}}
\toprule
\textbf{Symbol} & \textbf{Definition} & \textbf{Source} & \textbf{Used at} \\
\midrule
$\Gamma$ & Enforcement interface (physical boundary) & P1S FD1 & Throughout \\
$S_\Gamma$ & Substrate: set of configurations at $\Gamma$ & P1S FD1 & Throughout \\
$\mathcal{D}(\Gamma)$ & Admissible distinctions (binary predicates, $\eps > 0$) & P1S FD1 & \S\ref{sec:cauchy_supp} \\
$\eps(d)$ & Admissibility cost: $\inf_{\mathcal{P}(d)} \kappa$ & P1S FD4 & \S\ref{sec:cauchy_supp} \\
$C(\Gamma)$ & Total admissibility budget at $\Gamma$ & P1S A1 & Throughout \\
$\eps^*$ & Uniform cost floor: $\eps(d) \ge \eps^*$ for all $d$ & P1S $L_{\eps^*}$ & \S\ref{sec:cauchy_supp} \\
$M_d$ & Anchor set: minimal perturbation-support subspace & P1S FD3 & \S\ref{sec:nc_chain} \\
$\Pi$ & Pool: $S_\Gamma \ominus \bigoplus_d M_d$ & P1S $T_{\mathrm{sep}}$ & \S\ref{sec:nc_chain} \\
$E_d$ & Admissibility projection: $B$-orth.\ proj.\ onto $M_d$ & P1S FD5/$T_{\mathrm{adj}}$ & \S\ref{sec:nc_chain} \\
$B(\cdot,\cdot)$ & Sector-orthogonal bilinear form on $V$ & P1S $L_\omega$ & \S\ref{sec:T7B} \\
$\mathcal{A}$ & Admissibility algebra: $\mathrm{Alg}_\mathbb{R}(\{E_d, E_{\{d_i,d_j\}}\})$ & P1S O3 & \S\ref{sec:gauge_ident} \\
$\omega$ & Faithful positive state on $\mathcal{A}$ & P1S O4 & \S\ref{sec:gauge_ident} \\
$\mathcal{H}_\omega$ & GNS Hilbert space from $(\mathcal{A}, \omega)$ & P1S $T_{\mathrm{GNS}}$ & \S\ref{sec:gauge_ident} \\
$\Delta(d_1,d_2)$ & Superadditive surplus: $\eps(\{d_1,d_2\}) - \eps(d_1) - \eps(d_2)$ & P1S $L_\Delta$ & \S\ref{sec:nc_chain} \\
\bottomrule
\end{tabular}
\end{center}

\smallskip\noindent\textbf{APF $\leftrightarrow$ standard terminology.}

\smallskip
\begin{center}\small
\renewcommand{\arraystretch}{1.15}
\begin{tabular}{@{}ll@{}}
\toprule
\textbf{APF term} & \textbf{Standard / GPT equivalent} \\
\midrule
Distinction $d$ & Effect / binary test / observable outcome \\
Admissibility projection $E_d$ & Sharp effect / compression / filter \\
Pool $\Pi$ & Shared latent support / coupling sector \\
Superadditivity surplus $\Delta$ & Nonclassicality witness / coupling measure \\
Sector decomposition ($T_{\mathrm{sep}}$) & Face decomposition of state space \\
Admissibility algebra $\mathcal{A}$ & Observable algebra / $C^*$-algebra of observables \\
GNS representation & Gelfand--Naimark--Segal construction \\
\bottomrule
\end{tabular}
\end{center}


% ══════════════════════════════════════════════════════════════


\medskip\noindent\textbf{Dependency table.}  Each row gives a result proved in this supplement, its inputs, and the downstream results that depend on it.

\smallskip
\begin{center}\footnotesize
\renewcommand{\arraystretch}{1.15}
\begin{tabular}{@{}lp{4.7cm}p{4.5cm}@{}}
\toprule
\textbf{Theorem} & \textbf{Inputs} & \textbf{Used by} \\
\midrule
$L_{\mathrm{Cauchy}}$ (\S\ref{sec:cauchy_supp}) & A1, $L_{\mathrm{loc}}$, $L_{\eps^*}$ & cost ranking (\TSI) \\
Gauge ident.\ (\S\ref{sec:gauge_ident}) & $T_{\mathrm{alg}}$, $T_{\mathrm{sep}}$, $L_{\mathrm{loc}}$, T3 & Theorem R; closure theorems (\TSI) \\
Exact sequence / $T_{\mathrm{rel}}$ (\S\ref{sec:gauge_ident}) & gauge ident. & abelian closure (\TSI) \\
Matter category $+$ lift (\S\ref{sec:gauge_ident}) & exact sequence; charge rep (I-typing, named input) & $T_{\mathrm{field}}$'s category (\TSI) \\
EC / CM (\S\ref{sec:nc_chain}) & R-EC-inv $+$ I-typing (EC); A2 over A1$+$MD (CM) & closure theorems (\TSI) \\
$C(G)$ program (\S\ref{sec:nc_chain}, \S\ref{sec:cost_program}) & operational premise pair & cost ranking (\TSI) \\
$T_{7B}$ (\S\ref{sec:T7B}) & A1, K2, K3, SP, $L_\omega$, $T_{\mathrm{adj}}$ & $L_{\mathrm{irr,uniform}}$ \\
$L_{\mathrm{irr,uniform}}$ (\S\ref{sec:Lirr_uniform}) & $L_{\mathrm{loc}}$, $L_{\mathrm{nc}}$, $L_{\mathrm{irr}}$, $T_{7B}$, saturation bound & Theorem R (R2 chain) \\
$B1'$ (\S\ref{sec:B1prime}) & A1 (minimality), $T_M$, rep.\ theory & Theorem R (R1 chain) \\
\bottomrule
\end{tabular}
\end{center}

\tableofcontents

\medskip
\noindent\textit{On operational language.}\quad
Operational verbs are local readings of an underlying admissibility-space
structure under a temporal slicing.  They are retained because they are the
working vocabulary of physics; the structural reading is always available:
``acts,'' ``forms,'' ``selects,'' and ``evolves'' mean that the corresponding
admissible configuration, quotient, minimum, or continuation relation exists.
This convention is inherited from Paper~0 and Paper~1 Supplement v8.24.


\newpage

\part{The Foundational Chain}
\label{part:chain}

\section{Cost function uniqueness ($L_{\mathrm{Cauchy}}$)}
\label{sec:cauchy_supp}
% ══════════════════════════════════════════════════════════════

The admissibility cost function $F$ maps channel count to capacity committed.
Paper~2 \S4 states the uniqueness result; this section provides the full
formal treatment, specifying domain, regularity, and operational meaning
of each condition.

% ======================================================================
% Paper 2 Technical Supplement — §1
% Cost Function Uniqueness via Cauchy's Functional Equation
% ======================================================================
%
% PURPOSE.  Paper 2 §4 states that the admissibility cost function
% F(n) = n is unique under additivity, monotonicity, and unit
% normalization.  The reviewer asks (Q2): "What are the exact hypotheses
% under which Cauchy's functional equation yields F(d) = d in your
% setting?  What is the domain of d, and how is additivity defined
% operationally?"  This section provides the full formal treatment.
%
% UPSTREAM: A1, L_loc, L_cost (from Paper 1 Supplement).
% NO downstream results in Paper 2 require extension beyond N.
% The rational extension is needed only for Paper 4A (Weinberg angle).
% ======================================================================


\subsection{Domain and operational definitions}
\label{sec:cauchy_domain}

The admissibility cost function $F$ maps the \emph{number of admissibility
channels} consumed by a distinction to the capacity committed.  We now
specify its domain, its operational meaning, and the three conditions
that determine it uniquely.

\begin{definition}[Enforcement channel count]\label{def:channel_count}
At an interface $\Gamma$ with finite capacity $C(\Gamma)$ (A1), each
irreducible distinction $d \in \mathcal{D}_{\mathrm{irr}}(\Gamma)$ costs
$\varepsilon^*$ to enforce ($L_{\varepsilon^*}$, Paper~1 Supplement).
A composite distinction occupying $n$ independent irreducible channels
has cost $\varepsilon(\{d_1, \ldots, d_n\}) = n \cdot \varepsilon^*$
when the $n$ channels are at \emph{separate interfaces} ($L_{\mathrm{loc}}$),
or $n \cdot \varepsilon^* + \Delta$ when co-located ($L_\Delta$):
$\Delta \ge 0$ in general, and strictly $\Delta > 0$ exactly when the
co-located channels draw on a non-empty shared pool ($\Pi \neq \{0\}$;
the surplus is $\Delta = \kappa_\Pi > 0$, $L_\Delta$, and vanishes,
$\Delta = 0$, when $\Pi = \{0\}$).  The \emph{channel count} is the integer
$n(d) := \varepsilon(d) / \varepsilon^*$ evaluated at separate interfaces,
where the surplus $\Delta = 0$.
\end{definition}

\smallskip\noindent\textit{The additive/superadditive determinant is the shared pool.}\quad
Whether two co-located channels add ($\Delta = 0$) or super-add ($\Delta > 0$) is
not a free per-case choice but is fixed by one structural variable: the shared
pool $\Pi$ in the sector decomposition $S_\Gamma = \bigoplus_d M_d \oplus \Pi$
($T_{\mathrm{sep}}$).  The $L_\Delta$ surplus is $\Delta = \kappa_\Pi$ with
$\kappa_\Pi > 0 \iff \Pi \neq \{0\}$, so $\Pi = \{0\}$ (separated, no common
reservoir) gives additivity and $\Pi \neq \{0\}$ (genuine co-location) gives the
surplus.  The separated-interface additivity used by $L_{\mathrm{loc}}$ and C1 and
the superadditivity used by $L_\Delta$ and the confinement mechanism are therefore
one surplus functional evaluated on $\Pi = \{0\}$ versus $\Pi \neq \{0\}$
configurations, not two opposed cost primitives.  Accordingly the structural
existence ledger $L_{\mathrm{count}}$ is additive because its matter and Higgs
channels are separated ($T_M + L_{\mathrm{loc}}$, $\Pi = \{0\}$); the
gauge-generator contribution is additive at the level of resolution
channels (one independent distinction per group-manifold dimension; the
generator cost bound, Lemma~\ref{lem:gen_channel}), with the cost-exactness
$C(G) = \dim G$ carried as the composed operational statement --- lower
bound on the operational-encoding premise, upper bound on the
disjoint-anchor realizability premise --- banked as
\coderef{check\_L\_cost\_gauge}{core.py} and realized paper-side in
Proposition~4.1 of \TSI{}.

\noindent The domain of $F$ is therefore the positive integers:
$F: \mathbb{N} \to \mathbb{R}_{>0}$, where $F(n)$ gives the capacity
cost (in units of $\varepsilon^*$) of enforcing $n$ independent channels
at separate interfaces.

\begin{remark}[Why the domain is $\mathbb{N}$, not $\mathbb{R}^+$]
\label{rem:domain_N}
Enforcement channels are discrete.  This discreteness is not assumed ---
it is derived from A1 + MD:
\begin{enumerate}[label=(\roman*),nosep]
  \item A1 provides finite capacity $C(\Gamma) < \infty$.
  \item $L_{\varepsilon^*}$ (from MD) provides a positive cost floor
        $\varepsilon^* > 0$ for each irreducible channel.
  \item Therefore the maximum number of independently admissible channels
        at any interface is $\lfloor C(\Gamma) / \varepsilon^* \rfloor
        < \infty$, an integer.
\end{enumerate}
The domain $\mathbb{N}$ is a \emph{consequence}, not a choice.  The
classical Cauchy problem on $\mathbb{R}^+$ --- with its pathological
non-measurable solutions --- does not arise because the domain is
already discrete.
\end{remark}


\subsection{The three conditions and their derivation}
\label{sec:cauchy_conditions}

\begin{theorem}[$L_{\mathrm{Cauchy}}$: Cost Function Uniqueness]\label{thm:LCauchy}
Let $F: \mathbb{N} \to \mathbb{R}_{>0}$ satisfy:
\begin{align}
  &\textbf{C1 (Additivity):}  &&F(a + b) = F(a) + F(b)
    \quad \text{for all } a, b \in \mathbb{N}, \label{eq:C1_supp} \\
  &\textbf{C2 (Monotonicity):} &&a > b \;\Rightarrow\; F(a) > F(b)
    \quad \text{for all } a, b \in \mathbb{N}, \label{eq:C2_supp} \\
  &\textbf{C3 (Unit):}          &&F(1) = 1. \label{eq:C3_supp}
\end{align}
Then $F(n) = n$ for all $n \in \mathbb{N}$.
\end{theorem}

\begin{proof}
By C1 (induction): $F(n) = F(1 + 1 + \cdots + 1) = n \cdot F(1)$.
By C3: $F(1) = 1$.  Therefore $F(n) = n$ for all $n \in \mathbb{N}$.
C2 is satisfied automatically: $n > m \Rightarrow F(n) = n > m = F(m)$.
\end{proof}

\noindent The proof on $\mathbb{N}$ is trivial once the three conditions
are established.  The substantive content is the \emph{derivation of the
conditions from A1}:

\begin{proposition}[Derivation of C1--C3 from admissibility]\label{prop:C1C3}
\leavevmode
\begin{enumerate}[label=\textbf{C\arabic*},nosep,leftmargin=3em]
  \item \textbf{(Additivity).}  Let $d_1, d_2$ be distinctions at
  \emph{separate} interfaces $\Gamma_1, \Gamma_2$.  By $L_{\mathrm{loc}}$
  (Paper~1 Supplement~\cite{BrookeSupplement}, \S4): admissibility at separate
  interfaces is independent, with independent budgets.  The total
  admissibility cost is $\varepsilon(d_1) + \varepsilon(d_2)$ --- costs
  add exactly when the admissibility regions are disjoint.  In channel-count
  language: $F(n(d_1) + n(d_2)) = F(n(d_1)) + F(n(d_2))$.

  The operational meaning of ``additivity'' is therefore:
  \emph{independently admissible channels at separate interfaces
  contribute independently to the total capacity commitment.}  This is
  the content of $L_{\mathrm{loc}}$ applied to the cost functional.

  \textit{What C1 does not claim:} C1 does not assert additivity at a
  \emph{shared} interface.  At a shared interface, the superadditive
  surplus $\Delta > 0$ ($L_\Delta$) produces non-additive costs.
  C1 applies only to the separated-interface regime, which is precisely
  the regime relevant to the cost function $F$ (Definition~\ref{def:channel_count}).

  \item \textbf{(Monotonicity).}  More admissibility channels require
  more capacity.  This is immediate from A1: each channel costs
  $\varepsilon^* > 0$ ($L_{\varepsilon^*}$), so adding a channel
  increases the total cost.

  \item \textbf{(Unit).}  $F(1) = 1$ is a choice of capacity units:
  we measure cost in units of the single-channel cost $\varepsilon^*$.
  This is a normalization convention, analogous to $c = 1$ in natural
  units.  It introduces no physical content.
\end{enumerate}
\end{proposition}

\noindent\textit{Marginal-floor refinement (Paper~10 alignment).}\quad
The cost-function uniqueness theorem above operates on the separated-interface regime where, by $L_{\mathrm{loc}}$, costs add exactly: $\Delta = 0$.  In the broader continuation calculus of Paper~10 (\emph{The Calculus of Finite Continuability}) v1.11, \S3, this regime corresponds to a substrate with no shared structural commitments, where the in-isolation cost floor and the marginal floor coincide.  In regimes with shared substrate commitments --- broken vacuum, confined-phase QCD substrate, topological infrastructure --- the in-isolation cost overcounts (the shared piece is paid once by the carrying context, not once per distinction), and the marginal floor
\[
        \varepsilon^*_\Gamma \;:=\; \inf\bigl\{\kappa_\Gamma(S \otimes d) - \kappa_\Gamma(S) : S \text{ admissible},\ d \text{ tested},\ S\otimes d \text{ admissible}\bigr\}
\]
becomes the operative invariant.  Paper~2's $L_{\mathrm{Cauchy}}$ and the downstream $L_{\mathrm{count}}$ count $C_{\mathrm{total}} = 61$ pass through unchanged in either reading: $L_{\mathrm{Cauchy}}$ governs the separated-interface regime where the floors agree; the SM-interface count is read as a count of marginal channel-distinctions on top of the broken-vacuum and confined-phase substrate.  Paper~4 Supplement v1.5 \S2.2 applies the marginal-floor reading to $T_{4F}$ (the three-generations bound) explicitly.  No theorem in this section requires modification.

\noindent\textit{Connection to resource theories and GPTs.}\quad
The conditions C1--C3 have close analogues in resource-theoretic
frameworks: C1 corresponds to the tensor-product structure of
independent resources; C2 to monotonicity under free operations; C3
to the choice of a gold standard (unit resource).  The operational
stance --- deriving algebraic structure from resource constraints ---
is shared with GPT-style reconstructions (Hardy~\cite{Hardy2001},
CDP~\cite{CDP2011}) and the ``accessible fragments'' approach.  The
difference is scope: those programs reconstruct the state space; this
program uses the same starting point to derive gauge content and
matter spectrum.


\subsection{Rational extension (for Paper~4A)}
\label{sec:cauchy_rational}

Paper~2's results ($L_{\mathrm{gauge,unique}}$, $T_{\mathrm{field}}$,
$L_{\mathrm{count}}$, $P_{\mathrm{exhaust}}$, $L_{\mathrm{equip}}$) use
$F$ only on $\mathbb{N}$.  The Weinberg angle derivation (Paper~4A)
evaluates $F$ at the rational argument $1/d$ (the reciprocal channel
count), producing $\gamma = F(d) + F(1/d) = d + 1/d$.  We record the
rational extension here for completeness.

\begin{proposition}[Rational extension of $F$, with operational premise]\label{prop:rational_ext}
\textbf{Premise} (\textbf{Sharing-admissibility, S}): A single
admissibility channel evaluated at fractional occupancy $1/q$ ($q \in \mathbb{N}$)
contributes $F(1/q)$ to each of $q$ separated-interface sectors that
jointly share the channel, and the $q$ shared contributions combine
additively under C1.  Equivalently, $F$ is extended to $\mathbb{Q}_{>0}$
as the unique additive extension consistent with C1 applied $q$ times
across sharing.

\smallskip\noindent\textbf{Conclusion.}  Under premise S, $F(p/q) = p/q$
for all $p, q \in \mathbb{N}$.
\end{proposition}

\begin{proof}
For any $q \in \mathbb{N}$: by S, $F$ is defined at $p/q$; by C1
applied $q$ times, $q \cdot F(p/q) = F(p/q + \cdots + p/q) = F(p) = p$.
Therefore $F(p/q) = p/q$.
\end{proof}

\smallskip\noindent\textit{Logical status of the sharing premise.}\quad
Premise S is not derivable from C1 alone: C1 as stated in the integer
domain ($a, b \in \mathbb{N}$) is vacuous at fractional arguments, so
``C1 extended to $\mathbb{Q}_{>0}$'' is definitionally a new domain
assignment.  S supplies the operational content: it interprets $F(1/q)$
as the per-sector share of a channel decomposed across $q$ separated
interfaces.  The interpretation is consistent with $L_{\mathrm{loc}}$
(locality of capacity) and $T_M$ (interface monogamy)---sharing at
separated interfaces does not incur superadditive surplus---but it is
a reading beyond C1's native domain.  The honest flag is that the
Weinberg-angle derivation in Paper~4A rests on C1~$+$~S, not on C1
alone.  S is therefore listed alongside MD, BW, and the constitutive
admissibility semantics (SP, K3, FD1--FD6) in the assumption inventory
when Paper~4A uses the rational extension.

\noindent\textit{Operational justification for the rational extension.}\quad
The rational argument $1/d$ arises when the cost functional is evaluated
on a \emph{fractional} channel occupancy: a single admissibility channel
shared among $d$ sectors contributes $1/d$ of a full channel to each
sector.  This is the capacity-accounting interpretation of the Weinberg
angle formula's ``inverse channel count'' term.  The extension from
$\mathbb{N}$ to $\mathbb{Q}_{>0}$ requires no additional regularity
assumption (no measurability, no continuity): it follows from C1 applied
$q$ times, as the proof above shows.

\smallskip\noindent\textit{Exact status of the rational extension in the axiomatics.}\quad
There are two different statements, and the supplement keeps them separate.
First, on the native domain of $L_{\mathrm{Cauchy}}$, $F:\mathbb{N}\to\mathbb{N}$, C1 determines $F(n)=n$; no fractional arguments occur.  Second, once one \emph{extends the domain} to shared fractional channel occupancies, premise S supplies the operational meaning of $F(p/q)$ and C1 then forces the unique value $F(p/q)=p/q$.  Therefore the rational extension is not a theorem of C1 on $\mathbb{N}$ alone.  It is a theorem of \emph{C1 on the enlarged sharing-admissible domain} C1$+$S.  This is the version used in Paper~4A whenever $F(1/d)$ appears.  No measurability, continuity, or real-domain Cauchy assumption is added; the only added premise is the operational domain-extension premise S.

\smallskip\noindent\textit{No real extension needed.}\quad The classical
Cauchy problem on $\mathbb{R}^+$ admits pathological (non-measurable,
non-monotone, additive) solutions constructible via the Axiom of Choice.
These are excluded by C2 (Darboux, 1875): any Lebesgue-measurable (or
merely bounded on an interval) additive function $f: \mathbb{R}^+ \to
\mathbb{R}$ satisfying $f(x+y) = f(x) + f(y)$ and $f(1) = 1$ is
$f(x) = x$.  However, no result in this paper or in Paper~4A requires
the domain to extend beyond $\mathbb{Q}_{>0}$.  The Darboux theorem is
available as a safety net but is never load-bearing.
\coderef{check\_L\_Cauchy\_uniqueness}{L\_Cauchy\_uniqueness.py}


\subsection{What L$_{\mathrm{Cauchy}}$ eliminates}
\label{sec:cauchy_consequences}

Without $L_{\mathrm{Cauchy}}$, the cost function $F$ is a free parameter.
Any monotone function $F: \mathbb{N} \to \mathbb{R}_{>0}$ with $F(1) = 1$
would be compatible with A1.  With $L_{\mathrm{Cauchy}}$, the class
collapses to a single element:

\smallskip
\begin{center}
\small
\begin{tabular}{@{}lll@{}}
\toprule
\textbf{Candidate} & \textbf{Condition violated} & \textbf{Consequence if allowed} \\
\midrule
$F(n) = n^2$ & C1 ($9 = F(3) \ne F(1) + F(2) = 5$) & Weinberg angle shifts \\
$F(n) = \sqrt{n}$ & C1 ($\sqrt{3} \ne 1 + \sqrt{2}$) & Capacity non-additive \\
$F(n) = \log n$ & C1 ($\log 3 \ne 0 + \log 2$); C3 ($F(1) = 0$) & No unit cost \\
$F(n) = c$ (constant) & C2 ($F(2) = F(1)$); C1 ($c \ne 2c$) & Cost insensitive \\
$F(n) = n^{3/2}$ & C1 ($3\sqrt{3} \ne 1 + 2\sqrt{2}$) & Weinberg angle shifts \\
\bottomrule
\end{tabular}
\end{center}

\smallskip\noindent Each alternative is excluded by C1 alone (additivity).
The uniqueness is not a consequence of fine-tuning or accidental
cancellation: it is a theorem of algebra (the only additive function on
$\mathbb{N}$ with $F(1) = 1$ is the identity).  C2 is logically
redundant on $\mathbb{N}$ once C1 and C3 are established, but it is
stated explicitly because it is the physically motivated condition
(more channels cost more) and because it becomes non-redundant on
$\mathbb{R}^+$ where it excludes pathological solutions.


\newpage


% ══════════════════════════════════════════════════════════════
\section{From superadditivity to gauge structure}
\label{sec:nc_chain}
% ══════════════════════════════════════════════════════════════

This section restates the derivation chain $L_\Delta \to L_{\mathrm{blk}}
\to L_\Pi \to T_{\mathrm{alg}}$ (proved in the Paper~1 Supplement) in
condensed form, and then proves the gauge identification theorem that
separates internal (gauge) automorphisms from kinematic (diffeomorphism)
automorphisms.

The chain is summarized first; the proofs are in the Paper~1 Supplement
and are not reproduced.  The gauge identification theorem
(\S\ref{sec:gauge_ident}) is new to this supplement.

\smallskip\noindent\textbf{IJC routing.}\quad
The chain $L_\Delta \to L_{\mathrm{blk}} \to L_\Pi \to T_{\mathrm{alg}}$
recapitulated below assumes the (IJC) branch of the Sep/IJC
Representation Theorem (Paper~1 Supplement v8.9 \S``Sep/IJC
representation theorem'').
Branch (IJC) is the substrate-perturbation joint non-separability
condition: at a co-located pair $(d_1, d_2)$, the joint threat class
$\mathsf{T}(d_1, d_2)$ strictly exceeds the union
$\mathsf{T}(d_1) \cup \mathsf{T}(d_2)$.  In the alternative branch
(Sep), every joint perturbation reduces to a separable convex
combination of single-side perturbations, the pool $\Pi$ is inert,
and Link~1's $\Delta > 0$ degenerates to $\Delta = 0$ (the spectator
countermodel $V = M_{d_1} \oplus M_{d_2} \oplus \Pi$ with $\Pi$
inert; see Paper~1 Supplement v8.41 \S``Worked models: Sep, IJC,
and coupled-but-classical'').
PLEC alone does not force branch (IJC); it must be flagged as a
substrate-level premise.  The Sep/IJC dichotomy itself is structural
(certified by \texttt{check\_T\_IJC\_dichotomy} and the inseparability
bridge \texttt{check\_T\_inseparable\_IJC}, \texttt{apf/core.py}): branch
(Sep) is the substrate-factorizable case (a commuting hidden-variable
extension exists), branch (IJC) the inseparable one for which no such
factorization exists.  The branch verdict has two parts the framework keeps distinct.  The
\emph{bridge} --- that an inseparable (IJC) interface forces a
noncommutative record algebra --- is DERIVED internally (Paper~5
Supplement v6.8 Theorem \emph{general-finite-query-noncommutative-bridge},
graded $[\mathrm{P_{math}}{+}\mathrm{P_{APF}}]$): given an interface's
declared finite records, branch~(IJC) holds iff no faithful Boolean
global-section defender exists --- an internal Boole / LP feasibility
question --- and Bell, CHSH, Fine, and Kochen--Specker are RECOVERED as
special marginal presentations of this criterion, not imported as
theorems.  The \emph{occupancy}, by contrast, remains empirical: that a
physical quantum-capable interface presents records (its measured
correlation table) lying outside the Boole polytope is read off the
empirical marginals, not derived from A1 --- exactly as the framework
reads the Planck parameters and lattice-QCD constants off their records,
with Bell/KS serving as the in-house diagnostics for whether those records
clear the polytope.  (The earlier ``substrate-richness forces (IJC)''
reading is superseded: substrate-richness permits but does not force
inseparability.)  Theorems in this section therefore carry the
$[\mathrm{P}{+}\mathrm{IJC}]$ tag in the sense ``proved given the (IJC)
branch occupancy at this interface'' --- the bridge internal and derived,
the occupancy empirically certified.


\begin{proposition}[Worked example: the rotated three-dimensional Sep/IJC witness]
\label{prop:minimal_ijc_witness}
Let $V=\mathrm{span}\{e_1,e_2,p\}$ with individual anchor spaces
$M_1=\mathbb{R}e_1$, $M_2=\mathbb{R}e_2$, and pool
$\Pi=\mathbb{R}p$.  In the Sep branch the joint defender is
$E_{12}^{\mathrm{Sep}}=E_1+E_2$ and $[E_1,E_{12}^{\mathrm{Sep}}]=0$.
In the IJC branch the witness model's joint defender is taken in the
rotated two-parameter form
\[
  E_{12}^{\mathrm{IJC}}
  = \pi_{\mathrm{span}\{\cos\theta\,e_1+\sin\theta\,p,\;e_2\}},
  \qquad 0<\theta<\pi/2,
\]
--- the defining ansatz of this witness model, legitimate because the
proposition's role is to exhibit a witness, not to derive a normal form ---
so that
\[
  [E_1,E_{12}^{\mathrm{IJC}}]
  = \cos\theta\sin\theta\bigl(|e_1\rangle\langle p|
    - |p\rangle\langle e_1|\bigr) \ne 0.
\]
Thus the Sep countermodel is exactly the $\theta=0$ case; it is excluded
precisely when the substrate geometry forces nonzero pool rotation.
\end{proposition}

\begin{proof}
In the Sep branch every joint perturbation decomposes into separate
single-side perturbations.  The least-cost joint support is then the
direct sum $M_1\oplus M_2$, hence $E_{12}^{\mathrm{Sep}}=E_1+E_2$ and
all generated projections commute.  This is a genuine countermodel to
any attempted derivation of noncommutativity from A1 alone.

In the IJC branch there exists a joint threat not representable as a
convex mixture of single-side threats.  A direct-sum defender leaves a
pool-supported component undefended, so within the ansatz family the
least-cost joint defender rotates into $\Pi$ by an angle $\theta\ne0$.  Writing
$w_\theta=\cos\theta\,e_1+\sin\theta\,p$, the joint projection is
$|w_\theta\rangle\langle w_\theta|+|e_2\rangle\langle e_2|$.  A direct
multiplication gives the displayed commutator.  The obstruction is
therefore not definitional: Sep and IJC have the same individual anchor
spaces and the same finite ledger, but only IJC supplies the gluing
obstruction that prevents a faithful all-commuting representation.
\end{proof}

\noindent\textit{Interpretive consequence.}\quad
This is a worked example exhibiting the mechanism --- it is not a
universal derivation of the rotated normal form, and no downstream
result consumes it as one.  The toy model is the finite-dimensional
version of the reviewer-facing ``global section'' obstruction.  In Sep, local defenders glue to a single
Boolean refinement.  In IJC, the local defenders agree on their individual
marginals but the joint defender is rotated into the pool, so no faithful
Boolean refinement preserves both the tested marginals and the partial
product.  The bridge theorem therefore depends on an explicit, checkable
substrate condition: nonzero pool rotation, equivalently $\Delta>0$.

\begin{tcolorbox}[apfbox, title={$\Delta=0$ falsifiability branch}]
\footnotesize
The witness above is not itself a proof that every physical interface lies
in the IJC branch.  It proves that, once an interface has nonzero pool
rotation, the Sep Boolean-refinement defender is unavailable.  The physical
necessity of $\Delta>0$ is inherited from the Paper~1 Sep/IJC representation
theorem together with the empirical Bell--Kochen--Specker record: interfaces
that exhibit Bell/Tsirelson or contextuality signatures cannot be faithfully
represented by a single Boolean refinement.  Conversely, if an interface is
well approximated by $\Delta=0$, APF predicts the Sep branch: classically
composable records, no context-dependent holonomy, no Bell/Tsirelson
violation sourced by that interface, and no independent gauge curvature
from that interface.  Thus the branch condition is falsifiable rather than
smuggled.
\end{tcolorbox}

\begin{proposition}[Stand-alone proof spine: superadditivity to finite semisimple carrier]
\label{prop:superadd_to_semisimple_v35}
Assume a finite interface $\Gamma$ satisfying A1, MD, the IJC branch condition $\Delta>0$, the positive cost kernel $B$, closure of tested continuations under adjoint, and finite resolution.  Then the admissibility records generated by the interface form a finite-dimensional semisimple $*$-algebra whose noncommuting part is nontrivial.
\end{proposition}

\begin{proof}
A1 and finite resolution bound the number of mutually admissible primitive record distinctions at $\Gamma$; hence the algebra generated by tested finite continuations is finite-dimensional after quotienting by observational equivalence.  Positivity of $B$ and adjoint closure make that quotient a finite-dimensional complex $*$-algebra carrying a faithful positive state, hence (by GNS) a faithful $*$-representation on a finite-dimensional Hilbert space.  A finite-dimensional adjoint-closed complex subalgebra of $\mathcal{B}(\mathcal{H})$ is a finite-dimensional $C^*$-algebra, and every finite-dimensional $C^*$-algebra is semisimple (its Jacobson radical is a nil ideal, and a $C^*$-algebra has no nonzero nil ideals); Wedderburn--Artin then gives a direct sum of matrix algebras over the selected complex field.  This is the standard route, and it consumes nothing beyond the stated hypotheses.  (Earlier presentations reached semisimplicity through a separately named ``radical pricing'' premise; the $C^*$ route makes that premise unnecessary, and it is withdrawn --- MD's exclusion of zero-cost directions survives only as the physical gloss on what the observational-equivalence quotient removes.)  Finally, the banked non-commutativity chain ($L_\Delta \to L_{\mathrm{blk}} \to L_\Pi \to T_{\mathrm{alg}}$, \S\ref{sec:nc_summary}) shows that $\Delta>0$ forces a nonzero commutator among tested defenders --- the rotated three-dimensional model of Proposition~\ref{prop:minimal_ijc_witness} exhibits the mechanism --- so the semisimple quotient is not wholly commutative.  This proposition records the proof boundary explicitly: without IJC one may remain in the Sep/Boolean branch, and without the positivity/finite-resolution hypotheses one does not obtain the finite semisimple carrier used by Skolem--Noether.
\end{proof}

\smallskip\noindent\textit{Scope of the proposition.}\quad
The conclusion is conditional, and two of its steps are scoping
clauses rather than closed arguments.  First, the argument treats a
nonzero pool as cost-active.  An interface whose pool carries no
positive joint surplus is the Sep branch, and whether a physical
interface occupies the IJC branch is an empirical occupancy read off
its records --- the $\Delta = 0$ falsifiability box above owns exactly
this.  Second, the step from ``not a direct sum of individual
defenders'' to ``mixes sectors'' is not established for arbitrary
operators: operators that are block-diagonal in some finer
decomposition yet lie outside the generated algebra exist, and the
proposition's conclusion concerns the algebra generated by the tested
defenders, not all of $\End(V)$.  The proposition therefore sits
inside the hybrid keystone framing of \S\ref{sec:nc_summary}: the
bridge is derived given the IJC occupancy; the occupancy itself is
empirical, not derived.

\begin{lemma}[Boolean refinement blocks the IJC branch]
\label{lem:boolean_refinement_blocks_ijc}
If all queried finite contexts at an interface admit a faithful common
Boolean refinement preserving the tested marginals and the defined partial
products, then the Sep/IJC obstruction is absent: the pool rotation angle
is $\theta=0$, the joint defender is the direct-sum defender, and the
noncommutative admissibility-algebra branch is not activated.  Conversely,
a nonzero pool rotation supplies an explicit gluing obstruction to such a
faithful Boolean refinement.
\end{lemma}

\begin{proof}
A faithful Boolean refinement assigns a single joint record algebra whose
atoms refine every queried context and whose restrictions reproduce the
observed marginals.  In the notation of Proposition~\ref{prop:minimal_ijc_witness},
that means the joint defender is represented on the direct sum
$M_1\oplus M_2$ with no pool-supported component.  Hence
$E_{12}=E_1+E_2$ and all context projections commute.  If instead the
minimum joint defender contains the vector
$w_\theta=\cos\theta e_1+\sin\theta p$ with $\theta\ne0$, the restriction
of the joint defender to the individual contexts still reproduces the
single-context marginals.  Suppose, for contradiction, that a faithful
Boolean refinement $\mathcal B$ also preserved the defined partial product.
Then both $E_1$ and $E_{12}$ would be sums of atoms in the same Boolean
algebra $\mathcal B$, hence would commute.  But Proposition~\ref{prop:minimal_ijc_witness}
gives $[E_1,E_{12}]=\cos\theta\sin\theta(|e_1\rangle\langle p|-|p\rangle\langle e_1|)\ne0$.
Therefore no such atom assignment exists.  This is precisely the finite
global-section obstruction familiar from sheaf-theoretic contextuality: the
single-context marginals exist, but they do not glue to a faithful global
Boolean record preserving the tested product.
\end{proof}

\smallskip\noindent\textbf{Gauge as continuation-profile-preserving relabeling.}\quad
In the continuation-calculus vocabulary of Paper~10, gauge transformations
are zero-marginal-cost relabelings that preserve all tested continuation
profiles.  Here this becomes a finite-interface statement: internal
automorphisms preserve the admissibility algebra's sector decomposition,
while kinematic automorphisms move interface labels and are outer.  The
Paper~9 gravitational gauge-reduction analogy is methodological only;
Remark~\ref{rem:p26_parallel} records the parallel without making it
load-bearing for Paper~2.


\subsection{The non-commutativity chain (summary)}
\label{sec:nc_summary}

The chain from superadditivity to non-commutativity has four links,
with home as noted per link.  Link~1 is established in this supplement
(Appendix~\ref{app:LDelta}); Links~2--4 are the constructive chain
developed in Paper~5 and certified in the bank, whose endpoint
(non-commutativity) is the classification theorem of Paper~1 Supplement
v8.41 \S``Sep/IJC representation theorem'' (the \emph{Sep iff commutative
record representation} main theorem).  We restate the key result
of each without re-proof.

\smallskip\noindent\textbf{Link 1: $L_\Delta$ (Superadditivity).}
For co-located distinctions $d_1, d_2$ at an interface $\Gamma$ with
non-trivial pool $\Pi \ne \{0\}$:
$\Delta(d_1, d_2) := \eps(\{d_1, d_2\}) - \eps(d_1) - \eps(d_2) > 0$.
\textit{Source:} Appendix~\ref{app:LDelta} (this supplement).  \textit{Inputs:} A1, K2, K3,
$T_{\mathrm{sep}}$, SP.

\smallskip\noindent\textbf{Link 2: $L_{\mathrm{blk}}$ (Direct-Sum Failure).}
If $\Delta > 0$, the minimum-cost joint defense cannot be a direct sum
of individual defenses: $E_{\{d_1,d_2\}} \ne E_{d_1} + E_{d_2}$.
\textit{Source:} Paper~5 (non-commutativity chain); bank \texttt{check\_L\_nc} (\texttt{core.py}).

\smallskip\noindent\textbf{Link 3: $L_\Pi$ (Codespace Rotation).}
The minimum-cost joint defender $\pi_{W_*}$ has a codespace $W_*$ rotated
into the pool $\Pi$ by angle $\theta \ne 0$.  The off-diagonal component
$F_\Pi := \pi_{W_*} - E_{d_1} - E_{d_2} \ne 0$ is a new generator
supported on the pool.
\textit{Source:} Paper~5; bank \texttt{check\_L\_Pi} (\texttt{core.py}).

\smallskip\noindent\textbf{Link 4: $T_{\mathrm{alg}}$ (Non-Commutativity).}
The commutator $[E_{d_1}, F_\Pi] \ne 0$ witnesses non-commutativity.
The algebra $\mathcal{A} = \mathrm{Alg}_\mathbb{R}(\{E_d\} \cup
\{E_{\{d_i,d_j\}}\})$ admits no faithful commutative
$*$-representation.
\textit{Source:} classification --- Paper~1 Supplement v8.41
\S``Sep/IJC representation theorem'' (\emph{Sep iff commutative record
representation}); construction --- Paper~5 + bank \texttt{check\_T\_alg}
(\texttt{core.py}).  \textit{Inputs:} $L_\Delta$,
$L_{\mathrm{blk}}$, $L_\Pi$.

\smallskip\noindent\textit{The chain is pre-Hilbert:} no GNS
construction, no spectral theory, no Hilbert space is invoked.
The non-commutativity is an algebraic fact about $\End(V)$,
established before $\mathcal{H}_\omega$ is constructed.



\subsection{A two-sector admissibility-algebra example}
\label{sec:two_sector_example}

\begin{tcolorbox}[apfbox, title={Worked example: pool, anchors, projections}]
\footnotesize
Let a finite interface have sector space
\[
  V=M_c\oplus M_w\oplus \Pi,
  \qquad
  M_c\cong\mathbb C^3,
  \quad M_w\cong\mathbb C^2,
  \quad \Pi\cong\mathbb C .
\]
Here $M_c$ and $M_w$ are anchor spaces for two tested distinctions, and
$\Pi$ is the residual pool.  The individual admissibility projections are
$E_c=\pi_{M_c}$ and $E_w=\pi_{M_w}$.  In the Sep branch the joint
defender is $E_c+E_w$.  In the IJC branch the least-cost joint defender
has a graph component into the pool, for example
$\pi_{\mathrm{graph}(K)}$ for some nonzero $K:M_c\to\Pi$, so the algebra
generated by $E_c,E_w,$ and the joint projection is noncommutative.
After complex carrier closure, the corresponding finite admissibility
algebra has the canonical block form
\[
  \mathcal A\cong M_3(\mathbb C)\oplus M_2(\mathbb C)\oplus\mathbb C .
\]
Its inner automorphism group is $\mathrm{PU}(3)\times\mathrm{PU}(2)$;
its block-center supplies a relative-phase torus.  The later results of
this section show how the field-module lift replaces the projective
observable action by $\SU(3)\times\SU(2)$ on matter fields and how the
relative torus is pruned to the single hypercharge direction.  This toy
model is not the Standard Model derivation; it is only the smallest
notation map for $M_d$, $E_d$, $\Pi$, block centers, and inner
automorphisms.
\end{tcolorbox}
\begin{remark}[Example scope]\label{ex:two_sector_admissibility}
The example deliberately omits generations, anomaly equations, and the
full 1{,}680-template scan.  Its role is definitional: it shows how the
nonstandard words ``pool,'' ``anchor,'' and ``admissibility projection'' are
represented by ordinary finite-dimensional subspaces and projections before
any classification theorem is invoked.
\end{remark}

\subsection{Gauge identification theorem}
\label{sec:gauge_ident}

The central question is: how is the gauge group identified with the
automorphism group of the admissibility algebra, and how are
spatial/kinematic automorphisms (e.g., diffeomorphisms) excluded?

The answer rests on a structural separation: the admissibility algebra
$\mathcal{A}$ acts on the \emph{internal} sector decomposition at a
single interface $\Gamma$, which is finite-dimensional and defined by
$T_{\mathrm{sep}}$.  Diffeomorphisms are \emph{inter-interface} maps
($L_{\mathrm{loc}}$) and therefore commute with the internal algebra
by construction.

\begin{definition}[Local support: the algebra and the localized unital maps]\label{def:local_support}
Two distinct objects are localized at an interface $\Gamma$, and they
must not be conflated.
\begin{enumerate}[label=(LS\arabic*),nosep]
  \item \textbf{The ambient corner algebra.}
  $\mathcal{A}_{\mathrm{corner}}(\Gamma) := \{A \in \mathrm{End}(V) :
  A(S_\Gamma) \subseteq S_\Gamma,\ A|_{S_\Gamma^\perp} = 0\}
  \;\cong\; \mathrm{End}(S_\Gamma) \oplus 0$.
  This is a linear $*$-subalgebra of $\mathrm{End}(V)$ (closed under
  addition, scalar multiplication, products, and adjoints), with unit
  the $B$-orthogonal projection $P_\Gamma$ onto $S_\Gamma$ (a
  \emph{local} unit; $\mathcal{A}_{\mathrm{corner}}(\Gamma)$ is
  non-unital as a subalgebra of $\mathrm{End}(V)$).  The corner is
  \emph{simple} ($\cong \mathrm{End}(S_\Gamma)$); it is the ambient
  home for local observables, not the admissibility algebra itself,
  which is the generated proper subalgebra
  $\mathcal{A}(\Gamma) \subseteq \mathcal{A}_{\mathrm{corner}}(\Gamma)$
  defined below.
  \item \textbf{The localized unital transformations.}
  $\mathcal{U}(\Gamma) := \{U \in \mathrm{End}(V) :
  U(S_\Gamma) \subseteq S_\Gamma,\ U|_{S_\Gamma^\perp} = \mathrm{id}\}$.
  This set is closed under composition but \emph{not} under addition
  ($(U + U')|_{S_\Gamma^\perp} = 2\,\mathrm{id}$): it is a monoid, not
  an algebra, and its invertible $B$-unitary elements form the
  \emph{localized unitary group} $U_\Gamma \oplus \mathrm{id}$.
  Localized symmetries live here; observables live in (LS1).
\end{enumerate}
Earlier versions used a single definition demanding
$A|_{S_\Gamma^\perp} = \mathrm{id}$ for all localized operations; that
class is not closed under addition and cannot serve as the observable
algebra.  Where a tensor factorization of $V$ exists, (LS1) plays the
role of the conventional $\mathcal{B}(H_\Gamma) \otimes I$; here the
substrate decomposition is a direct sum, so the correct algebra is the
corner $\mathrm{End}(S_\Gamma) \oplus 0$.
\end{definition}

\noindent\textit{Derivation of local support from FD3.}\quad
FD3 (anchor-set locality) states that the anchor set $M_d$ of any
distinction $d \in \mathcal{D}(\Gamma)$ lies within $S_\Gamma$.
The admissibility projection $E_d$ has $\mathrm{im}(E_d) = M_d
\subseteq S_\Gamma$ and $\ker(E_d) \supseteq S_\Gamma^\perp$ (by
$T_{\mathrm{adj}}$: $E_d$ is the $B$-orthogonal projection onto $M_d$,
hence $E_d|_{M_d^\perp} = 0$, and $S_\Gamma^\perp \subseteq M_d^\perp$).
Therefore $E_d(v) = v$ for $v \in M_d$ and $E_d(v) = 0$ for
$v \in S_\Gamma^\perp$: each generator annihilates the complement and
so satisfies (LS1).  The \emph{admissibility algebra}
$\mathcal{A}(\Gamma) = \mathrm{Alg}(\{E_d : d \in \mathcal{D}(\Gamma)\})$
is generated inside the corner $\mathcal{A}_{\mathrm{corner}}(\Gamma)
\cong \mathrm{End}(S_\Gamma) \oplus 0$, and
(LS1) \emph{is} closed under sums, scalar multiples, products, and
adjoints.  Therefore every $A \in \mathcal{A}(\Gamma)$ has local
support in the (LS1) sense.  The two algebras are distinct objects
carrying distinct symbols: $\mathcal{A}_{\mathrm{corner}}(\Gamma)$ is
the simple ambient corner, while $\mathcal{A}(\Gamma)$ is the
generated proper subalgebra, and it is the generated algebra --- with
its non-simple block structure --- that every downstream result (the
gauge identification, the exact sequence, EC/CM) consumes.  (The v3.9 text asserted that the $E_d$
act as the identity on the complement and satisfy the unital form of
local support; that was wrong --- zero and the identity are different
operators --- and is repaired by the two-object definition above.)

\begin{proposition}[Block structure of the generated admissibility algebra]
\label{prop:generated_block_structure}
The generated admissibility algebra
$\mathcal{A}(\Gamma) = \mathrm{Alg}(\{E_d\} \cup \{E_{\{d_i,d_j\}}\})$
is a finite-dimensional $*$-subalgebra of the simple corner
$\mathcal{A}_{\mathrm{corner}}(\Gamma)$, and it --- not the corner ---
carries the Wedderburn block decomposition consumed downstream:
$\mathcal{A}(\Gamma) \cong \bigoplus_k M_{n_k}(\mathbb{C})$, with the
Standard-Model specialization $\{n_k\} = \{3,2,1\}$, i.e.\
$M_3(\mathbb{C}) \oplus M_2(\mathbb{C}) \oplus \mathbb{C}$.
\end{proposition}

\begin{proof}
The generated algebra is a $*$-subalgebra of the corner by the FD3
derivation above (each generator annihilates $S_\Gamma^\perp$).
Semisimplicity of the generated algebra is supplied by
Proposition~\ref{prop:superadd_to_semisimple_v35}, under that
proposition's stated premises, by the standard finite-dimensional
$C^*$ route.  Wedderburn--Artin (T2a) with field selection
(T2c) then gives $\mathcal{A}(\Gamma) \cong \bigoplus_k
M_{n_k}(\mathbb{C})$.  The specialization $\{n_k\} = \{3,2,1\}$ is not
proved at this point of the document: it is delivered by Theorem~R's
carrier requirements together with the fermion scan, and enters
downstream as hypothesis H1 of
Theorem~7.7 of \TSI{}.  The corner, being simple,
admits no such decomposition; any statement attributing block
structure to $\mathcal{A}_{\mathrm{corner}}(\Gamma)$ would be false,
which is why the two symbols are kept apart.
\end{proof}

\begin{lemma}[Commutation from disjoint support]
\label{lem:loc_commut}
Let $\Gamma_1, \Gamma_2$ be interfaces with $B$-orthogonal substrates:
$B(u, v) = 0$ for all $u \in S_{\Gamma_1}$, $v \in S_{\Gamma_2}$
(from $L_\omega$ applied to disjoint interfaces via $L_{\mathrm{loc}}$).
Then:
\begin{enumerate}[label=(\roman*),nosep]
  \item (\emph{Observable algebras.})  For all
  $A \in \mathcal{A}_{\mathrm{corner}}(\Gamma_1)$ and
  $B \in \mathcal{A}_{\mathrm{corner}}(\Gamma_2)$
  (the (LS1) corners), $AB = BA = 0$; in particular $[A, B] = 0$.
  The generated admissibility algebras
  $\mathcal{A}(\Gamma_i) \subseteq
  \mathcal{A}_{\mathrm{corner}}(\Gamma_i)$ inherit the conclusion as
  subalgebras.
  \item (\emph{Localized unital transformations.})  For all
  $U \in \mathcal{U}(\Gamma_1)$ and $U' \in \mathcal{U}(\Gamma_2)$,
  $[U, U'] = 0$.
\end{enumerate}
\end{lemma}

\begin{proof}
Decompose $V = S_{\Gamma_1} \oplus S_{\Gamma_2} \oplus W$, where
$W = (S_{\Gamma_1} \oplus S_{\Gamma_2})^\perp$ (exists because $V$ is
finite-dimensional and $B$ is positive-definite).

(i)  By (LS1), $A = A_1 \oplus 0 \oplus 0$ and
$B = 0 \oplus B_2 \oplus 0$ with
$A_1 := A|_{S_{\Gamma_1}}$, $B_2 := B|_{S_{\Gamma_2}}$.
Then $AB$ annihilates $S_{\Gamma_1}$ and $W$ (where $B$ vanishes) and
annihilates $S_{\Gamma_2}$ (since $B$ maps it into $S_{\Gamma_2}$,
which $A$ annihilates); so $AB = 0$, and symmetrically $BA = 0$.

(ii)  By (LS2), $U = U_1 \oplus \mathrm{id}_{S_{\Gamma_2}} \oplus
\mathrm{id}_W$ and $U' = \mathrm{id}_{S_{\Gamma_1}} \oplus U'_2 \oplus
\mathrm{id}_W$.  For $v = v_1 + v_2 + w$:
\begin{align*}
  UU'(v) &= U(v_1 + U'_2 v_2 + w) = U_1 v_1 + U'_2 v_2 + w,\\
  U'U(v) &= U'(U_1 v_1 + v_2 + w) = U_1 v_1 + U'_2 v_2 + w.
\end{align*}
Therefore $UU' = U'U$.  Item (i) is the locality statement for
observables (Einstein-causality analogue); item (ii) is the statement
consumed by the net-automorphism and gauge arguments, which act by
localized unitaries.
\end{proof}

\begin{remark}[Two bilinear structures: $B$ and $g_{\mu\nu}$]%
\label{rem:two_bilinear}
There are two bilinear structures in the supplement and they live at
separate levels.  The form $B$ is the positive-definite admissibility-cost
kernel on the finite internal sector space
$S_\Gamma=\bigoplus_d M_d\oplus\Pi$ of a single interface.  By K2 and SP,
$B(v,v)>0$ for every nonzero internal perturbation.  It has no timelike
direction and is the form used for local support, sector orthogonality, and
Lemma~\ref{lem:loc_commut}.

The spacetime metric $g_{\mu\nu}$ is introduced later on the tangent space
of the continuum manifold.  Its spatial block inherits the positive cost
kernel, while the irreversible time direction contributes the negative
term:
\[
  g_{\mu\nu}dx^\mu dx^\nu=-N^2dt^2+g^{(\mathrm{sp})}_{ij}dx^idx^j,
  \qquad g^{(\mathrm{sp})}_{ij}\sim B_{ij}.
\]
Thus Lorentzian signature comes from adjoining the irreversible temporal
ordering to the positive spatial cost form, not from changing the sign of
$B$.  The finite-interface locality theorem therefore uses $B$; the later
spacetime reading of spacelike separation is a derived bridge, not an
input to the gauge-algebra argument.
\end{remark}

\noindent\textit{Relationship to AQFT.}\quad
In AQFT, Einstein causality ($[\mathcal{A}(\mathcal{O}_1),
\mathcal{A}(\mathcal{O}_2)] = 0$ for spacelike-separated regions) is
a separate axiom.  In the APF, it is a theorem: FD3 (anchor-set
locality) forces local support (Definition~\ref{def:local_support}),
which forces commutativity (Lemma~\ref{lem:loc_commut}).  FD3 is the
APF's analogue of Einstein causality, but it is a constitutive
definition of what ``admissibility at $\Gamma$'' means, not an
independent physical postulate.

\smallskip\noindent\textit{Scope of this result.}\quad
Lemma~\ref{lem:loc_commut} is a theorem \emph{internal to the
finite-capacity interface algebra}.  It is not claimed to reproduce
full AQFT locality in the continuum/type-III regime.  Specifically:
it applies to admissibility algebras $\mathcal{A}(\Gamma) =
M_n(\mathbb{C})$ at finite $n$, using the positive-definite cost
kernel $B$ for the orthogonal decomposition.  The continuum-limit
analogue (commutativity of type~III local algebras for
spacelike-separated regions) is a separate result that requires the
full machinery of AQFT (Haag--Kastler axioms, nuclearity estimates,
split property).  The finite-capacity result is the \emph{precursor}
from which the continuum result can be recovered in the
$C / \varepsilon^* \to \infty$ limit
(\S\ref{sec:finite_vs_typeIII}), but the two should not be
conflated.

\begin{lemma}[Internal--kinematic exact sequence]\label{lem:aut_decomp}
Let $\{\mathcal{A}(\Gamma)\}_{\Gamma \in \mathcal{I}}$ be a net of
finite-dimensional $*$-algebras indexed by a set $\mathcal{I}$, with
isotony $\mathcal{A}(\Gamma) \hookrightarrow \mathcal{A}(\Gamma')$ for
$\Gamma \subseteq \Gamma'$, and with independence: $\mathcal{A}(\Gamma)$
and $\mathcal{A}(\Gamma')$ commute when $\Gamma \cap \Gamma' = \emptyset$
(Lemma~\ref{lem:loc_commut}).  Every net automorphism $\alpha$ induces a
permutation $\pi(\alpha)$ of the index set $\mathcal{I}$ (preserving the
isotony order), and $\pi$ is a group homomorphism.  Let
$\mathrm{Aut}_{\mathrm{int}} := \ker(\pi)$ (automorphisms fixing
$\mathcal{I}$ pointwise, mapping each $\mathcal{A}(\Gamma)$ to itself)
and let $K := \mathrm{im}(\pi)$ be the induced group of admissible
index transformations.  Then the sequence
\[
  1 \;\longrightarrow\; \mathrm{Aut}_{\mathrm{int}}
  \;\longrightarrow\; \mathrm{Aut}(\text{net})
  \;\xrightarrow{\;\pi\;}\; K \;\longrightarrow\; 1
\]
is exact: $\mathrm{Aut}_{\mathrm{int}}$ is a normal subgroup and
$\mathrm{Aut}(\text{net})/\mathrm{Aut}_{\mathrm{int}} \cong K$.
\end{lemma}

\begin{proof}
$\pi$ is a homomorphism because index actions compose; its kernel is by
definition the fibrewise automorphisms, and kernels are normal.
Exactness at the ends is the definition of $K$ and injectivity of the
inclusion.
\end{proof}

\begin{remark}[No semidirect splitting is claimed]\label{rem:no_splitting}
Earlier versions asserted
$\mathrm{Aut}(\text{net}) = \mathrm{Aut}_{\mathrm{int}} \rtimes
\mathrm{Aut}_{\mathrm{kin}}$ with a \emph{unique} factorization of every
automorphism into internal and kinematic parts.  Both claims exceed what
the construction delivers: a semidirect decomposition requires a group
homomorphism $K \to \mathrm{Aut}(\text{net})$ splitting $\pi$, and no
such section is constructed here; and a putative ``kinematic part''
built from fiber isomorphisms extracted from $\alpha$ itself is not
determined by the index action alone, since different fiber
isomorphisms cover the same permutation --- the ambiguity is exactly a
coset of $\mathrm{Aut}_{\mathrm{int}}$.  Nothing downstream needs the
splitting: the gauge identification consumes only the kernel
$\mathrm{Aut}_{\mathrm{int}}$, which the exact sequence supplies.
\end{remark}

\begin{theorem}[Gauge Identification]\label{thm:gauge_ident}
Let $\mathcal{A}$ be the admissibility algebra at interface $\Gamma$
(from $T_{\mathrm{alg}}$), and let $\mathrm{Aut}(\mathcal{A})$ denote
its $*$-automorphism group.  Then:
\begin{enumerate}[label=(\roman*),nosep]
  \item The net automorphisms sit in the exact sequence of
        Lemma~\ref{lem:aut_decomp}:
        $1 \to \mathrm{Aut}_{\mathrm{int}} \to
        \mathrm{Aut}(\text{net}) \xrightarrow{\pi} K \to 1$, where
        $\mathrm{Aut}_{\mathrm{int}}$ (the kernel of the index action)
        consists of the automorphisms preserving the interface label
        $\Gamma$ and $K$ is the induced group of index transformations.
        No splitting of this sequence is claimed or needed
        (Remark~\ref{rem:no_splitting}).
  \item $\mathrm{Aut}_{\mathrm{int}}(\mathcal{A})$ contains the inner
        automorphism group $\mathrm{Inn}(\mathcal{A})$, and equals it
        blockwise.  By T2c ($\mathbb{F} = \mathbb{C}$) and T2a
        (Wedderburn--Artin), $\mathcal{A} \cong \bigoplus_{k=1}^{K}
        M_{n_k}(\mathbb{C})$ is \emph{semisimple} (central simple only
        when $K = 1$); Skolem--Noether applies per block, giving
        $\mathrm{Inn}(\mathcal{A}) \cong \prod_{k=1}^{K}
        \mathrm{PU}(n_k)$.
  \item The physical gauge group $G$ is identified with
        $\mathrm{Inn}(\mathcal{A})$.  Kinematic automorphisms
        (diffeomorphisms) project non-trivially to $K$ and are
        excluded from $G$.
  \item $G$ is non-abelian: $\mathrm{Inn}(\mathcal{A})$ is non-trivial
        and non-commutative because $\mathcal{A}$ is non-commutative
        ($T_{\mathrm{alg}}$) with at least one block of size
        $n_k \ge 2$.
\end{enumerate}
\end{theorem}

\begin{proof}

\textit{Step 1: Internal vs.\ kinematic automorphisms.}\quad
By Lemma~\ref{lem:aut_decomp} applied to the admissibility net
$\{\mathcal{A}(\Gamma)\}_{\Gamma \in \mathcal{I}}$ (with independence
from $L_{\mathrm{loc}}$), the index-action homomorphism $\pi$ places
every automorphism in the exact sequence
$1 \to \mathrm{Aut}_{\mathrm{int}} \to \mathrm{Aut}(\text{net})
\to K \to 1$.  The internal automorphisms are canonically the kernel
--- no choice of complement or splitting is involved.  Diffeomorphisms
act non-trivially on $\mathcal{I}$ and hence project non-trivially to
$K$; the gauge group is the kernel.  This gives (i).

\textit{Step 2: Wedderburn--Artin and Skolem--Noether.}\quad
After field selection (T2c: $\mathbb{F} = \mathbb{C}$), the admissibility
algebra is a finite-dimensional $*$-algebra over $\mathbb{C}$.
Wedderburn--Artin (T2a) gives
$\mathcal{A} \cong \bigoplus_{k=1}^{K} M_{n_k}(\mathbb{C})$.
By Skolem--Noether, every automorphism of $M_{n_k}(\mathbb{C})$ is inner:
$\alpha(a) = UaU^*$ for some unitary $U$.  The inner automorphism group
of each block is $\mathrm{Inn}(M_{n_k}) \cong \mathrm{PU}(n_k) =
\mathrm{U}(n_k) / \mathrm{U}(1)$.  For the full semisimple algebra:
$\mathrm{Inn}(\mathcal{A}) = \prod_{k=1}^K \mathrm{PU}(n_k)$.
This gives (ii).

\textit{Step 2b: Block permutations.}\quad
For $K \ge 2$, the full automorphism group of a finite-dimensional
semisimple algebra is:
\[
  \mathrm{Aut}\Bigl(\bigoplus_{k=1}^K M_{n_k}(\mathbb{C})\Bigr)
  \;=\;
  \Bigl(\prod_{k=1}^K \mathrm{PU}(n_k)\Bigr) \;\rtimes\; S_{\mathrm{iso}},
\]
where $S_{\mathrm{iso}} \subseteq S_K$ is the group of permutations among
isomorphic blocks (blocks with the same $n_k$).  The inner automorphisms
are $\mathrm{Inn}(\mathcal{A}) = \prod_k \mathrm{PU}(n_k)$; the block
permutations $S_{\mathrm{iso}}$ are outer.

The physical status of $S_{\mathrm{iso}}$:
\begin{enumerate}[label=(\alph*),nosep]
  \item Each Wedderburn block corresponds to a distinct admissibility sector
  under $T_{\mathrm{sep}}$.  Sectors are distinguished by their anchor sets
  $M_{d_k}$, which carry distinct admissibility costs $\eps(d_k)$.
  \item A block permutation $\sigma \in S_{\mathrm{iso}}$ exchanges sectors.
  This is an automorphism of the \emph{abstract} algebra but not of the
  \emph{cost-equipped} algebra $(\mathcal{A}, \eps)$: it maps
  $\eps(d_k) \mapsto \eps(d_{\sigma(k)})$, which differs from $\eps(d_k)$
  unless the permuted sectors are cost-degenerate.
  \item Cost-degenerate sectors (identical admissibility costs at every
  state), if any occurred, would make $S_{\mathrm{iso}}$ a nontrivial
  discrete relabeling of repeated summands.  For the derived algebra
  $M_3 \oplus M_2 \oplus \mathbb{C}$ no two summands are isomorphic, so
  $S_{\mathrm{iso}}$ is trivial and no such relabeling exists.  In
  either case $S_{\mathrm{iso}}$ is discrete, carries no connection
  one-form, and does not contribute to the gauge group $G$.
  Generations are not carried here: they are multiplicities within one
  representation type of the matter category
  (Definition~\ref{def:matter_category}; the copy clause of
  Definition~\ref{def:admissible_gauge}), and the generation count is
  the imported capacity--beta derivation
  (\S8 of \TSI{}), not a count of algebra
  blocks.
  \item Non-degenerate blocks: $\eps(d_k) \ne \eps(d_j)$ for $k \ne j$.
  The permutation $\sigma$ is not an automorphism of $(\mathcal{A}, \eps)$,
  so $S_{\mathrm{iso}}$ is trivial on the cost-equipped algebra.
\end{enumerate}

\noindent\textit{Conclusion.}\quad The continuous gauge group is
$G = \mathrm{Inn}(\mathcal{A}) = \prod_k \mathrm{PU}(n_k)$,
regardless of whether $S_{\mathrm{iso}}$ is trivial (non-degenerate
costs) or nontrivial (degenerate costs).  The reason is that
$S_{\mathrm{iso}}$ is \emph{discrete and outer}: it permutes
Wedderburn blocks but does not generate any continuous family of
automorphisms.  The continuous gauge group sees only the connected
component of $\mathrm{Aut}(\mathcal{A})$, which is
$\mathrm{Inn}(\mathcal{A})$.

\smallskip\noindent\textit{What survives of $S_{\mathrm{iso}}$.}\quad
$S_{\mathrm{iso}}$ is discrete and outer, generates no continuous
family of automorphisms, and contributes no gauge bosons, in any
regime.  For the derived algebra it is trivial outright:
$M_3 \oplus M_2 \oplus \mathbb{C}$ has no repeated summands, so there
are no cost-degenerate blocks to permute.  Flavor structure lives
elsewhere.  The observed generations are copies within one
representation type of the matter category --- multiplicity spaces,
not algebra summands (Definition~\ref{def:admissible_gauge}'s copy
clause) --- and their splitting by Yukawa couplings is Paper~4B's
territory.

\textit{Step 3: Identification.}\quad
The physical gauge group acts on the internal degrees of freedom at a
single interface, preserving the cost-equipped algebraic structure.  By
Steps~1--2, this is precisely $\mathrm{Inn}(\mathcal{A})$.  Kinematic
automorphisms are outer (they change the interface label) and are not
part of the gauge group.  Block permutations among isomorphic summands are outer and discrete,
and carry no generation content (the copy clause of
Definition~\ref{def:admissible_gauge}; \S8 of \TSI{}).
This gives (ii) and (iii).

\textit{Step 4: Non-abelian.}\quad
$\mathcal{A}$ is non-commutative ($T_{\mathrm{alg}}$).  A commutative
algebra has $\mathrm{Inn}(\mathcal{A}) = \{1\}$ (every inner
automorphism is trivial).  Since $\mathrm{Inn}(\mathcal{A})$ is
non-trivial, $G$ is non-trivial.  For $M_n(\mathbb{C})$ with $n \ge 2$,
$\mathrm{PU}(n)$ is non-abelian (e.g., $\mathrm{PU}(2) \cong
\mathrm{SO}(3)$).  This gives (iv).
\end{proof}


\begin{theorem}[Continuous APF gauge redundancies are inner; equality under full invariance]
\label{thm:continuous_gauge_inner_v34}
Let $\mathcal A(\Gamma)$ be a finite admissibility algebra equipped with
its interface support data, cost kernel $B$, anchor assignment $M$, and
admissible continuation profile $E$.  Write $\mathrm{Gau}_0(\Gamma)$ for
the identity component of the group of physical gauge redundancies
preserving these data (the identity component of
$\mathrm{Aut}_0(\mathcal A(\Gamma),B,M,E)$).  Then, unconditionally,
\[
  \mathrm{Gau}_0(\Gamma) \;\subseteq\; \mathrm{Inn}(\mathcal A(\Gamma))
  \;=\; \prod_k \mathrm{PU}(n_k),
\]
so every continuous physical gauge direction is generated by an inner
derivation of the finite admissibility algebra.  Equality,
$\mathrm{Gau}_0(\Gamma) = \mathrm{Inn}(\mathcal A(\Gamma))$, holds under
an explicit \emph{full-invariance hypothesis}: the anchor set, cost
kernel, and continuation data are $\mathrm{Inn}$-invariant as a family
--- conjugation by any unitary of $\mathcal A(\Gamma)$ maps the
preserved data onto itself.  The hypothesis is substantive, not
bookkeeping: a fixed anchor projector onto a proper sub-summand
subspace $M_d$ is preserved only by its stabilizer, which reduces the
group to a proper subgroup of that block's $\mathrm{PU}(n_k)$.  Outer
block permutations, including accidental permutations among isomorphic
blocks, are kinematic relabelings and cannot contribute gauge bosons.
\end{theorem}

\begin{proof}
Write the Wedderburn decomposition as
$\mathcal A=\bigoplus_k M_{n_k}(\mathbb C)$.  The automorphism group of the
abstract algebra is
\[
  \mathrm{Aut}(\mathcal A)\cong
  \Bigl(\prod_k \mathrm{PU}(n_k)\Bigr)\rtimes S_{\rm iso},
\]
where $S_{\rm iso}$ permutes isomorphic simple summands.  The first factor is
the identity component and has Lie algebra $\bigoplus_k\mathfrak{su}(n_k)$.
The second factor is finite and discrete, so it has no infinitesimal
generators.

\textit{The inclusion (unconditional).}\quad
A continuous physical gauge redundancy is in particular a
$*$-automorphism of $\mathcal A(\Gamma)$ in the identity component of
$\mathrm{Aut}(\mathcal A)$; the discrete factor $S_{\rm iso}$ is
separated from the identity, so the identity component is
$\prod_k \mathrm{PU}(n_k) = \mathrm{Inn}(\mathcal A(\Gamma))$ (Skolem--Noether
per block).  Hence $\mathrm{Gau}_0(\Gamma) \subseteq
\mathrm{Inn}(\mathcal A(\Gamma))$, with no hypothesis beyond the
finite complex semisimple carrier.

\textit{The reverse inclusion (conditional).}\quad
More is needed than summand preservation.  Inner automorphisms rotate
states within a fixed simple summand and preserve the \emph{summand's}
support projector and cost profile; but the preserved data list
includes the anchor assignment, and anchors are sub-summand objects.
For a generic $\mathrm{Ad}_U$ acting inside a block, a proper anchor
subspace $M_d$ is not mapped to itself --- it is preserved only by the
stabilizer of $M_d$.  The full-invariance hypothesis is exactly what
licenses the reverse inclusion: if the anchor set, cost kernel, and
continuation data are $\mathrm{Inn}$-invariant as a family, then every
inner automorphism preserves the data, so
$\mathrm{Inn}(\mathcal A(\Gamma)) \subseteq \mathrm{Gau}_0(\Gamma)$ and
equality holds.

\textit{Block permutations.}\quad
A nontrivial element of $S_{\rm iso}$ exchanges summand labels.  If the
exchanged summands are not identical as cost-equipped local codespaces,
support or cost preservation fails.  If they are identical, the
permutation remains a discrete symmetry of repeated copies; it is
separated from the identity component and has no connection one-form or
curvature.  It therefore behaves as a kinematic relabeling rather than
a continuous local gauge redundancy.

\textit{What downstream consumes.}\quad
The gauge-identification chain consumes the unconditional inclusion
together with the exact sequence (Lemma~\ref{lem:aut_decomp},
Proposition~\ref{prop:exact_seq}): the abelian-source and closure
arguments need every continuous gauge direction to be inner (the
inclusion) and the central structure of $\mathrm{U}(\mathcal A)$ (the
exact sequence).  No downstream result consumes the unconditional
equality.  Where the identification register $G =
\mathrm{Inn}(\mathcal A)$ is used (Theorem~\ref{thm:gauge_ident}), the
gauge-preserved data are read at block resolution --- summand
granularity --- which is the full-invariance hypothesis in the form the
physical criterion intends; the hypothesis boundary is stated here so
that the reader can see exactly where it enters.  If an additional
continuous symmetry is postulated that acts outside $\mathcal A$, it is
not derived by the APF gauge identification theorem and must pay for a
new admissibility record.
\end{proof}

\begin{remark}[Essential hypotheses for $G_{\rm int}=\mathrm{Inn}(\mathcal A)$]
\label{rem:inn_conditions_v35}
The conclusion uses five finite-algebra hypotheses: finite resolution, a faithful positive state/GNS representation, unitary implementation of zero-marginal-cost relabelings, complex-sector selection, and --- for the equality direction only --- the full-invariance hypothesis (anchor set, cost kernel, and continuation data $\mathrm{Inn}$-invariant as a family; without it only the inclusion $\mathrm{Gau}_0 \subseteq \mathrm{Inn}(\mathcal A)$ is available).  Relaxing them changes the theorem.  Extra commutative summands contribute central phases but no nonabelian inner automorphisms; real or quaternionic blocks move the problem to the field-selection gate rather than the gauge-identification theorem; nonfaithful states quotient out invisible summands and must be replaced by the faithful quotient before Skolem--Noether applies.  Thus the theorem is robust under small perturbations preserving the faithful finite complex $*$-algebra, but not under changing the algebraic category.
\end{remark}

\begin{definition}[The matter category]\label{def:matter_category}
A \emph{matter multiplet} is a finite-dimensional irreducible unitary
representation of the lifted gauge group
$\widetilde{G} = \prod_k \SU(n_k) \times \U(1)$ --- equivalently, a
choice of one irreducible representation per factor together with a
$\U(1)$ charge, including tensor products across factors
(bifundamentals such as $(\mathbf{3}, \mathbf{2})_{1/6}$).  The matter
category is an \emph{independent structural input} at this point of
the derivation: it is the category of representations of the group
that the algebra determines, not a category read off from the algebra's
own modules.
\end{definition}

\begin{remark}[Why left $\mathcal{A}$-modules cannot carry the matter]
\label{rem:no_left_modules}
Earlier versions took a matter multiplet to be an irreducible left
$\mathcal{A}$-module.  That category is too small: for
$\mathcal{A} = \bigoplus_k M_{n_k}(\mathbb{C})$ the central idempotents
act on any irreducible left module as scalar idempotents summing to
$1$, so exactly one acts as the identity (Schur), and every irreducible
left module factors through a \emph{single} block --- the irreducibles
are $\mathbb{C}^{n_1}, \ldots, \mathbb{C}^{n_K}$, never a simultaneous
bifundamental.  A quark doublet $(\mathbf{3}, \mathbf{2})$ therefore
cannot arise as an irreducible left module of
$M_3 \oplus M_2 \oplus \mathbb{C}$.  Carrying bifundamentals requires
commuting representations of the several factors on one space ---
equivalently, representations of the product group (the category of
Definition~\ref{def:matter_category}), a bimodule structure, or an
equivalent device.  The fermion scan of
Part~II of \TSI{} searches Definition~\ref{def:matter_category}'s
category; the algebra fixes the \emph{group}, and the group's
representation theory --- not the algebra's module theory --- carries
the matter.
\end{remark}

\begin{proposition}[Gauge automorphisms act on fields by the lifted matter representation]
\label{prop:gauge_acts_on_fields}
Let $(\mathcal{A},\omega)$ be the finite admissibility algebra with faithful
GNS representation $\rho_\omega:\mathcal{A}\to\mathcal{B}(\mathcal{H}_\omega)$,
and let $\Psi$ denote a matter multiplet carrying an irreducible unitary
representation of the lifted gauge group
$\widetilde{G} = \prod_k \SU(n_k) \times \U(1)$
(Proposition~\ref{prop:PU_to_SU}; the matter category is specified in
Definition~\ref{def:matter_category} below --- it is \emph{not} the
category of irreducible left $\mathcal{A}$-modules).  A
zero-marginal-cost internal relabeling preserving all tested
expectations is represented by a unitary normalizer
$U\in U(\mathcal{A})$.  The matter action itself is \emph{defined on
$\widetilde{G}$}: for $\tilde g \in \widetilde{G}$,
\[
  a\mapsto \pi(\tilde g)\,a\,\pi(\tilde g)^{*},\qquad
  \psi\mapsto \rho_{\Psi}(\tilde g)\,\psi ,
\]
where $\pi:\widetilde{G}\to \mathrm{Inn}(\mathcal{A})$ is the composite
of the blockwise covering maps $\SU(n_k)\twoheadrightarrow\mathrm{PU}(n_k)$
(the abelian coordinate acts trivially on observables), and
$\rho_\Psi$ --- the symbol is only ever applied to elements of
$\widetilde{G}$ --- is $\Psi$'s representation from the matter category.
Three clauses fix what the algebra side does and does not determine.
\begin{enumerate}[label=(\roman*),nosep]
  \item \emph{Observable side.}  A zero-marginal-cost relabeling acts
  on observables by conjugation $a \mapsto UaU^*$; after quotienting
  the central kernel this is an element of $\mathrm{Inn}(\mathcal{A})$
  (Proposition~\ref{prop:exact_seq}).
  \item \emph{Determination up to the central torus.}  The algebra
  side determines the corresponding matter transformation only up to
  the central data: conjugation data fix an element of
  $\mathrm{Inn}(\mathcal{A})$ and leave the action of
  $Z(\widetilde{G})$ --- in particular of the relative-phase torus
  $T_{\mathrm{rel}}$ --- undetermined.  No global homomorphic splitting
  $U(\mathcal{A}) \to \widetilde{G}$ is claimed: the blockwise
  phase--$\SU$ factorizations are multivalued (each determined only up
  to $\mathbb{Z}_{n_k}$), and no continuous global section exists.
  \item \emph{The charge representation closes the gap.}  Which
  one-parameter subgroup of $T_{\mathrm{rel}}$ acts on $\Psi$, and with
  what charge, is supplied by the charge representation $\rho_\Psi$
  itself --- a named input of the matter category, per the entry-typed
  inventory I-typing (Definition~\ref{def:matter_category};
  Definition~\ref{def:admissible_gauge}).
\end{enumerate}
Thus the algebraic
identification $G_{\mathrm{alg}}=\mathrm{Inn}(\mathcal{A})$ and the
physical gauge action on matter fields are two presentations of the same
internal symmetry --- adjoint/projective on observables, linear on
fields --- with the central (abelian) data carried by the matter
category's charge representation, not derived from the algebra side.
\end{proposition}

\begin{proof}
(i): Finite dimensionality and faithfulness imply that every tested
expectation has the form $\omega(b^*ab)$ for $a,b\in\mathcal A$.  A
relabeling that preserves all such expectations and does not move the
interface label is, by the finite-dimensional Kadison/Wigner form,
implemented by a unitary or antiunitary normalizer of the represented
algebra; the oriented complex carrier $B1'$ excludes the antiunitary
branch for continuous gauge transformations.  Conjugation by $U$ is the
map of Proposition~\ref{prop:exact_seq}; scalars in $Z(U(\mathcal A))$
act trivially by conjugation, so the induced automorphism group is
$\mathrm{Inn}(\mathcal A)$.

(ii): Exactness of $1 \to Z(U(\mathcal A)) \to U(\mathcal A) \to
\mathrm{Inn}(\mathcal A) \to 1$ says precisely that conjugation data
determine $U$ only up to $Z(U(\mathcal A))$; passing to the lifted
group, the corresponding matter transformation is determined only up to
the central torus's action.  A global continuous homomorphic section of
the blockwise phase--$\SU$ factorization would invert the $n_k$-fold
covers $\U(1)\times\SU(n_k) \to \U(n_k)$, which admit no continuous
global sections; the up-to-central-data statement is therefore the
strongest algebra-side conclusion available, and (ii) is sharp.

(iii): The matter category is an independent structural input
(Definition~\ref{def:matter_category}); $\rho_\Psi$ assigns each
multiplet its $\widetilde{G}$-representation, including its abelian
charge.  The action displayed in the statement is therefore well
defined on $\widetilde{G}$ by construction, and it descends to the
faithful quotient identified in Proposition~\ref{prop:global_structure}
because the derived template's charges satisfy the $\mathbb{Z}_6$
condition (Corollary~\ref{cor:fermion_Z6_compat}).  This is the
R-EC-inv $+$ I-typing conditionality in its lift form: the physicality
of the abelian direction is carried by the charged matter that sees it,
not by the algebra's conjugation action, and the surviving relative
phase is fixed downstream by Proposition~7.1 of \TSI{}.
\end{proof}

\begin{remark}[Methodological parallel: Paper~9 v1.2 gauge-reduction]\label{rem:p26_parallel}
The gauge-identification theorem follows a reusable APF audit template:
start with a larger automorphism/response group, quotient continuation-
profile-preserving relabelings, and name the residual obstruction that must
be closed.  Paper~9 applies the same template in a weak-field comparison
sector, reducing a local $GL^+(2,\mathbb R)$ clock-ruler response by gauge
equivalence and isolating the trace obstruction.  The analogy is useful for
wayfinding only.  Paper~2's gauge derivation depends on the finite
Wedderburn/Skolem--Noether chain above, not on the Paper~9 gravitational
normal-form program.
\end{remark}

\begin{remark}[Emergence of abelian $\U(1)$ factors]\label{rem:U1}
The inner automorphism group $\mathrm{Inn}(\mathcal{A}) = \prod_k
\mathrm{PU}(n_k)$ contains no abelian factors: $\mathrm{PU}(n)$ is
semisimple for $n \ge 2$, and trivial for $n = 1$.  This raises the
question: where does the $\U(1)_Y$ hypercharge factor of the Standard
Model originate?

The answer involves the relative-phase torus of the multi-block unitary
group, together with a no-redundancy theorem.
\end{remark}

The algebraic relationship between the groups acting on the algebra
(by conjugation) and on the Hilbert space (linearly) is captured by
an exact sequence.

\begin{proposition}[Exact sequence for gauge representations]%
\label{prop:exact_seq}
Let $\mathcal{A} = \bigoplus_{k=1}^K M_{n_k}(\mathbb{C})$ with
$\mathrm{U}(\mathcal{A}) = \prod_k \mathrm{U}(n_k)$.  Define
the conjugation map $\pi: \mathrm{U}(\mathcal{A}) \to
\mathrm{Aut}(\mathcal{A})$ by $\pi(U)(a) = UaU^*$.  Then:
\begin{enumerate}[label=(\roman*),nosep]
  \item The image of $\pi$ is $\mathrm{Inn}(\mathcal{A}) = \prod_k
  \mathrm{PU}(n_k)$.
  \item The kernel of $\pi$ is the center
  $Z(\mathrm{U}(\mathcal{A})) = \prod_k \mathrm{U}(1)_k$,
  consisting of block-diagonal phase matrices
  $(e^{i\theta_1}\mathbf{1}_{n_1}, \ldots,
  e^{i\theta_K}\mathbf{1}_{n_K})$.
  \item The sequence
  \[
    1 \;\longrightarrow\; Z(\mathrm{U}(\mathcal{A}))
    \;\longrightarrow\; \mathrm{U}(\mathcal{A})
    \;\xrightarrow{\;\pi\;}\;
    \mathrm{Inn}(\mathcal{A}) \;\longrightarrow\; 1
  \]
  is exact.
  \item The overall diagonal phase
  $\mathrm{U}(1)_{\mathrm{diag}} = \{(e^{i\theta}\mathbf{1}_{n_1},
  \ldots, e^{i\theta}\mathbf{1}_{n_K}) : \theta \in \mathbb{R}\}$
  acts trivially on both $\mathcal{A}$ (by conjugation) and on
  $\mathcal{H}_\omega$ (as a global phase).  The quotient
  $Z / \mathrm{U}(1)_{\mathrm{diag}}$ is a torus of rank $K - 1$, the
  relative-phase torus $T_{\mathrm{rel}}$.
  \item The algebra's GNS carrier $\mathcal{H}_\omega$ transforms
  under $\mathrm{U}(\mathcal{A})$ linearly (each block entering with
  its left-regular multiplicity: $n_k$ copies of the fundamental
  $\mathbb{C}^{n_k}$), not under $\mathrm{Inn}(\mathcal{A})$ (which
  acts on $\mathcal{A}$ by conjugation, i.e., the adjoint
  representation).  Matter fields do \emph{not} live in
  $\mathcal{H}_\omega$: they carry the independently specified matter
  category (Definition~\ref{def:matter_category}); the GNS space is
  the algebra-side bookkeeping object in which the fundamental of each
  block is exhibited as a subrepresentation.
  The two actions agree on the Lie algebra
  ($\mathfrak{u}(\mathcal{A}) \to
  \mathfrak{inn}(\mathcal{A})$ is surjective with kernel
  $\mathfrak{z}$), but differ globally: $\mathrm{U}(\mathcal{A})$
  admits representations that $\mathrm{Inn}(\mathcal{A})$ does not.
\end{enumerate}
\end{proposition}

\begin{proof}
(i) By Skolem--Noether, every automorphism of $M_n(\mathbb{C})$ has
the form $a \mapsto UaU^*$ for some $U \in \mathrm{U}(n)$.  For the
semisimple algebra, each block is mapped to itself by an inner
automorphism (outer block permutations are handled separately in
Step~2b of Theorem~\ref{thm:gauge_ident}), so
$\mathrm{Im}(\pi) = \prod_k \mathrm{Inn}(M_{n_k}) = \prod_k
\mathrm{PU}(n_k)$.  (ii)~$\pi(U) = \mathrm{id}$ iff
$UaU^* = a$ for all $a \in \mathcal{A}$, iff $U$ is central.
For $M_n(\mathbb{C})$, the center is $\mathrm{U}(1) \cdot
\mathbf{1}_n$; for the direct sum, $Z = \prod_k \mathrm{U}(1)_k$.
(iii)~Surjectivity of $\pi$ onto $\mathrm{Inn}(\mathcal{A})$ is (i);
the kernel statement is (ii); the sequence is exact at
$\mathrm{U}(\mathcal{A})$ by construction.  (iv)~The diagonal
$\mathrm{U}(1)_{\mathrm{diag}}$ is a subgroup of $Z$ with
$K$-dimensional $Z$ quotient of dimension $K - 1$.  Triviality on
$\mathcal{H}_\omega$ is because states are defined up to global
phase: $|\psi\rangle$ and $e^{i\theta}|\psi\rangle$ represent the
same physical state.  (v)~The GNS representation maps each
$M_{n_k}(\mathbb{C})$ block to operators on $\mathbb{C}^{n_k}
\subset \mathcal{H}_\omega$; the natural group action on
$\mathbb{C}^{n_k}$ is $v \mapsto Uv$ (fundamental), not
$a \mapsto UaU^*$ (adjoint).  The fundamental is a linear
representation of $\mathrm{U}(n_k)$ but only a projective
representation of $\mathrm{PU}(n_k)$.
\end{proof}

\begin{definition}[Admissible gauge structure]\label{def:admissible_gauge}
Let $G$ be a compact Lie group acting faithfully on a finite-dimensional
Hilbert space $\mathcal{H}$, and let $\mathcal{M} = \{r_1, \ldots, r_N\}$
be the set of physically distinct irreducible multiplet types occurring in
$\mathcal{H}$.  \emph{Type} means an isomorphism class of
$\widetilde{G}$-representations, not an individual copy: replicated
copies of one representation (the observed generations are the
physical example) are one type with multiplicity, and EC does
\emph{not} demand that gauge labels separate copies within a type.
That generation identity is carried outside the gauge labels
(by mass and mixing structure) is a scoping \emph{premise} of this
definition, stated here explicitly rather than derived.  The copy
clause is scoped to the \emph{full realized} $\widetilde{G}$: types
are isomorphism classes of $\widetilde{G}$-representations, and it is
this inventory --- with $u^c$ and $d^c$ antecedently distinct
entry-types (the entry-typed inventory I-typing, a named input
consumed at Theorem~7.7 of \TSI{} H3; see
Lemma~\ref{lem:EC_CM_from_A1}) --- that CM's restriction tests act
on.  Applied instead to a proper subgroup $H \subsetneq \widetilde{G}$,
the clause would recount $u^c/d^c$ as one type at multiplicity two and
void the restriction test; the scoping forbids that reading.

\smallskip\noindent\textbf{Enforcement completeness (EC).}\quad
$G$ is \emph{admissibility-complete} if the map
$r_i \mapsto [r_i]_G$ from multiplet types to $G$-representation labels
is injective: distinct physical types carry distinct gauge quantum numbers.

\smallskip\noindent\textbf{Capacity minimality (CM).}\quad
$G$ is \emph{capacity-minimal} if no generator channel of its
enforcing configuration can be removed while preserving admissibility
completeness.  Removal is an operation on channel configurations, not
on groups; the group-level comparison space is the closed proper
subgroups of $G$ (with their restricted faithful representations).
Formally: for every closed connected subgroup $H \subsetneq G$ enforced
by a configuration with at least one generator channel removed, the
label map $r_i \mapsto [r_i]_H$ is \emph{not} injective.

\smallskip\noindent\textbf{Admissibility.}\quad
$G$ is \emph{admissible} if it is both admissibility-complete and
capacity-minimal.
\end{definition}

Before deriving EC and CM, we establish the bridge from
Lie-algebra generators to admissibility channels, which CM requires.
The bridge is not as simple as ``one generator = one channel'': the
relationship between Lie-algebra generators and admissibility channels
involves incompatibility (non-commuting generators cannot be
simultaneously monitored), coarse-graining (degenerate eigenvalues
reduce the information content of a single generator), and the
distinction between Cartan (commuting) and root (non-commuting)
generators.  We treat these systematically.

\begin{definition}[Enforcement monitoring of a generator]%
\label{def:monitoring}
Let $g_a \in \mathfrak{g}$ be a Hermitian generator of a compact Lie
algebra acting on a finite-dimensional Hilbert space $\mathcal{H}$ of
dimension $n$.  The \emph{admissibility monitoring} of $g_a$ is the
projection-valued measure (PVM) $\{P_{a,\lambda}\}_\lambda$ on the
spectrum $\sigma(g_a) \subseteq \mathbb{R}$, where $P_{a,\lambda}$
projects onto the $\lambda$-eigenspace of $g_a$.  The \emph{distinction
count} of $g_a$ is $d(g_a) := |\sigma(g_a)| - 1$ --- one less than
the number of distinct eigenvalues.  This is the number of binary
questions needed to determine which eigenspace a state occupies
(using a binary search on the ordered eigenvalue list).
\end{definition}

\begin{lemma}[Generator cost bound by resolution rank]\label{lem:gen_channel}
Let $G$ be a compact Lie group of the class fixed by the gauge
identification --- $\mathrm{Inn}(\mathcal{A}) = \prod_k \mathrm{PU}(n_k)$
and its closed subgroups (closed subgroups of a Lie group are Lie, by
Cartan's closed-subgroup theorem) --- acting faithfully and
non-trivially on the admissibility structure at $\Gamma$, with
$\dim(G) = p$.  Then the admissibility cost of \emph{enforcing}
$G$-equivariance is at least $p \cdot \varepsilon^*$:
\[
  C(G) \;\ge\; \dim(G)\cdot\varepsilon^* .
\]
This is the paper-side statement of the banked gauge clause
$n(G) = \dim(G)$, $C(G) = \dim(G)\cdot\varepsilon$
\coderef{check\_L\_cost\_gauge}{core.py}; the two independent proof
routes below are the bank's.
\end{lemma}

\begin{proof}
\textit{Route A (resolution rank; primary).}\quad
Enforcing $G$-equivariance requires the admissibility structure to
distinguish automorphisms that act differently on observables: if
$\alpha_{g_1}(A) \ne \alpha_{g_2}(A)$ for some observable $A$ and the
structure cannot separate $g_1$ from $g_2$, then only a quotient of
$G$ is enforced, not $G$ (orbit separation).  $G$ is a $p$-dimensional
manifold: every group element in a neighborhood of the identity is
specified by $p$ real coordinates, and by invariance of domain
(Brouwer, local form: a continuous injection from an open
$U \subseteq \mathbb{R}^p$ into $\mathbb{R}^k$ forces $k \ge p$), no
continuous injective resolution of that neighborhood uses fewer than
$p$ independent distinctions.  This step consumes a named premise,
\emph{operational encoding}: the enforced group action is registered
by a continuous \emph{injective} encoding of an identity neighborhood
into the admissibility structure's distinction coordinates ---
invariance of domain applies to that encoding, not to the abstract
group manifold.  The premise is motivated by orbit separation (a
relabeling the structure cannot separate from a neighbor is not
enforced as part of $G$), and it is a premise of this lemma, not a
theorem of the framework.  Resolving the enforced group therefore
requires at least $p$ independent admissibility distinctions, each an
irreducible channel costing at least $\varepsilon^*$
($L_{\varepsilon^*}$).  The count is a property of the group manifold
--- it is invariant under change of Lie-algebra basis, and it does not
assert that any particular generator corresponds to a separate
physical measurement.

\textit{Route B (bracket non-closure; confirmatory).}\quad
The bracket $[T_a, T_b]$ is realized by composition (four
exponentials), and composition is non-free ($L_{\mathrm{nc}}$:
interaction cost $\ge 0$, generically positive).  Each
bracket-generated direction therefore costs at least $\varepsilon^*$
($L_{\varepsilon^*}$); after closure all $\dim(G)$ directions are
populated, each costing at least $\varepsilon^*$.  Total
$\ge \dim(G)\cdot\varepsilon^*$.

Both routes give a resolution count of at least $\dim(G)$, hence the
lower bound of the statement: $C(G) \ge \dim(G)\cdot\varepsilon^*$ ---
in $L_{\mathrm{Cauchy}}$'s unit normalization, $C(G) \ge \dim(G)$ in
$\varepsilon^*$ units.  Equality holds under the disjoint-anchor
realizability premise; the composed operational statement below is
where the equality lives.
\end{proof}

\begin{remark}[The withdrawn measurement-channel argument]
\label{rem:withdrawn_channel_argument}
Earlier versions argued this bound by assigning one \emph{operational
measurement channel} per Lie-algebra basis generator, claiming that
commuting linearly independent generators require separate spectral
interfaces because no basis element can be a function of another.
That argument is withdrawn: it is false.  For
$g_1 = \operatorname{diag}(1, 0, -1)$ and
$g_2 = \operatorname{diag}(0, 1, -1)$, the operators commute and are
linearly independent, yet $g_2$ is a polynomial in the nondegenerate
$g_1$, and a single rank-one projective measurement in the common
eigenbasis determines both eigenvalues --- linear independence in the
Lie algebra does not imply separate measurement channels, and
``one generator, one channel'' is not invariant under change of
generator basis.  The corrected route above counts \emph{resolution
rank of the enforced group manifold}, not measurements on states: an
informationally complete measurement may characterize many
observables, but it cannot reduce the number of independent
distinctions needed to resolve a $p$-dimensional group of enforced
relabelings below $p$.  The bank's statement
(\codeid{check\_L\_cost\_gauge}) has carried the resolution-rank route
since v24.3.404; the withdrawn presentation was this document's, not
the bank's.
\end{remark}

\smallskip\noindent\textit{Scope and caveats.}\quad
The lemma establishes a \emph{lower bound} on the admissibility cost,
not an exact cost: co-located non-abelian channels may incur
superadditive surplus ($L_\Delta$), making the actual cost strictly
greater.  The lower bound suffices for CM: a redundant generator
wastes \emph{at least} $\varepsilon^*$ capacity, which A2 then
eliminates (Lemma~\ref{lem:EC_CM_from_A1}).  The scope rider travels
with the bank statement: the argument holds for
$\mathrm{Inn}(\mathcal{A})$-class groups and their closed subgroups,
which is exactly the class the classification of
Part~II of \TSI{} ranges over.

\smallskip\noindent\textit{The composed operational statement
$C(G) = \dim(G)$.}\quad
The equality the cost ranking consumes is a composition of two named
operational premises, stated once here.  The \emph{lower} bound is the
lemma above, a theorem on the operational-encoding premise (invariance
of domain applied to a continuous injective encoding of an identity
neighborhood).  The \emph{upper} bound rests on the disjoint-anchor
realizability premise: the $\dim(G)$ resolution distinctions can be
anchored on disjoint supports, so C1-additivity applies and no
superadditive surplus accrues across generator channels.  Under the
pair, $C(G) = \dim(G)\cdot\varepsilon^*$ exactly.  The pair is the
operational cost premise of the gauge-cost program: a cost functor
deriving both premises from the admissibility semantics is the named
open lane, and until it closes, every consumer of the equality --- the
Lie-algebra cost ranking and the classification theorem's cost
convention --- carries the pair by name.  The paper-side realization of
the composed statement, with the same two premises named, is
Proposition~4.1 of \TSI{}\coderef{check\_L\_cost\_gauge}{core.py}.


\begin{lemma}[EC and CM from the admissibility semantics]\label{lem:EC_CM_from_A1}
EC and CM are derived principles of the admissibility framework.
They are not consequences of A1 (finite capacity) alone; their
dependencies on the two axioms (A1, A2) and the regularity input MD
are stated cleanly under the dual-axiom framing of Paper~1, \S2:
\begin{enumerate}[label=(\roman*),nosep]
  \item \textbf{EC} requires: A1 (finite capacity) $+$
  $T_{\mathrm{sep}}$ (sector decomposition) $+$
  $L_{\mathrm{loc}}$ (locality) $+$ FD3 (anchor-set locality) $+$
  the gauge identification (Theorem~\ref{thm:gauge_ident}).
  EC is the statement that the internal gauge labels at a single
  interface must resolve all physically distinct internal types.
  It is a \emph{completeness principle for the interface-local
  admissibility semantics}.  Its structural-separation ingredient is
  A1-derived, but EC itself rests on two named, non-derived elements:
  the \emph{kinematic type-inventory reading} R-EC-inv (at a single
  interface, kinematic type-distinctness is exhausted by the
  interface's counted labels; label-identical species are copies, not
  types --- A1-motivated via the admissibility referent, not
  A1-derived) and, where EC is consumed against a fixed template, the
  \emph{entry-typed inventory} I-typing (the template's entries are
  antecedently distinct types: $u^c \ne d^c$ prior to any abelian
  label).  A2 plays no role in EC.
  \item \textbf{CM} requires: A2 (minimum cost) applied to the
  admissible set defined by A1 $+$ MD (via $L_{\varepsilon^*}$) $+$
  the generator--channel bridge (Lemma~\ref{lem:gen_channel}).
  A configuration carrying a redundant generator is \emph{not}
  inadmissible under A1 --- it still satisfies $\Sigma\varepsilon \le C$.
  But A1 $+$ MD (via $L_{\varepsilon^*}$) guarantees that removing the
  redundant generator yields a strictly cheaper configuration in the
  same admissible set, by $\varepsilon^* > 0$.  A2 then selects the
  cheaper one as the realised gauge structure.  CM is therefore the
  statement that the realised admissible structure carries no
  redundant generators --- a direct consequence of A2.
\end{enumerate}
\end{lemma}

\begin{proof}
\textit{EC from the structural separation.}\quad
The argument for EC requires a premise that is not A1 alone but is
\emph{derived from} A1: the internal--external structural separation.
By $T_{\mathrm{sep}}$ (Paper~1 Supplement), the admissibility structure
at each interface $\Gamma$ decomposes into sectors.  By
Definition~\ref{def:local_support} and Lemma~\ref{lem:loc_commut}
(both derived from FD3 and $L_{\mathrm{loc}}$, which are themselves
derived from A1): at a single interface, the admissibility algebra
$\mathcal{A}(\Gamma)$ acts on the \emph{internal} sector decomposition
of $S_\Gamma$, while inter-interface (spacetime) distinctions
--- spatial separation, momentum, mass --- are properties of the
\emph{external} structure (the inter-interface map
$L_{\mathrm{loc}}$).  At a single interface, the only admissibility
operations available are those in $\mathcal{A}(\Gamma)$, whose
automorphism group is the gauge group $G$
(Theorem~\ref{thm:gauge_ident}).

Now suppose two physically distinct multiplet types $r_i \ne r_j$
carry identical gauge quantum numbers: $[r_i]_G = [r_j]_G$.  Since
$G = \mathrm{Aut}(\mathcal{A}(\Gamma))$ and the $G$-representation
labels exhaust the distinctions available within
$\mathcal{A}(\Gamma)$, no admissibility operation at $\Gamma$
distinguishes $r_i$ from $r_j$.  (The exhaustion step just used is,
within the algebra-module vocabulary, a triviality --- isomorphic
modules are internally indistinguishable; the \emph{physical} claim,
that the algebra's counted module structure exhausts kinematic
type-distinctness at the interface, is the reading R-EC-inv, consumed
exactly here.)  But A1 requires that all physically
distinct states be admissible (this is the operational content of
``admissibility capacity'': the budget $C(\Gamma)$ is the capacity to
enforce distinctions at $\Gamma$).  If two types are indistinguishable
at every interface at which they co-occur, they are admissibility-equivalent
and therefore physically identical --- contradicting $r_i \ne r_j$.

The objection ``perhaps they are distinguished by mass or momentum
rather than gauge quantum numbers'' is addressed by the structural
separation: mass and momentum are spacetime (external) properties
resolved by the inter-interface structure, not by the admissibility
algebra at a single interface.  EC is the statement that the
\emph{internal} sector must resolve all internal types; it does not
claim that gauge quantum numbers are the only physically meaningful
labels, only that they must be complete within the internal sector.

\textit{CM from A2 over the admissible set.}\quad
CM requires three ingredients in addition to A2:
\begin{enumerate}[label=(\alph*),nosep]
  \item $L_{\varepsilon^*}$: each irreducible distinction costs at
  least $\varepsilon^* > 0$.  This is derived from A1 $+$ MD
  (admissibility isotropy) in the Paper~1 Supplement.  It is
  \emph{not} derivable from A1 alone: without the isotropy assumption
  MD, one cannot exclude the possibility that some channels have
  arbitrarily small cost.
  \item Lemma~\ref{lem:gen_channel}: resolving the enforced group
  costs at least one channel per group-manifold dimension
  ($n(G) = \dim(G)$, the resolution-rank route; the banked gauge
  clause \codeid{check\_L\_cost\_gauge}).  This is derived from
  faithfulness and non-triviality of the $G$-action (which follow from
  Theorem~\ref{thm:gauge_ident}, itself derived from A1) and invariance
  of domain on the group manifold --- not from any
  one-generator-one-measurement claim
  (Remark~\ref{rem:withdrawn_channel_argument}).
  \item A1 $+$ MD jointly define the admissible set
  $\mathcal{A}(\rho, R)$ of \emph{channel configurations} satisfying
  $\Sigma\varepsilon \le C$ with $\varepsilon(d) \ge \mu^* n(d)$.  A
  channel configuration is a finite set $\mathcal{C}$ of costed
  generator-channels; the group it enforces is the closed subgroup
  $G = \langle \mathcal{C} \rangle$ generated by the configuration's
  channels, with its faithful representation on the interface carrier.
\end{enumerate}
Given (a)--(c): CM is stated at the channel-configuration level (the
semantics of the admissible set above).  Consider any admissible
configuration $\mathcal{C}$ containing a channel $X$ whose removal
preserves injectivity of the label map
$r_i \mapsto [r_i]_{\langle \mathcal{C} \setminus \{X\}\rangle}$ ---
a channel that does not refine any multiplet distinction.  The
configuration $\mathcal{C}' := \mathcal{C} \setminus \{X\}$ is a
well-defined set-theoretic removal of one costed channel.  Removal is
defined on configurations, never on groups: where a group-level
comparison is needed, the comparison space is closed subgroups,
quotients, and central-factor removals with specified faithful
representations, and the group enforced by $\mathcal{C}'$ is
$\langle \mathcal{C} \setminus \{X\}\rangle$, a closed subgroup of
$\langle \mathcal{C} \rangle$ (Cartan).  $\mathcal{C}'$ is admissible
(it satisfies the same capacity bound and isotropy condition; removing
a channel only relaxes them).  By (a) and (b),
$\Sigma_{d \in \mathcal{C}} \varepsilon(d) - \Sigma_{d \in \mathcal{C}'}
\varepsilon(d) = \varepsilon(X) \ge \varepsilon^* > 0$, so
$\mathcal{C}'$ is strictly cheaper than $\mathcal{C}$ inside the same
admissible set.  By A2 (minimum cost), the realised configuration
$\mathcal{C}_{\mathrm{realized}}(\rho, R)
= \arg\min_{\mathcal{C} \in \mathcal{A}(\rho, R)} \Sigma\varepsilon(d)$
cannot contain $X$.  Iterating over all redundant channels: the
realised admissible structure contains no redundant generators.
\end{proof}

\noindent\textit{Status of the premises: derived, not additional.}\quad
A careful reading of Lemma~\ref{lem:EC_CM_from_A1} might suggest that
EC and CM import ``extra'' assumptions beyond A1.  This would be a
genuine vulnerability if the required premises were independently
postulated.  In fact, every premise in the dependency chain is derived:
\begin{center}\small
\renewcommand{\arraystretch}{1.15}
\begin{tabular}{@{}p{3.0cm}p{3.5cm}p{3.5cm}p{2.5cm}@{}}
\toprule
\textbf{Premise} & \textbf{Content} & \textbf{Derived from} & \textbf{Reference} \\
\midrule
$T_{\mathrm{sep}}$ & Sector decomposition of $S_\Gamma$ &
  A1 $+$ $T_{\mathrm{form}}$ & P1S~\S2 \\
$L_{\mathrm{loc}}$ & Independent interfaces &
  A1 $+$ FD1--FD2 & P1S~\S4 \\
FD3 & Anchor-set locality &
  Foundational definition & P1S~\S1 \\
Gauge ident. & $G = \mathrm{Inn}(\mathcal{A})$ &
  $T_{\mathrm{alg}}$ $+$ T2a $+$ Skolem--Noether & This supplement \\
$L_{\varepsilon^*}$ & Cost floor $\varepsilon^* > 0$ &
  A1 $+$ MD & P1S~\S3 \\
Gen.-channel bridge & $\dim(\mathfrak{g})$ channels &
  Faithfulness $+$ $T_{\mathrm{form}}$ & This supplement \\
A2 & Min-cost selection over $\mathcal{A}(\rho, R)$ &
  Second axiom (P1, \S2) & P1, \S2; P1S~\S2 \\
\bottomrule
\end{tabular}
\end{center}
\noindent The only external inputs beyond the two named axioms (A1, A2)
are the regularity condition MD (admissibility isotropy), which enters
through $L_{\varepsilon^*}$, and FD3 (anchor-set locality), a
foundational definition rather than an empirical assumption.  The
remaining premises are mathematical consequences of A1 (or A1 $+$ MD)
and the upstream chain.  The dependency for CM in the table now reads
A1 $+$ MD $+$ $L_{\varepsilon^*}$ $+$ generator--channel bridge (which
together define the admissible set with cost floor) $+$ A2 (which
selects the min-cost member).

The correct characterization is therefore: \emph{the structural
separation that EC rests on is A1-derived, but EC itself is not a
consequence of A1 alone.  Two of its elements are named rather than
derived: the inventory-exhaustiveness step --- the claim that the
counted labels exhaust the kinematic distinctions available at the
interface, i.e.\ the reading R-EC-inv --- and the type/copy scoping
premise of Definition~\ref{def:admissible_gauge} (with its
entry-typed inventory I-typing where a fixed template is consumed).
CM follows from A2
applied to the admissible set defined by A1 $+$ MD: A1 bounds total
cost and admits every configuration satisfying that bound; MD gives
the positive cost floor (via $L_{\varepsilon^*}$) that makes a
redundant generator strictly costly; A2 then selects the minimum-cost
member of the admissible set, which by construction contains no
redundant generators.}  That the named elements are genuinely
non-derived is witnessed by a finite counter-construction: two
interfaces with independent budgets on which the species-to-label map
is non-injective while every distinction is enforced and priced within
the shapes of the banked premises --- A1 $+$ MD $+$ $T_{\mathrm{sep}}$
$+$ $L_{\mathrm{loc}}$ do not entail EC-injectivity over
species\coderef{check\_L\_EC\_inventory\_reading}{ec\_inventory\_reading.py}.
EC's honest grade is therefore
[P\textsubscript{structural\_reading} $\mid$ R-EC-inv $+$ I-typing];
CM's is unconditional given A2 over the A1 $+$ MD admissible set.
The earlier-version claim ``EC and CM follow
from A1'' was an overclaim on both counts: EC consumes the named
reading and input above, and
CM is properly an A2 result (operating on the A1 $+$ MD-defined
admissible set).  Neither correction imports a modeling assumption
beyond the framework's stated axioms and the two named elements: A2 is
one of the two axioms
named in the Paper~1, \S2 box, and MD is the regularity input also
listed in the Paper~1 axiom inventory.

\smallskip\noindent\textit{CM is not a selection principle (descriptive
reading).}\quad
Capacity minimality is sometimes described as ``Occam's razor for gauge
structure,'' but this risks making it sound like an aesthetic preference
or, worse, a dynamical mechanism that ranges over candidate structures
and rejects the non-minimal ones.  Neither reading is correct.
Lemma~\ref{lem:EC_CM_from_A1} establishes CM as a \emph{consequence
of A2} applied to the admissibility set defined by A1 $+$ MD: a
configuration with a redundant generator is admissible (it satisfies
the capacity bound and isotropy condition), but it sits at strictly
higher total cost than the same configuration with the redundant
generator removed (also admissible, by $\varepsilon^* > 0$).  A2
selects the cheaper one as the realised structure.  The
generator--channel bridge (Lemma~\ref{lem:gen_channel}) makes the
gap quantitative: each redundant generator costs exactly
$\varepsilon^*$.  Read correctly, CM is therefore a \emph{descriptive
label} on the realised admissible structure, not an algorithm: the
realised configuration is the $\arg\min_G \dim G$ subject to EC, and
that argmin is exactly what A2 fixes.  No dynamical search is
performed and none is needed; the framework reads off admissibility
space's argmin (A1 admits, A2 selects).

\noindent\textbf{Physical gauge group.}\quad By the gauge identification
theorem (Theorem~\ref{thm:gauge_ident}) and
Proposition~7.1 of \TSI{}, the Lie algebra of the physical gauge
group is $\mathfrak{su}(3) \oplus \mathfrak{su}(2) \oplus
\mathfrak{u}(1)_Y$.  At the global (topological) level:
\[
  G_{\mathrm{phys}} = \frac{\SU(3) \times \SU(2) \times \U(1)_Y}
  {\mathbb{Z}_6}\,,
\]
where the $\mathbb{Z}_6$ quotient is derived in
Proposition~\ref{prop:global_structure} below, $\SU(n_\alpha)$ is the
universal cover of $\mathrm{PU}(n_\alpha)$, and $\U(1)_Y$ is the
unique relative-phase factor from Proposition~7.1 of \TSI{}.

The derivation chain produces $\mathrm{Inn}(\mathcal{A}) = \prod_\alpha
\mathrm{PU}(n_\alpha)$ as the automorphism group of the algebra, but the
physical gauge group that acts faithfully on matter fields is
$[\SU(3) \times \SU(2) \times \U(1)_Y] / \mathbb{Z}_6$.  The passage between these two descriptions involves
(a)~a lift from $\mathrm{PU}(n)$ to $\SU(n)$ required by the
representation theory of matter, and (b)~a question about global
structure (center quotients) that determines the topology of the gauge
group.  We address both formally.

\begin{proposition}[Lift from $\mathrm{PU}(n)$ to $\SU(n)$]%
\label{prop:PU_to_SU}
The physical gauge group acting on matter fields in $\mathcal{H}_\omega$
is $\SU(n)$, not $\mathrm{PU}(n)$, for each Wedderburn block
$M_n(\mathbb{C})$ with $n \ge 2$.
\end{proposition}

\begin{proof}
\textit{Step 1 (Projective obstruction).}\quad
$\mathrm{Inn}(\mathcal{A}) = \prod_\alpha \mathrm{PU}(n_\alpha)$ acts
on $\mathcal{A}$ by conjugation: $\alpha_U(a) = UaU^*$.  This is the
\emph{adjoint} action.  Matter multiplets, however, live in the
independently specified matter category
(Definition~\ref{def:matter_category}), which requires the
\emph{fundamental} representation of each block; the GNS Hilbert space
$\mathcal{H}_\omega$ is the algebra's carrier (of dimension $\sum_k
n_k^2$, containing $n_k$ copies of each fundamental,
Proposition~\ref{prop:exact_seq}(v)) and serves here only to exhibit
that the fundamental exists as a subrepresentation: $v \mapsto Uv$ for
$v \in \mathbb{C}^n \subset \mathcal{H}_\omega$.  The fundamental is a linear representation of
$\mathrm{U}(n)$ but only a \emph{projective} representation of
$\mathrm{PU}(n)$: the lift $\mathrm{PU}(n) \to \mathrm{GL}(n,
\mathbb{C})$ requires choosing a phase for each element, and there is
no continuous choice that is globally consistent (the obstruction is
the second group cohomology $H^2(\mathrm{PU}(n), \mathrm{U}(1))
\cong \mathbb{Z}_n$).

\textit{Step 2 (Linearization by lift to universal cover).}\quad
The universal covering group of $\mathrm{PU}(n)$ is $\SU(n)$, with
covering map $\SU(n) \twoheadrightarrow \mathrm{PU}(n)$ and kernel
$Z(\SU(n)) \cong \mathbb{Z}_n$ (the center, generated by
$e^{2\pi i / n}\,\mathbf{1}_n$).  Every projective representation of
$\mathrm{PU}(n)$ lifts to a genuine (linear) representation of
$\SU(n)$.  The Lie algebras are isomorphic:
$\mathfrak{su}(n) = \mathfrak{pu}(n)$, so the gauge field content
(connection 1-form $A_\mu$, curvature $F_{\mu\nu}$, covariant
derivative $D_\mu$) is identical regardless of whether one uses
$\SU(n)$ or $\mathrm{PU}(n)$.  The distinction is purely global:
$\SU(n)$ admits the fundamental representation while $\mathrm{PU}(n)$
does not.

\textit{Step 3 (Carrier requirement forces the fundamental).}\quad
The carrier requirement R1 (Theorem~R, \S\ref{sec:carrier_prereqs})
demands a faithful complex carrier of dimension $n$ (from $B1'$).
R1 is a requirement on the matter category
(Definition~\ref{def:matter_category}), and this forcing argument runs
there --- $\mathcal{H}_\omega$ entered only in Step~1, to exhibit that
the fundamental exists.  The fundamental representation of $\SU(n)$ on
$\mathbb{C}^n$ is the unique faithful irreducible representation of
dimension $n$; the adjoint ($\mathrm{PU}(n)$ on $\mathfrak{su}(n)
\cong \mathbb{C}^{n^2-1}$) has dimension $n^2 - 1$ and is not faithful
on $\mathbb{C}^n$.  R1 therefore forces the fundamental, which in turn
forces $\SU(n)$ as the gauge group acting on matter.
\end{proof}

\begin{proposition}[Global structure of the gauge group]%
\label{prop:global_structure}\label{prop:global_Z6}
The gauge group acting on the matter content of the APF is:
\[
  G_{\mathrm{phys}} \;=\;
  \frac{\SU(3) \times \SU(2) \times \U(1)_Y}{\mathbb{Z}_6}\,,
\]
where $\mathbb{Z}_6 \cong Z(\SU(3)) \times Z(\SU(2)) \cap \ker(Y)$
is the subgroup of the center that acts trivially on all matter
representations.
\end{proposition}

\begin{proof}
The group that acts faithfully on matter is not the direct product
$\SU(3) \times \SU(2) \times \U(1)_Y$ but its quotient by the
subgroup of center elements acting trivially on all derived multiplets.

\textit{Step 1 (Center identification).}\quad
The center of $\SU(3) \times \SU(2) \times \U(1)_Y$ is
$\mathbb{Z}_3 \times \mathbb{Z}_2 \times \U(1)_Y$.  A central element
$(e^{2\pi i k/3}\,\mathbf{1}_3,\; (-1)^m\,\mathbf{1}_2,\;
e^{iY\theta})$ acts on a multiplet $(r_3, r_2, Y_f)$ by the phase
$e^{2\pi i k\, c(r_3)/3}\, (-1)^{m\, t(r_2)}\, e^{iY_f\theta}$,
where $c(r_3)$ is the $\mathbb{Z}_3$ triality of the $\SU(3)$
representation and $t(r_2)$ is the $\mathbb{Z}_2$ character of the
$\SU(2)$ representation.

\textit{Step 2 (Trivially acting subgroup --- explicit verification).}\quad
A central element acts trivially on \emph{all} derived multiplets iff
$e^{2\pi i k\, c(r_3)/3}\, (-1)^{m\, t(r_2)}\, e^{iY_f\theta} = 1$
for every multiplet in the SM template.  The derived multiplet list
(from the fermion scan, $T_{\mathrm{field}}$) assigns each multiplet
definite quantum numbers.

\textit{Normalization.}\quad We work with the integer-valued charge
$Y_{\mathrm{int}} := 6Y$ (so that $Y_Q = 1/6$ becomes
$Y_{\mathrm{int}} = 1$, etc.).  The $\U(1)_Y$ element $e^{i\theta}$
acts on a multiplet with charge $Y_{\mathrm{int}}$ by the phase
$e^{iY_{\mathrm{int}}\theta}$.  A center element
$(e^{2\pi i k/3}\,\mathbf{1}_3,\; (-1)^m\,\mathbf{1}_2,\;
e^{i\theta})$ acts on multiplet $(c, t, Y_{\mathrm{int}})$ by:
\begin{equation}\label{eq:center_phase}
  \phi(k,m,\theta) \;=\; e^{2\pi i k\, c/3}\;
  e^{\pi i m\, t}\; e^{i Y_{\mathrm{int}}\,\theta}.
\end{equation}
Setting $\theta = \pi/3$ (the candidate $\mathbb{Z}_6$ generator)
and $(k, m) = (1, 1)$:
\[
  \phi\bigl(1,1,\tfrac{\pi}{3}\bigr) = \exp\Bigl[\,
  i\pi\,\frac{2c + 3t + Y_{\mathrm{int}}}{3}\Bigr].
\]
This equals $1$ iff $2c + 3t + Y_{\mathrm{int}} \equiv 0 \pmod{6}$.
We verify this for every multiplet type in one generation of the
derived fermion template:

\smallskip
\begin{center}\small
\renewcommand{\arraystretch}{1.15}
\begin{tabular}{@{}lccccccr@{}}
\toprule
\textbf{Multiplet} & $\SU(3)$ & $\SU(2)$ & $Y$ &
  $c$ & $t$ & $Y_{\mathrm{int}}$ &
  $2c{+}3t{+}Y_{\mathrm{int}}$ \\
\midrule
$Q_L$ & $\mathbf{3}$ & $\mathbf{2}$ & $1/6$ & 1 & 1 & 1 & $6 \equiv 0$ \\
$u_R$ & $\mathbf{3}$ & $\mathbf{1}$ & $2/3$ & 1 & 0 & 4 & $6 \equiv 0$ \\
$d_R$ & $\mathbf{3}$ & $\mathbf{1}$ & $-1/3$ & 1 & 0 & $-2$ & $0 \equiv 0$ \\
$L_L$ & $\mathbf{1}$ & $\mathbf{2}$ & $-1/2$ & 0 & 1 & $-3$ & $0 \equiv 0$ \\
$e_R$ & $\mathbf{1}$ & $\mathbf{1}$ & $-1$ & 0 & 0 & $-6$ & $-6 \equiv 0$ \\
\midrule
$H$ (Higgs) & $\mathbf{1}$ & $\mathbf{2}$ & $1/2$ & 0 & 1 & 3 & $6 \equiv 0$ \\
\bottomrule
\end{tabular}
\end{center}

\smallskip\noindent All entries satisfy
$2c + 3t + Y_{\mathrm{int}} \equiv 0 \pmod{6}$.  The Higgs doublet
$(1, 2, 1/2)$ is included because it participates in gauge
interactions; its compatibility is necessary for consistency of the
Yukawa sector.

The $\mathbb{Z}_6$ generator is therefore
$g := (\omega_3\,\mathbf{1}_3,\; -\mathbf{1}_2,\;
e^{i\pi/3})$, and the six elements $\{g^0, g^1, \ldots, g^5\}$
exhaust the kernel.  To confirm that the kernel is exactly
$\mathbb{Z}_6$ and not a proper subgroup, note that
$\mathbb{Z}_3 \times \mathbb{Z}_2 \cong \mathbb{Z}_6$ (since
$\gcd(2,3) = 1$), so $\mathbb{Z}_6$ is the largest discrete
subgroup of $Z(\SU(3)) \times Z(\SU(2))$.  It remains to verify that
elements of the center with \emph{trivial} $\U(1)$ phase do not
act trivially: $(\omega_3,\, \mathbf{1}_2,\, 1)$ acts on $Q_L$ as
$\omega_3 \ne 1$, and $(\mathbf{1}_3,\, -\mathbf{1}_2,\, 1)$ acts
on $Q_L$ as $-1 \ne 1$.  The nontrivial $\U(1)$ phase
$e^{i\pi/3}$ in the generator $g$ is essential: the kernel is a
\emph{diagonal} $\mathbb{Z}_6$ inside
$Z(\SU(3)) \times Z(\SU(2)) \times \U(1)_Y$, not the direct
product $\mathbb{Z}_3 \times \mathbb{Z}_2 \times \{1\}$.

\textit{Step 3 (Quotient).}\quad
The faithful gauge group is therefore
$[\SU(3) \times \SU(2) \times \U(1)_Y] / \mathbb{Z}_6$.  This is
the global form of the Standard Model gauge group, consistent with the
analysis of Aharony, Seiberg, and Tachikawa (2013) and
Tong (2017).
\end{proof}

\begin{corollary}[Fermion template compatible with quotient]%
\label{cor:fermion_Z6_compat}
Every multiplet type in the derived fermion template
($T_{\mathrm{field}}$, \S5 of \TSI{}) defines a
well-defined (linear, not merely projective) representation of
$G_{\mathrm{phys}} = [\SU(3) \times \SU(2) \times \U(1)_Y] /
\mathbb{Z}_6$.
\end{corollary}

\begin{proof}
A representation of $\SU(3) \times \SU(2) \times \U(1)_Y$ descends
to a representation of the quotient
$G_{\mathrm{phys}}$ if and only if the $\mathbb{Z}_6$ kernel acts
trivially on that representation.  The table above verifies this
condition ($2c + 3t + Y_{\mathrm{int}} \equiv 0 \pmod{6}$) for each
of the five fermion multiplet types in one generation and for the
Higgs doublet.  Since the three generations carry identical quantum
numbers, the check extends to all 45 Weyl fermions (15 types
$\times$ 3 generations).  No exotic representation outside the
$\mathbb{Z}_6$-compatible sublattice appears in the derived matter
content.

Conversely, the $\mathbb{Z}_6$ compatibility is not an accident:
it follows from the anomaly cancellation conditions
(\S6 of \TSI{}), which relate the $\U(1)_Y$ charges to
the $\SU(3)$ and $\SU(2)$ representations.  Specifically, the
anomaly equation
$z^2 - 2z - (N_c^2 - 1) = 0$ with $N_c = 3$ yields
$z \in \{4, -2\}$; the resulting charge assignments automatically
satisfy the $\mathbb{Z}_6$ condition because the anomaly system
forces $Y_{\mathrm{int}} \equiv -(2c + 3t) \pmod{6}$ for every
multiplet.  This is a derived consistency, not a coincidence.
\end{proof}

\noindent\textit{Why the PSU($n$) issue is not a loophole.}\quad
The reviewer's concern can be sharpened: since $\mathrm{Inn}(\mathcal{A})
= \mathrm{PU}(3) \times \mathrm{PU}(2) \cong \mathrm{PSU}(3) \times
\mathrm{SO}(3)$, and neither $\mathrm{PSU}(3)$ nor $\mathrm{SO}(3)$
admits the fundamental representation, the matter content might appear
incompatible with the group that the algebra naturally produces.  The
resolution is that the algebra and the Hilbert space carry
\emph{different} representations of the same underlying structure:
\begin{enumerate}[label=(\alph*),nosep]
  \item The algebra $\mathcal{A} = M_3(\mathbb{C}) \oplus
  M_2(\mathbb{C}) \oplus \mathbb{C}$ is acted on by
  $\mathrm{Inn}(\mathcal{A})$ via conjugation (adjoint
  representation).  This action is well-defined on
  $\mathrm{PU}(n)$ because the phase ambiguity cancels:
  $e^{i\theta} U a U^* e^{-i\theta} = UaU^*$.
  \item The GNS Hilbert space $\mathcal{H}_\omega$ carries a linear
  action of $\mathrm{U}(\mathcal{A}) = \mathrm{U}(3) \times
  \mathrm{U}(2) \times \mathrm{U}(1)$
  (Proposition~\ref{prop:exact_seq}).  The fundamental representation
  $v \mapsto Uv$ requires the full unitary group, not the quotient.
  $\mathcal{H}_\omega$ is the algebra-side bookkeeping object (of
  dimension $\sum_k n_k^2$, with $n_k$ copies of each fundamental); it
  exhibits the fundamental of each block as a subrepresentation, but
  it is not the matter space.
  \item The physical gauge group $G_{\mathrm{phys}}$ is the image of
  the lifted group in the general linear group of the \emph{matter}
  representation space --- the carrier of the derived multiplet
  content (Definition~\ref{def:matter_category}) --- modulo the center
  acting trivially there; the GNS space plays no role in this
  definition beyond the bookkeeping of item~(b).  The exact sequence
  (Proposition~\ref{prop:exact_seq}) shows that this image is
  $[\SU(3) \times \SU(2) \times \U(1)_Y] / \mathbb{Z}_6$:
  larger than $\mathrm{Inn}(\mathcal{A})$ (because it includes the
  $\U(1)_Y$ grading from the center) but smaller than
  $\mathrm{U}(\mathcal{A})$ (because the diagonal $\mathrm{U}(1)$
  and the $\mathbb{Z}_6$ kernel are removed).
\end{enumerate}
One premise rides this list and is named here (it is consumed
wherever the $\U(1)_Y$ grading is called physical): central unitaries
act trivially by conjugation on $\mathcal{A}$, so the relative phases
of $T_{\mathrm{rel}}$ are physical exactly because the independently
specified charged matter representations
(Definition~\ref{def:matter_category}) see them; absent charged
matter, the grading would be unobservable.
The passage $\mathrm{Inn}(\mathcal{A}) \to G_{\mathrm{phys}}$ is
therefore not a ``choice'' to work with $\SU(n)$ instead of
$\mathrm{PU}(n)$; it is a derived consequence of the map
$\mathrm{U}(\mathcal{A}) \twoheadrightarrow G_{\mathrm{phys}}$
(Propositions~\ref{prop:exact_seq}--\ref{prop:global_structure}),
forced by the requirement that matter fields transform linearly
and that the kernel of the representation is exactly
$\mathbb{Z}_6$.  One convention rides this conclusion and is named
here: the computation identifies the kernel of the \emph{chosen
matter representation} and hence the maximal faithful quotient
\emph{for that representation}.  Calling this quotient ``the'' global
gauge group is the convention of taking the maximal quotient
compatible with the derived matter content --- the standard choice,
adopted explicitly, not an additional theorem.

\smallskip\noindent\textit{Physical significance of the global
structure.}\quad
The distinction between $\SU(3) \times \SU(2) \times \U(1)_Y$ and
$[\SU(3) \times \SU(2) \times \U(1)_Y] / \mathbb{Z}_6$ is invisible
at the level of the Lie algebra (perturbative physics), but has
observable consequences in the non-perturbative sector:
\begin{enumerate}[label=(\alph*),nosep]
  \item \textit{Allowed monopole charges.}\quad The Dirac quantization
  condition $eg = n/2$ depends on the fundamental group
  $\pi_1(G_{\mathrm{phys}})$.  For the quotient group,
  $\pi_1([\SU(3) \times \SU(2) \times \U(1)] / \mathbb{Z}_6) \cong
  \mathbb{Z}$, and the minimal monopole charge is $6$ times the naive
  Dirac quantum.
  \item \textit{Allowed line operators.}\quad In the path-integral
  formulation, the global form determines which Wilson and 't~Hooft
  line operators are allowed.  The $\mathbb{Z}_6$ quotient restricts
  the allowed electric and magnetic charges to a sublattice of the
  weight lattice.
  \item \textit{Topological sectors.}\quad The distinct homotopy
  groups of the quotient group vs.\ the product group affect the
  classification of instantons and the vacuum structure.
\end{enumerate}
The APF predicts the $\mathbb{Z}_6$ quotient because the derivation
proceeds through $\mathrm{Inn}(\mathcal{A})$ --- which is
$\mathrm{PU}(3) \times \mathrm{PU}(2)$, already a quotient --- and
then lifts to $\SU(n)$ only as far as the matter representations
require.  The lift introduces exactly the covering needed for the
fundamental representation (Proposition~\ref{prop:PU_to_SU}), but the
center elements that act trivially on all matter multiplets remain
identified.  The global form is therefore a prediction, not an
assumption: it follows from the algebraic derivation chain
$\mathrm{Inn}(\mathcal{A}) \to \mathrm{PU}(n) \to \SU(n) / \Gamma$
combined with the specific matter content from the fermion scan.

\smallskip\noindent\textit{Reconciliation summary.}\quad
The three levels of description are:
\begin{center}\small
\begin{tabular}{@{}lll@{}}
\toprule
\textbf{Group} & \textbf{Acts on} & \textbf{Representation type} \\
\midrule
$\mathrm{Inn}(\mathcal{A}) = \prod_k \mathrm{PU}(n_k)$ &
  $\mathcal{A}$ (algebra) & Adjoint (conjugation) \\
$\mathrm{U}(\mathcal{A}) = \prod_k \mathrm{U}(n_k)$ &
  $\mathcal{H}_\omega$ (Hilbert space) & Fundamental (linear) \\
$G_{\mathrm{phys}} = [\SU(3){\times}\SU(2){\times}\U(1)_Y]/\mathbb{Z}_6$ &
  Matter multiplets & Fundamental, faithful \\
\bottomrule
\end{tabular}
\end{center}
The passage from the first row to the third involves two steps: lifting
from $\mathrm{PU}(n)$ to $\SU(n)$ (forced by the fundamental
representation, Proposition~\ref{prop:PU_to_SU}), then quotienting by
the trivially acting center $\mathbb{Z}_6$ (forced by faithfulness on
the derived matter content, Proposition~\ref{prop:global_structure}).
The Lie algebras are identical at all three levels; the distinctions are
global (topological).

\smallskip\noindent\textit{Status of $\U(1)_Y$.}\quad
The abelian factor is determined by a two-layer argument:
\begin{enumerate}[nosep,label=\textbf{Layer \arabic*:}]
  \item \textbf{How many abelian directions?}\quad
  Proposition~7.1 of \TSI{} (this section): exactly one $\U(1)$
  from the relative-phase torus $T_{\mathrm{rel}}$ is admissible.
  Enforcement completeness (EC) forces at least one; capacity
  minimality (CM) excludes all but one.

  \item \textbf{Which direction?}\quad
  Theorem~6.1 of \TSI{} (\S6 of \TSI{}): among all
  one-dimensional subgroups of $T_{\mathrm{rel}}$, the unique one
  consistent with all four anomaly-cancellation conditions for the
  surviving fermion template is hypercharge $Y$, with assignments
  $Y_Q = 1/6$, $Y_u = 2/3$, $Y_d = -1/3$, $Y_L = -1/2$, $Y_e = -1$,
  up to $u \leftrightarrow d$ relabeling and overall normalization
  (the ratios are derived; the unit is the compact-$\U(1)$
  charge-lattice convention, front-matter ledger).
\end{enumerate}
Neither layer alone suffices.  Layer~1 without Layer~2 gives ``one
abelian factor exists'' without identifying it.  Layer~2 without
Layer~1 gives ``these hypercharges cancel anomalies'' without
explaining why a second independent $\U(1)$ is excluded.  Together
they uniquely determine the Lie algebra of
$G_{\mathrm{phys}}$ as $\mathfrak{su}(3) \oplus \mathfrak{su}(2)
\oplus \mathfrak{u}(1)_Y$; the global form
$[\SU(3) \times \SU(2) \times \U(1)_Y] / \mathbb{Z}_6$ is determined
by Proposition~\ref{prop:global_structure}.

$\U(1)_Y$ is not an inner automorphism of $\mathcal{A}$; it is a
central element of the unitary group that acts non-trivially on matter
representations.  The non-abelian gauge factors ($\SU(3)$, $\SU(2)$)
are structural (from $\mathrm{Inn}(\mathcal{A})$); the abelian factor
is a grading channel (from $T_{\mathrm{rel}}$,
selected by Proposition~7.1 of \TSI{} and fixed by
Theorem~6.1 of \TSI{}), physical because the charged matter
representations see it (Definition~\ref{def:matter_category}).
The non-abelian factors are derived from A1 via non-closure; the
abelian factor follows from admissibility completeness under its
named reading (R-EC-inv $+$ I-typing,
Theorem~7.7 of \TSI{} H3--H4), minimality, and
anomaly cancellation.  Neither is postulated, and the conditionality
of the abelian leg is stated where it is proved.

\smallskip\noindent\textit{Connection to NCG.}\quad
This mechanism parallels Connes' ``unimodular'' condition in NCG, where
the same $\U(1)_Y$ emerges from the determinant constraint on the
automorphism group of the spectral triple~\cite{Connes1996}.  In the APF,
the selection principle is admissibility completeness
(Proposition~7.1 of \TSI{}) rather than unimodularity, but the
algebraic mechanism is identical: the abelian factor comes from the
\emph{center} of the multi-block unitary group, not from inner
automorphisms.

\noindent\textit{Connection to T3.}\quad T3 (Paper~1 Supplement) derives
the gauge structure from the Doplicher--Roberts reconstruction theorem,
which recovers a compact group from its representation category.  The
gauge identification theorem above provides the complementary result:
it identifies \emph{which} automorphisms constitute the gauge group,
and separates them from kinematic automorphisms.  Together, T3 and
Theorem~\ref{thm:gauge_ident} establish that the admissibility algebra
produces a non-abelian compact Lie gauge group acting on internal
degrees of freedom at each interface.

\smallskip\noindent\textit{Connection to noncommutative geometry.}\quad
The identification $G = \mathrm{Inn}(\mathcal{A})$ parallels a central
theme of Connes' NCG program, where the gauge group of the Standard
Model arises as the inner automorphism group of an almost-commutative
spectral triple~\cite{Connes1996}.  The present framework arrives at
the same structural identification from different premises: the algebra
$\mathcal{A}$ is derived from A1 ($T_{\mathrm{alg}}$), not postulated
as a spectral triple; the inner/outer separation follows from interface
locality ($L_{\mathrm{loc}}$) rather than from the order-one condition;
and the carrier constraints (Theorem~R) replace the KO-dimension and
reality conditions.  A comparison of the constraint structures is in
Appendix~A of \TSI{}.

\smallskip\noindent\textit{State independence.}\quad
The gauge group $G = \mathrm{Inn}(\mathcal{A})$ depends only on the
abstract $*$-algebra $\mathcal{A}$, not on the choice of faithful state
$\omega$ used in the GNS construction.  This is because
$\mathrm{Inn}(\mathcal{A})$ is defined purely algebraically: it is the
group of automorphisms $\alpha_U: a \mapsto UaU^*$ for $U$ unitary in
$\mathcal{A}$.  This definition references only the algebra's
multiplication and involution, not any representation or state.

For the same reason, the Wedderburn block structure
$\mathcal{A} \cong \bigoplus_k M_{n_k}(\mathbb{C})$, the block sizes
$\{n_k\}$, and the product $\mathrm{Inn}(\mathcal{A}) = \prod_k
\mathrm{PU}(n_k)$ are all $\omega$-independent.  For any \emph{faithful}
state the GNS null ideal is zero, so $\mathcal{H}_\omega \cong
\mathcal{A}$ as a vector space and
$\dim \mathcal{H}_\omega = \sum_k n_k^2$ --- the same dimension for
every faithful $\omega$, with left-regular multiplicities ($n_k$ copies
of $\mathbb{C}^{n_k}$).  What varies with $\omega$ is the inner
product's weighting across blocks, not the dimension.  (Earlier
versions claimed different faithful states could produce different GNS
dimensions; that is false in this finite-dimensional setting.)  The
gauge group is the same in every realization.

\smallskip\noindent\textit{Which AQFT algebra does $\mathcal{A}$ correspond to?}\quad
In the Doplicher--Haag--Roberts (DHR) framework of AQFT, two algebras
play distinct roles: the \emph{observable algebra}
$\mathcal{A}_{\mathrm{obs}}$ (gauge-invariant quantities accessible to
local measurements) and the \emph{field algebra} $\mathcal{F}$ (which
includes charged fields that transform non-trivially under the gauge
group).  The gauge group $G$ acts on $\mathcal{F}$ by automorphisms
that fix $\mathcal{A}_{\mathrm{obs}}$ pointwise:
$\mathcal{A}_{\mathrm{obs}} = \mathcal{F}^G$, the fixed-point
subalgebra.

The APF admissibility algebra $\mathcal{A}$ corresponds to the
\emph{field algebra} $\mathcal{F}$, not the observable algebra.
This identification is forced by the construction:
\begin{enumerate}[label=(\alph*),nosep]
  \item $\mathcal{A}$ is generated by all admissibility projections
  $\{E_d\}$ and joint defenders $\{E_{\{d_i,d_j\}}\}$, including those
  that act on individual charged sectors ($M_{d_k}$ in the Wedderburn
  decomposition).  An operator that projects onto a single color-triplet
  subspace is not gauge-invariant: it distinguishes between colors, and
  a gauge rotation maps it to a different projection.  Such operators
  belong to $\mathcal{F}$, not to $\mathcal{A}_{\mathrm{obs}}$.
  \item Gauge transformations are inner automorphisms of $\mathcal{A}$
  (Theorem~\ref{thm:gauge_ident}):
  $\alpha_U(a) = UaU^*$ for $U \in \mathcal{U}(\mathcal{A})$.  This
  is precisely the AQFT statement that gauge transformations are inner
  on the field algebra.  In AQFT, gauge transformations are
  \emph{not} inner automorphisms of $\mathcal{A}_{\mathrm{obs}}$ ---
  they act trivially there (by definition of gauge invariance).  No
  conflict arises because the APF identifies
  $\mathcal{A} = \mathcal{F}$, not
  $\mathcal{A} = \mathcal{A}_{\mathrm{obs}}$.
  \item The observable algebra in the APF is the gauge-invariant
  subalgebra: $\mathcal{A}_{\mathrm{obs}} := \mathcal{A}^G =
  \{a \in \mathcal{A} : \alpha_U(a) = a \;\forall\, U \in G\}$.
  This is derived, not postulated: once $G = \mathrm{Inn}(\mathcal{A})$
  is identified (Theorem~\ref{thm:gauge_ident}),
  $\mathcal{A}^G$ is determined.  The
  Doplicher--Roberts reconstruction (T3, Paper~1 Supplement) then
  recovers $G$ from the superselection structure of
  $\mathcal{A}_{\mathrm{obs}}$, closing the circle.
\end{enumerate}
In summary: the APF's $\mathcal{A}$ is the field algebra
$\mathcal{F}$; gauge $=$ Inn($\mathcal{A}$) is the standard
DHR statement that gauge transformations act innerly on
$\mathcal{F}$; and the observable algebra
$\mathcal{A}_{\mathrm{obs}} = \mathcal{A}^G$ is the
gauge-invariant fixed-point subalgebra, consistent with the
standard AQFT framework.
\coderef{check\_T3}{core.py}


\subsection{Why the gauge group is a compact Lie group}
\label{sec:compact_lie}

\noindent The gauge group $G = \mathrm{Inn}(\mathcal{A})$ is a compact
Lie group by direct construction, not by an appeal to general
topological-group theory:

\smallskip\noindent\textit{Step 1 (Matrix Lie group).}\quad
$\mathcal{A} \cong \bigoplus_k M_{n_k}(\mathbb{C})$ (Wedderburn--Artin,
T2a).  By the gauge identification theorem
(Theorem~\ref{thm:gauge_ident}),
$G = \prod_k \mathrm{PU}(n_k)$.  Each $\mathrm{PU}(n_k) =
\mathrm{U}(n_k) / \mathrm{U}(1)$ is a matrix Lie group: it is a
closed subgroup of $\mathrm{GL}(n_k^2, \mathbb{R})$ (via the adjoint
representation), hence a Lie group by the closed-subgroup theorem
(Cartan, 1930).

\smallskip\noindent\textit{Step 2 (Compactness).}\quad
A1 (finite capacity) bounds $\dim(\mathcal{H}_\omega) < \infty$
($T_{\mathrm{form}}$).  $\mathrm{U}(n)$ is compact (closed and bounded
in $M_n(\mathbb{C})$); $\mathrm{PU}(n)$ is the continuous image of a
compact group, hence compact.  A finite product of compact groups is
compact.  Therefore $G = \prod_k \mathrm{PU}(n_k)$ is compact.

\smallskip\noindent\textit{Consequence.}\quad The gauge group $G$ is a
non-abelian compact Lie group acting faithfully on a finite-dimensional
Hilbert space.  The classification of such groups reduces to the
classification of compact simple Lie algebras (Section~4 of \TSI{}).


\subsection{Finite capacity vs.\ type III local algebras}
\label{sec:finite_vs_typeIII}

Finite dimensionality is not a convenience in this supplement; it is the
mathematical expression of A1 at a finite interface.  It enters the
substrate theorem, Wedderburn--Artin decomposition, compactness of
$\prod_k\mathrm{PU}(n_k)$, finite Lie-algebra classification, Theorem~R,
the fermion scan, and the integer capacity count.  The comparison with
AQFT should therefore be read as a regime comparison, not as an attempted
finite-dimensional truncation of a type~III theory.

In AQFT, local algebras for continuum regions are generically type~III
factors.  APF instead treats such behavior as a continuum-refinement shadow
of finite-capacity interface algebras.  The finite statement is deliberately
stronger and simpler:

\begin{proposition}[Type III incompatibility]\label{prop:typeIII_incompat}
Let $\mathcal A(\Gamma)=M_n(\mathbb C)$ with
$n=\lfloor C(\Gamma)/\varepsilon^*\rfloor<\infty$.  Then
$\mathcal A(\Gamma)$ has a faithful normal trace, contains minimal
projections, and satisfies the split property for disjoint finite
interfaces.  Hence it is type~I$_n$, not type~III.
\end{proposition}

\begin{proof}
All three statements are standard for a finite matrix algebra: the
normalized trace is faithful and tracial; rank-one projections are minimal;
and commuting finite subalgebras have an algebraic tensor product already
complete in norm.  Type~III factors have no faithful normal semifinite
trace and no nonzero finite projections, so the regimes are incompatible at
finite $n$.
\end{proof}

\begin{proposition}[Controlled limit to type III]\label{prop:typeIII_limit}
For a refinement sequence $\mathcal A_n=M_n(\mathbb C)$ equipped with
compatible KMS states, the inductive/GNS limit can generate a
Powers--Araki--Woods type~III$_\lambda$ factor; in the saturation limit
$\lambda\to1$ this approaches type~III$_1$.  The finite trace then fails to
extend to a normal faithful trace on the limit factor.
\end{proposition}

\begin{proof}
This is the standard UHF/CAR-to-Powers-factor mechanism: finite matrix
algebras embed into an inductive system, the product KMS state defines the
GNS representation, and the modular spectrum fixes the Connes parameter
$\lambda=e^{-\beta\varepsilon^*}$.  Taking the infinite-refinement limit
before asking local-algebra questions destroys the finite trace.
\end{proof}


\begin{proposition}[Continuum-limit audit plan]
\label{prop:continuum_plan}
A finite APF net is compatible with type-III AQFT only under an explicit
three-part limiting plan: (i) nested finite subalgebras embed by
inclusion/refinement maps; (ii) the limiting state satisfies the usual
regularity hypotheses used in split/nuclearity arguments; and (iii) the
sector labels locked at finite capacity persist rather than opening new
independent low-cost sectors.  Under these hypotheses, strict tensor
factorization is a finite-regime tool, while the continuum limit may have
the standard type-III obstruction to traces and exact local subsystems.
\end{proposition}

\begin{proof}
The first condition supplies the inductive net.  The second prevents a
pathological limit in which the GNS closure loses the local regularity
needed for AQFT comparison.  The third is the APF-specific stability
condition: if refinement introduced a new independent low-cost sector, the
finite gauge/matter classification would have been performed on an
incomplete carrier.  When all three hold, the type-III limit affects the
idealized local-algebra type but not the already locked finite block labels
used by the gauge and matter derivations.
\end{proof}

\noindent\textit{Robustness audit under the limit.}\quad
The continuum passage is not allowed to re-open the finite classification.
The structures below are robust because they are locked at finite capacity
before the limiting operation; the non-robust entries are explicitly bridge
or regulator questions:
\begin{center}
\scriptsize
\renewcommand{\arraystretch}{1.08}
\begin{tabular}{@{}p{3.3cm}p{5.0cm}p{4.0cm}@{}}
\toprule
\textbf{Structure} & \textbf{Status under type-III limit} & \textbf{Reason} \\
\midrule
Block labels $M_n(\mathbb C)$ and $\prod\mathrm{PU}(n)$ & Robust, provided sector-stability hypothesis (iii) holds. & They are finite-capacity superselection data before refinement. \\
Matter multiplet template and hypercharge ray & Robust under stable embedding. & The anomaly equations and center kernel are algebraic on the finite modules. \\
Strict tensor factorization & Not robust. & Replaced by split/nuclearity approximations and Definition~\ref{def:finite_capacity_interface_independence}. \\
Trace on local algebra & Not robust. & The finite trace need not extend to the type-III GNS closure. \\
Regulator-level beta coefficients & Limit-dependent bridge. & Paper~4A must supply the RG regulator/renormalization step. \\
\bottomrule
\end{tabular}
\end{center}
\subsubsection{Gravitational dressing and the scope of locality}
\label{sec:dressing_scope}

The finite locality lemma is pre-gravitational and gauge-fixed at the level
of interface algebras.  It says: if two finite substrates are
$B$-orthogonal and disjoint, their admissibility operations commute.  This is
compatible with gravitational dressing results, where diffeomorphism-
invariant dressed observables necessarily carry long-range metric data and
therefore fail to factorize strictly.  In APF language, the dressing field
is part of the shared interface substrate; once it is included, the
interfaces are not independent in the sense required by
Lemma~\ref{lem:loc_commut}.  Thus gravitational nonlocality is not a
counterexample to the finite lemma; it identifies the regime where the
lemma's disjoint-support hypothesis is false.

\begin{tcolorbox}[apfbox, title={Scope of the finite locality lemma}]
\footnotesize
Lemma~\ref{lem:loc_commut} is a finite-interface theorem with an explicit
hypothesis: disjoint $B$-orthogonal substrates.  It is not a theorem about
fully diffeomorphism-invariant gravitational observables.  Once a dressing
tail is required to define an observable, the tail is part of the substrate
and supplies shared support.  The correct APF replacement for exact
subsystem independence in that regime is finite-capacity interface
independence, not strict tensor factorization.
\end{tcolorbox}

\begin{definition}[Finite-capacity interface independence]
\label{def:finite_capacity_interface_independence}
Two finite interface algebras $\mathcal A(\Gamma_1)$ and
$\mathcal A(\Gamma_2)$ are finite-capacity independent at tolerance
$\eta$ when every joint continuation cost decomposes as
\[
  C_{12}(x,y)=C_1(x)+C_2(y)+I_{12}(x,y),\qquad
  |I_{12}(x,y)|\le \eta
\]
for all tested records $x,y$, and when the interface term $I_{12}$ is
accounted for as shared substrate capacity rather than as a tensor-product
factor.  Exact tensor factorization is the special case $\eta=0$.  A
gravitational dressing tail gives $\eta>0$: it is not a violation of the
finite locality lemma but evidence that the disjoint-support hypothesis has
failed.
\end{definition}

\begin{proposition}[A1-bounded Reeh--Schlieder mechanism]\label{prop:RS_modification}
At finite capacity, cyclicity/separating behavior is limited by the number
of admissible channels.  Reeh--Schlieder-type entanglement can be
approximated on accessible subalgebras, but exact continuum cyclicity for
arbitrarily fine local excitations requires the $C/\varepsilon^*\to\infty$
limit.
\end{proposition}

\begin{proof}
A finite algebra has finite rank and a finite-dimensional state space, so a
local algebra cannot generate an infinite independent set of excitations
from the vacuum vector.  Refinement increases the accessible rank; the
usual AQFT behavior is recovered only after the unbounded limit.
\end{proof}

\subsubsection{Modular theory and the Unruh effect as a limit}
\label{sec:modular_bridge}

Finite matrix algebras have modular automorphism groups for faithful
states, but the full continuum Unruh/KMS structure belongs to the limiting
factor.  The APF reading is therefore: finite interfaces carry regulated
modular flow; the familiar thermal wedge behavior is the large-refinement
limit of that flow.

\subsubsection{Haag duality and locality}
\label{sec:haag_duality}

For disjoint finite interfaces, Haag-duality-like factorization is trivial:
commuting finite algebras split as tensor factors once the substrate
orthogonality assumptions hold.  In the continuum, Haag duality and the
split property require the usual AQFT nuclearity/regularity hypotheses.
APF claims only that the finite theorem is the regulated precursor.

\subsubsection{Independent evidence for finite local capacity}
\label{sec:independent_evidence}

The physical motivation for A1 is the convergence of entropy bounds rather
than a preference for finite-dimensional algebra.  Bekenstein--Hawking
entropy, Bousso's covariant bound, holography, and black-hole information
unitarity all point toward finitely many distinguishable states associated
with a finite-area interface.  APF takes this as fundamental and treats
continuum type~III behavior as idealized emergent behavior.

\begin{theorem}[Stability of structural predictions under
$n \to \infty$]\label{thm:structural_stability}
The finite structural predictions used in Paper~2---the block dimensions,
compact gauge factors, anomaly equations, and template-scan survivor---are
stable under refinement provided the finite stage has already locked the
same separated sectors and no new independent low-cost sector opens in the
limit.
\end{theorem}

\begin{proof}
Refinement adds resolution within already separated sectors.  It does not
change the finite block labels, anomaly equations, or template admissibility
unless a new sector with independent admissibility support appears.  The
latter case is exactly excluded by the persistent-pooling/no-new-sector
hypothesis used in the refinement statement below.
\end{proof}

\begin{definition}[Refinement Hypothesis (RH)]\label{def:RH}
Interfaces admit arbitrarily fine operational refinements whose finite
subalgebras embed compatibly into later refinements.
\end{definition}

\begin{definition}[Capacity Density (CD)]\label{def:CD}
The number of available channels scales with interface area according to a
finite density bound, so the channel count diverges only in the continuum
limit.
\end{definition}

\begin{definition}[Persistent Pooling (PP)]\label{def:PP}
The shared-pool/nonseparability structure persists under refinement; no
refinement splits the locked interface into independent cost-free sectors.
\end{definition}

\begin{proposition}[Effective type-III behavior in the refinement
limit]\label{prop:effective_typeIII}
Under RH, CD, and PP, a sequence of finite APF interface algebras can have
continuum observables whose GNS closure is effectively type~III while the
finite-stage gauge and matter predictions remain unchanged.
\end{proposition}

\begin{proof}
RH supplies the inductive system, CD supplies unbounded channel growth in
the limit, and PP preserves the finite locked sector structure.  The limit
therefore has the usual type~III obstructions to a faithful trace, while
the finite-stage algebraic invariants used for the Standard-Model template
are stable by Theorem~\ref{thm:structural_stability}.
\end{proof}

\begin{tcolorbox}[apfbox, title={Regime: finite-capacity admissibility algebras}]
Paper~2 derives the gauge/matter structure in finite APF interface
algebras.  AQFT type~III locality is treated as a controlled continuum
limit, not as an input.  Gravitational dressing is accommodated by denying
the disjoint-support hypothesis whenever the dressing field supplies a
shared long-range substrate.
\end{tcolorbox}

\newpage


% ══════════════════════════════════════════════════════════════
% ══════════════════════════════════════════════════════════════
\section{Carrier prerequisites: metric, sector-uniform irreversibility, and complex type}
\label{sec:carrier_prereqs}
% ══════════════════════════════════════════════════════════════

Theorem~R (Paper~2, \S3.5) derives three carrier requirements (R1--R3)
from admissibility.  The derivation chains for R1 and R2 depend on three
lemmas originally developed in companion papers.  This section provides
self-contained proofs of all three, eliminating external dependencies.

The dependency order is: $T_{7B} \to L_{\mathrm{irr,uniform}} \to B1'$.
$T_{7B}$ is used by $L_{\mathrm{irr,uniform}}$ (Step~2);
$L_{\mathrm{irr,uniform}}$ closes the R2 chain (chiral carrier);
$B1'$ closes the R1 chain (complex ternary carrier).
$B1'$ is logically independent of $L_{\mathrm{irr,uniform}}$;
both appear in this section because both close formerly external
dependencies.

All upstream requirements (A1, K2, K3, SP, $L_{\mathrm{nc}}$,
$L_{\mathrm{loc}}$, $L_{\mathrm{irr}}$, $T_{\mathrm{sep}}$,
$T_{\mathrm{GNS}}$, T3) are proved in the Paper~1 Supplement.
No result in this section imports from Papers~3--7.

% ======================================================================
% Paper 2 Technical Supplement — §3
% Carrier Prerequisites: Metric, Sector-Uniform Irreversibility,
% and Complex Type
% ======================================================================
%
% PURPOSE.  Theorem~R (Paper~2, §3.5) requires three carrier properties
% (R1–R3) whose proofs depend on three lemmas originally developed in
% companion papers.  This section provides self-contained proofs of all
% three, eliminating external dependencies and making the Theorem~R
% chain fully checkable within Paper~2 and its supplement.
%
% DEPENDENCY ORDER.  T7B → L_irr_uniform → B1'.
% T7B is used by L_irr_uniform (Step 2).
% L_irr_uniform is used by R2 (chiral carrier derivation).
% B1' is used by R1 (complex ternary carrier derivation).
% B1' is logically independent of L_irr_uniform; both are placed
% in this section because both close formerly external dependencies.
%
% UPSTREAM REQUIREMENTS (all proved in Paper 1 / Paper 1 Supplement):
%   A1, K2, K3, SP, L_nc, L_loc, L_irr, T_sep, T_GNS, T3.
% No result in this section imports from Papers 3–7.
% ======================================================================


\subsection{Cost kernel from shared-interface non-factorization ($T_{7B}$)}
\label{sec:T7B}

At separate interfaces, admissibility costs are additive by $L_{\mathrm{loc}}$
(Paper~1 Supplement~\cite{BrookeSupplement}, \S4).  At a shared interface~$\Gamma$, co-located
distinctions produce a superadditive surplus $\Delta > 0$ ($L_\Delta$, Paper~1 Supplement~\cite{BrookeSupplement}).  The question is: what mathematical object
tracks this non-factorization when admissibility operates across
\emph{multiple} directions at a single interface?

The answer is already implicit in the Paper~1 Supplement's derivation
chain: the sector-orthogonal bilinear form $B$ (constructed by $L_\omega$
from K3) and the characterization of admissibility costs as $B$-norms
(from $T_{\mathrm{adj}}$) together force the external feasibility cost
to be an \emph{exact} quadratic form.  $T_{7B}$ makes this explicit.

\begin{theorem}[$T_{7B}$: Positive-Definite Cost Kernel from
Shared-Interface Non-Factorization]\label{thm:T7B}
Let $\Gamma$ be an interface with substrate $S_\Gamma =
\bigoplus_d M_d \oplus \Pi$ ($T_{\mathrm{sep}}$) and sector-orthogonal
bilinear form $B$ ($L_\omega$, Paper~1 Supplement~\cite{BrookeSupplement}).
Let $\varepsilon_{\mathrm{ext}}: S_\Gamma \to \mathbb{R}_{\ge 0}$ denote
the external feasibility cost.  Then:
\begin{enumerate}[label=(\roman*),nosep]
  \item $\varepsilon_{\mathrm{ext}}(v) = B(v,v)$ for all $v \in S_\Gamma$:
        the external feasibility cost is an exact quadratic form.
  \item The associated bilinear form $g := B$ is positive-definite.
  \item $g$ is unique up to one positive scalar \emph{per sector} of the
        $T_{\mathrm{sep}}$ decomposition (the within-sector gauge freedom
        of $B$); it is unique up to a single overall scalar exactly when
        an inter-sector normalization is imposed, which this theorem
        does not supply.  (v3.9 stated ``an overall scalar'' while the
        proof delivers the per-sector freedom; the statement now matches
        the proof.)
\end{enumerate}
No additional regularity assumption is needed.
\end{theorem}

\begin{proof}
\textit{Step 1: Defense costs are $B$-norms.}\quad
By $T_{\mathrm{adj}}$ (Paper~1 Supplement~\cite{BrookeSupplement}, \S5.7),
the admissibility projection for distinction $d$ is the $B$-orthogonal
projection $E_d = \pi_{M_d}^{(B)}$ onto its anchor set $M_d$.  The cost
of defending $d$ against a perturbation $v$ is the $B$-norm of the
component of $v$ in $M_d$:
\[
  \kappa_d(v) = B(E_d\, v,\, E_d\, v) = \|E_d\, v\|_B^2.
\]
This is the minimum-cost defense: $T_{\mathrm{adj}}$ proves that the
$B$-orthogonal projection minimizes the defense cost among all
idempotents with image $M_d$.

\textit{Step 2: K3 gives exact additivity over disjoint supports.}\quad
By the sector decomposition $S_\Gamma = \bigoplus_d M_d \oplus \Pi$ and
K3 (additivity on disjoint supports), the total defense cost decomposes
as:
\[
  \varepsilon_{\mathrm{ext}}(v)
  = \sum_{d} \|E_d\, v\|_B^2 + \|\pi_\Pi\, v\|_B^2
  = \sum_{d} B(v_d, v_d) + B(v_\Pi, v_\Pi),
\]
where $v = \sum_d v_d + v_\Pi$ is the sector decomposition of $v$,
and $\pi_\Pi$ is the $B$-orthogonal projection onto $\Pi$.

\textit{Step 3: Sector orthogonality collapses the sum.}\quad
By $L_\omega$ ($B(v_d, v_{d'}) = 0$ for $d \ne d'$, and
$B(v_d, v_\Pi) = 0$ for all $d$):
\[
  \varepsilon_{\mathrm{ext}}(v)
  = \sum_d B(v_d, v_d) + B(v_\Pi, v_\Pi)
  = B\!\Bigl(\sum_d v_d + v_\Pi,\; \sum_d v_d + v_\Pi\Bigr)
  = B(v, v).
\]
The cross-terms vanish by sector orthogonality.  Therefore
$\varepsilon_{\mathrm{ext}}(v) = B(v,v)$ \emph{exactly}: the quadratic
form is not a leading-order approximation but an identity.  This
gives~(i).

\textit{Step 4: Positive-definiteness.}\quad
For $v \ne 0$: $B(v,v) > 0$ because $B$ is positive-definite
($L_\omega$, Paper~1 Supplement).  Therefore $g = B$ is
positive-definite, giving~(ii).

\textit{Step 5: Uniqueness.}\quad
$L_\omega$ constructs $B$ with forced inter-sector orthogonality and
free within-sector gauge ($B_d \mapsto \alpha_d B_d$, any $\alpha_d > 0$).
The choice of capacity units (A1) fixes an overall scale; the remaining
gauge is the within-sector normalization convention, which does not
affect the geometric structure.  This gives~(iii).
\end{proof}

\smallskip\noindent\textit{Regularity conditions used:}
A1 (finite capacity), K2 (positive perturbation cost), K3 (disjoint-support
additivity), SP (no null directions), $T_{\mathrm{sep}}$ (sector
decomposition), $L_\omega$ (sector-orthogonal bilinear form),
$T_{\mathrm{adj}}$ (admissibility projections are $B$-orthogonal).

\smallskip\noindent\textit{Status:} [P].  No additional regularity
assumption (R7B, K3-bilinear, or differentiability) is needed.  The
quadraticity is \emph{derived} from $L_\omega + T_{\mathrm{adj}}$, not
assumed.

\begin{lemma}[No metric assumed upstream of $T_{7B}$]\label{lem:T7B_provenance}
The bilinear form $B$ used in $T_{7B}$ is constructed from A1, K3, and
FD3 alone.  No norm, inner product, quadratic cost law, or metric
structure is assumed at any point in the construction chain.
\end{lemma}

\begin{proof}
The chain (each step uses only preceding outputs):
\begin{enumerate}[nosep,label=(\arabic*)]
  \item A1 provides real-valued costs $\varepsilon(d) > 0$.
  \item $T_{\mathrm{embed}}$ maps the convex state space into
  $\mathbb{R}^{N+1}$ using cost coordinates (convexity from OR1; no
  inner product invoked).
  \item $T_{\mathrm{sep}}$ decomposes $S_\Gamma = \bigoplus_d M_d
  \oplus \Pi$ from K3 + FD3.  This is a \emph{pre-metric algebraic
  direct sum}: no inner product is used (Appendix~\ref{app:Tsep}).
  \item $L_\omega$ constructs $B$ as a direct sum of free within-sector
  forms with \emph{forced} zero cross-terms.  The vanishing
  $B(u_d, v_{d'}) = 0$ for $d \ne d'$ is a \emph{theorem} of K3
  (Appendix~\ref{app:Lomega}), not a definition.  The within-sector
  forms $B_d$ are free (normalization gauge).
  \item $T_{\mathrm{adj}}$ identifies admissibility maps as $B$-orthogonal
  projections (Appendix~\ref{app:Tadj}).  Uses only $L_\omega$'s
  inter-sector orthogonality.
\end{enumerate}
$T_{7B}$ then shows $\varepsilon_{\mathrm{ext}}(v) = B(v,v)$ by
combining (4) and (5) with K3 (Step 2 of $T_{7B}$'s proof).  The
quadratic form is the \emph{output} of this chain, not an input at
any stage.
\end{proof}

\smallskip\noindent\textit{Why the cost kernel is positive-definite,
not Lorentzian.}\quad
The cost kernel $g = B$ measures admissibility cost at a single
interface~$\Gamma$.  Admissibility cost is positive for every nonzero
perturbation (K2 + SP), so $g$ is positive-definite.  Lorentzian
signature ($-,+,+,+$) arises at a different level: it is the signature
of the \emph{inter-interface} causal structure, where the timelike
direction is determined by $L_{\mathrm{irr}}$ (irreversibility provides
the arrow of time) and the spatial directions are determined by
$L_{\mathrm{loc}}$ (independent interfaces provide spatial separation).
The positive-definite $g$ at each interface is the internal
(Riemannian) metric; the Lorentzian spacetime metric is constructed in
the exited spacetime bridge (Paper~6; boundary in \S3 of \TSI{})
by combining $g$ with the causal ordering of interfaces.  The
disambiguation of these two bilinear structures is detailed in
Remark~\ref{rem:two_bilinear}.  This
is the same structural split as in standard GR, where the spatial metric
is positive-definite and the Lorentzian signature comes from including
the timelike direction.

\smallskip\noindent\textit{Roadmap.}\quad
The full derivation of Lorentzian spacetime from the cost kernel $g$ is
carried by Paper~6 (boundary statement in \S3 of \TSI{}).  The chain runs:
$L_{\mathrm{irr}} \to$ causal partial order $\to$ smooth manifold
(Kolmogorov) $\to$ Lorentzian conformal class (Hawking--King--McCarthy)
$\to$ $d = 4$ (Lovelock uniqueness + DOF minimality) $\to$ Einstein equations.
\coderef{check\_T7B}{gravity.py}


% ------------------------------------------------------------------
\subsection{Sector-uniform irreversibility ($L_{\mathrm{irr,uniform}}$)}
\label{sec:Lirr_uniform}

$L_{\mathrm{irr}}$ (Paper~1 Supplement~\cite{BrookeSupplement}, \S5) derives
irreversibility from A1 + $L_{\mathrm{nc}}$ + $L_{\mathrm{loc}}$: cross-interface
correlations commit capacity that no local observer can recover.
That proof establishes irreversibility \emph{generically} --- it
guarantees that irreversible channels exist at interfaces where
cross-interface correlations are superadditive.  The question for
Theorem~R's R2 requirement is whether this irreversibility extends to
the \emph{gauge sector specifically}: does the sector of internal
(non-gravitational) admissibility contain its own irreversible channels,
or could it remain fully reversible while inheriting apparent
irreversibility from gravity alone?

\begin{theorem}[$L_{\mathrm{irr,uniform}}$: Gauge-Sector Irreversibility at Saturated Interfaces]\label{thm:Lirr_uniform}
Let $\Gamma$ be a shared interface with non-trivial pool ($\Pi \ne \{0\}$),
let $G$ denote the set of gauge-sector distinctions at $\Gamma$, and let
$R_\Gamma$ denote the gravitational record distinctions at $\Gamma$.
Suppose:
\begin{enumerate}[label=(\alph*),nosep]
  \item The gauge sector is non-trivially coupled to gravity at $\Gamma$
  (i.e., the cost kernel $g$ from $T_{7B}$ has nonzero cross-terms
  between $G$ and $R_\Gamma$).
  \item The interface saturation exceeds the bound
  \begin{equation}
    s \;:=\; \frac{\varepsilon(G) + \varepsilon(R_\Gamma)}{C_\Gamma}
    \;>\; 1 \;-\; \frac{\Delta(G, R_\Gamma)}{C_\Gamma - \varepsilon(R_\Gamma)},
    \label{eq:saturation_bound}
  \end{equation}
  where $\Delta(G, R_\Gamma) \ge \varepsilon^* > 0$ is the superadditive
  surplus from $L_\Delta$.
  \item A gauge-sector transition $G_1 \to G_2$ commits positive shared
  surplus with the record sector: $\Delta(G_2, R_\Gamma) > 0$.
\end{enumerate}
Then the gauge sector contains an irreversible channel at $\Gamma$:
there exist gauge-sector histories $H, H'$ differing only in gauge
distinctions such that the reversal $H' \to H$ is inadmissible.
\end{theorem}

\begin{proof}[Proof (three steps)]

\textit{Step 1: Records are interface-local.}\quad
By $L_{\mathrm{loc}}$ (Paper~1 Supplement~\cite{BrookeSupplement}, \S4), admissibility
is distributed over independent interfaces, each with its own budget.
Irreversibility arises because cross-interface correlations commit
capacity that no \emph{single-interface} observer can recover
($L_{\mathrm{irr}}$, Step~3).  In particular, irreversible records at
$\Gamma$ are distinction sets \emph{at} $\Gamma$, not global objects.

\textit{Step 2: Gauge--gravity coupling at shared interfaces.}\quad
By hypothesis~(a) and $T_{7B}$ (Theorem~\ref{thm:T7B}), gauge
distinctions $G$ contribute to the cross-terms of the cost kernel $g$
at $\Gamma$.  Therefore gauge and gravitational admissibility share
interface $\Gamma$.  This sharing implies: there exist admissible
histories $H, H'$ that differ by gauge-sector distinctions and yield
different gravitational records $R_\Gamma(H) \ne R_\Gamma(H')$.  If
no such histories existed, gauge distinctions would have no recordable
consequences at $\Gamma$, contradicting hypothesis~(a).

\textit{Step 3: Budget overrun from non-closure.}\quad
Since $G$ and $R_\Gamma$ coexist at $\Gamma$ with $\Pi \ne \{0\}$,
$L_{\mathrm{nc}}$ gives:
\[
  \varepsilon_\Gamma(G \cup R_\Gamma)
  \;>\;
  \varepsilon_\Gamma(G) + \varepsilon_\Gamma(R_\Gamma).
\]
Consider the reversal of a gauge transition $G_1 \to G_2$.  The
forward transition commits cross-interface surplus
$\Delta(G_2, R_\Gamma) > 0$ (hypothesis~(c)).  The reversal
$G_2 \to G_1$ while $R_\Gamma$ persists requires recovering this
surplus.  By $L_{\mathrm{irr}}$, the surplus includes cross-interface
terms that no single-interface operation can free.  The reversal cost
satisfies:
\[
  \varepsilon_\Gamma(\text{reverse } G_2 \to G_1 \mid R_\Gamma)
  \;\ge\;
  \varepsilon_\Gamma(G_1) + \Delta(G_1, R_\Gamma).
\]
The remaining budget after the forward transition is
$C_\Gamma - \varepsilon(R_\Gamma) - \varepsilon(G_2)$.  By
hypothesis~(b) (saturation bound~(\ref{eq:saturation_bound})),
this remaining budget is strictly less than the reversal cost.
The reversal is inadmissible.
\end{proof}

\begin{corollary}[Near-capacity interfaces]
At any interface within one admissibility channel of full capacity
($s > 1 - \varepsilon^*/C_\Gamma$), condition~(b) is automatically
satisfied because $\Delta \ge \varepsilon^* > 0$ ($L_\Delta$ +
$L_{\varepsilon^*}$).  Conditions~(a) and~(c) hold at any interface
where the gauge sector is non-trivially coupled and undergoes
transitions.  Therefore gauge-sector irreversibility is generic at
near-capacity shared interfaces.
\end{corollary}

\noindent\textit{Regularity conditions used:}
$L_{\mathrm{loc}}$ (locality), $L_{\mathrm{nc}}$ (non-closure),
$L_{\mathrm{irr}}$ (irreversibility from cross-interface correlations),
$T_{7B}$ (cost kernel from shared-interface non-factorization),
A1 (finite budget), SP (non-trivial gauge coupling).

\smallskip\noindent\textit{Scope of the sufficient condition.}\quad
$L_{\mathrm{irr,uniform}}$ is a sufficient-condition theorem: it
guarantees gauge-sector irreversibility at interfaces satisfying
hypotheses~(a)--(c), not at all interfaces.  This is deliberate.
The downstream use (R2: chiral carrier requirement) needs only the
\emph{existence} of irreversible gauge-sector channels at physically
relevant interfaces --- not their universality.  The corollary shows
that the hypotheses are automatically satisfied at near-capacity
shared interfaces, which is the regime where the gauge architecture
is determined.  A universal irreversibility theorem (holding at
arbitrarily low saturation) would be false: at very low saturation
($s \ll 1$), the budget is loose enough that all transitions are
reversible, and the interface is effectively classical.

\smallskip\noindent\textit{Key structural point.}\quad The argument does
\emph{not} assume that the gauge sector has any specific algebraic
structure beyond being a non-trivially coupled set of distinctions at a
shared interface.  The irreversibility is inherited from the \emph{interface
geometry} ($T_{7B}$) and the \emph{non-closure mechanism} ($L_\Delta$),
not from any property specific to non-abelian gauge theories.  Chirality
(R2) then follows because a \emph{vector-like} gauge sector, despite
sharing the interface, produces zero \emph{intrinsic} gauge irreversibility
(all gauge-level processes are CPT-reversible; see Paper~2, \S3.4).  The existence of gauge-sector irreversible
channels (this theorem) combined with the impossibility of intrinsic
irreversibility in a vector-like theory (Paper~2, §3.5) forces chirality.
\coderef{check\_L\_irr\_uniform}{core.py}


% ------------------------------------------------------------------
\subsection{Complex type from oriented composite robustness ($B1'$)}
\label{sec:B1prime}

The R1 carrier requirement (Theorem~R) demands a complex carrier:
$V \not\cong V^*$ as $G$-modules.  The physical reason is that composites
must be \emph{oriented} --- robustly distinguishable from their
anti-composites --- to provide independent admissibility channels under $T_M$.
This subsection proves that real and pseudoreal carriers cannot satisfy
this requirement: the equivariant antilinear map $J$ collapses composite
orientation at zero capacity cost.

\begin{definition}[Oriented composite]\label{def:oriented_composite}
Let $V$ be a faithful irreducible carrier for a gauge group $G$, and
let $B \in (V^{\otimes k})^G$ be a gauge-singlet composite
(a $G$-invariant tensor in $V^{\otimes k}$).  The conjugate composite
is $\bar{B} \in (V^{*\otimes k})^G$.  The composite $B$ is
\emph{oriented} if $B$ and $\bar{B}$ are robustly distinguishable:
no admissibility-preserving relabeling internal to the carrier maps
$B \mapsto \bar{B}$.
\end{definition}

\begin{theorem}[$B1'$: Complex Carrier from Oriented Composite Robustness]\label{thm:B1prime}
Let $V$ be a faithful irreducible carrier for a compact group~$G$.
Suppose:
\begin{enumerate}[label=(B\arabic*),nosep]
  \item $L_{\mathrm{irr}}$-stable gauge-singlet composites $B$ and
        $\bar{B}$ exist whose oriented distinction is robust under
        admissible refinements.
  \item This distinction is not enforced by an independent external
        grading channel (minimality clause from A1).
\end{enumerate}
Then $V$ must be of \emph{complex type}: $V \not\cong V^*$ as $G$-modules.
\end{theorem}

\begin{proof}[Proof (by contradiction)]

Suppose $V \cong V^*$ as $G$-modules.  Then $V$ is either \emph{real type}
(admits a $G$-equivariant $\mathbb{R}$-linear isomorphism $\phi: V \to V^*$
with $\phi = \bar\phi$) or \emph{pseudoreal type} (admits a $G$-equivariant
antilinear map $J: V \to V$ with $J^2 = -\mathrm{id}$, but no real structure).
In either case, there exists a $G$-equivariant antilinear map
$J: V \to V$.

\textit{Step 1: Tensorial extension.}\quad Extend $J$ to $V^{\otimes k}$ by
\[
  J^{\otimes k}(v_1 \otimes v_2 \otimes \cdots \otimes v_k)
  \;=\; J(v_1) \otimes J(v_2) \otimes \cdots \otimes J(v_k).
\]
Since $J$ is $G$-equivariant ($J \circ g = g \circ J$ for all $g \in G$),
$J^{\otimes k}$ is $G$-equivariant on $V^{\otimes k}$.  In particular,
$J^{\otimes k}$ maps gauge-invariant subspaces to gauge-invariant subspaces:
if $B \in (V^{\otimes k})^G$, then $J^{\otimes k}(B) \in (V^{\otimes k})^G$.

\textit{Step 2: Identification of composite with conjugate.}\quad
We split cases.

\smallskip\noindent\textbf{Real-type case.}\quad If $V$ is real-type,
$V \cong V^*$ via a $G$-equivariant real structure $\phi$ with
$\phi = \bar\phi$.  The associated $J: V \to V$ satisfies $J^2 = +\mathrm{id}$
(a genuine involution), so $J^{\otimes k}$ is an involution on $V^{\otimes k}$
for every $k$ and provides a direct $G$-equivariant identification
$(V^{\otimes k})^G \leftrightarrow (V^{*\otimes k})^G$.  The map
$B \mapsto \bar B$ is well-defined as a square-trivial involution, and
the identification is orientation-collapsing in the sense required below.

\smallskip\noindent\textbf{Pseudoreal-type case.}\quad If $V$ is pseudoreal-type,
then by the standard classification (Humphreys 1972; Fulton--Harris 1991,
\S IV.23) $V \cong V^*$ as $G$-modules \emph{by definition of pseudoreal
type}.  The carrier-complexity hypothesis $V \not\cong V^*$ of the theorem
statement is therefore already false for pseudoreal $V$, and the proof
does not need to invoke any $J^{\otimes k}$ argument in this case.
The $J$-extension machinery used below is only necessary in the real-type
case.  The earlier formulation, which attempted to handle both types
under a single $J^{\otimes k}$ identification and raised a potential
signed-tensor subtlety for odd $k$, is therefore unnecessary: pseudoreal
carriers are ruled out by the structural hypothesis $V \not\cong V^*$
directly.

\smallskip\noindent Accordingly, the remainder of the proof treats only
the real-type case, where $J^{\otimes k}$ is a genuine involution and
the identification $(V^{\otimes k})^G \leftrightarrow (V^{*\otimes k})^G$
is unambiguous.

Concretely: let $B = \sum_i c_i\, v_{i_1} \otimes \cdots \otimes v_{i_k}$
be a gauge singlet.  Then
\[
  J^{\otimes k}(B)
  \;=\; \sum_i \bar{c}_i\, J(v_{i_1}) \otimes \cdots \otimes J(v_{i_k})
\]
is also a gauge singlet (by equivariance), and it is the image of $B$ under
conjugation composed with the equivariant isomorphism $V \cong V^*$.
Therefore $J^{\otimes k}(B)$ is identified with $\bar{B}$.

\textit{Step 3: The identification is admissibility-preserving.}\quad
The map $J$ is internal to the carrier --- it is a relabeling of the carrier
basis vectors, not a change of physical state.  It costs zero admissibility
capacity: the admissibility structure at the interface is unchanged because
$J$ commutes with all $G$-actions (equivariance), so the admissibility
projections $E_d$ and the cost functional $\varepsilon$ are invariant under
$J$.  Therefore $J^{\otimes k}: B \mapsto \bar{B}$ is an
admissibility-preserving relabeling.

\textit{Step 4: Oriented distinction collapses.}\quad
By Steps~2--3, there exists an admissibility-preserving relabeling
(internal to the carrier, at zero capacity cost) that maps $B \mapsto \bar{B}$.
By Definition~\ref{def:oriented_composite}, the composite $B$ is
\emph{not} oriented: the ``forward'' and ``backward'' versions collapse
into the same equivalence class under admissible relabelings.

\textit{Step 5: Contradiction with hypotheses.}\quad
Hypothesis (B1) requires that the oriented distinction $B \ne \bar{B}$ be
robust.  Step~4 shows it is not robust if $V \cong V^*$.  The only escape
is to introduce an independent external grading channel that ``protects''
the $B / \bar{B}$ distinction from the $J$-identification.  But this
violates hypothesis~(B2) (minimality under A1: no additional admissibility
channel may be introduced solely to protect a distinction that the carrier
itself should sustain).

Therefore $V \not\cong V^*$: the carrier must be of complex type.
\end{proof}

\noindent\textit{Counterexample (pseudoreal confinement).}\quad
$G = \SU(2)$ with fundamental representation $V = \mathbb{C}^2$
(pseudoreal: $V \cong V^*$ via the $\varepsilon$-tensor, $J = i\sigma_2 K$
where $K$ is complex conjugation).  The ``baryon'' $B = \varepsilon_{ij}\, v^i w^j$
is a bilinear singlet.  $J^{\otimes 2}(B) = \bar{B}$, and the identification
is $G$-equivariant: there is no robust oriented distinction.
``Baryon'' and ``antibaryon'' in an $\SU(2)$ confining theory are the
same singlet reached by different paths.

\smallskip\noindent\textit{Positive example.}\quad
$G = \SU(3)$ with fundamental representation $V = \mathbb{C}^3$
(complex: $V \not\cong V^*$, since the Dynkin labels $[1,0]$ and
$[0,1]$ are distinct).  No $G$-equivariant antilinear $J: V \to V$ exists.
The baryon $B = \varepsilon_{ijk}\, v^i w^j x^k$ and the antibaryon
$\bar{B} = \varepsilon^{ijk}\, \bar{v}_i \bar{w}_j \bar{x}_k$ are
robustly distinct: no internal relabeling maps one to the other.

\smallskip\noindent\textit{Regularity conditions used:}
A1 (minimality clause, hypothesis B2), $T_M$ (admissibility independence,
which motivates hypothesis~B1), $L_{\mathrm{irr}}$ (stability of composites,
hypothesis~B1).  The proof itself is pure $G$-module theory --- the only
non-algebraic input is the admissibility-preservation criterion (Step~3),
which invokes equivariance and zero-cost relabeling.

\smallskip\noindent\textit{Consequence for Theorem~R.}\quad
$B1'$ closes the R1 derivation chain.  The ternary carrier ($k = 3$, from
the trilinear invariant requirement) must be complex type to sustain
oriented composites.  Combined with faithfulness and minimality
($\dim V = 3$ is the minimum for an irreducible trilinear invariant),
R1 is fully derived: a faithful complex 3-dimensional carrier.
\coderef{check\_B1\_prime}{gauge.py}


% ------------------------------------------------------------------
\subsection{Worked example: three-dimensional interface}
\label{sec:worked_example}

We illustrate $T_{7B}$ and $L_{\mathrm{irr,uniform}}$ with the simplest
nontrivial interface: $\dim(S_\Gamma) = 3$ with two one-dimensional
sectors and a one-dimensional pool.

\smallskip\noindent\textit{Setup.}\quad
$S_\Gamma = \mathbb{R}^3$ with standard basis $\{e_1, e_2, e_3\}$.
Sector decomposition ($T_{\mathrm{sep}}$):
$M_{d_1} = \spn(e_1)$, $M_{d_2} = \spn(e_2)$, $\Pi = \spn(e_3)$.
Admissibility projections:
$E_{d_1} = \diag(1,0,0)$, $E_{d_2} = \diag(0,1,0)$.
Individual admissibility costs: $\eps(d_1) = \eps(d_2) = \eps^* = 1$.

\smallskip\noindent\textit{Computing $\Delta$.}\quad
The pool $\Pi = \spn(e_3) \ne \{0\}$, so $L_\Delta$ applies.
A pool-supported perturbation $p = \alpha\, e_3$ with $|\alpha| > 0$
threatens joint admissibility (SP: physically meaningful DOF;
K2: positive cost).  Defense costs $\kappa_\Pi = \eps^* = 1$ (one pool
channel).  By K3: $\eps(\{d_1, d_2\}) = 1 + 1 + 1 = 3$.
Therefore $\Delta = 1 > 0$.

\smallskip\noindent\textit{Computing $\varepsilon_{\mathrm{ext}}$ and
verifying $T_{7B}$.}\quad
The external feasibility cost for a general perturbation
$v = (v_1, v_2, v_3)^T$ is specified by the cost matrix:
\[
  G \;=\; \begin{pmatrix}
    1 & \lambda & 0 \\
    \lambda & 1 & 0 \\
    0 & 0 & 1
  \end{pmatrix},
  \qquad
  \varepsilon_{\mathrm{ext}}(v) = v^T G\, v
  = v_1^2 + v_2^2 + v_3^2 + 2\lambda\, v_1 v_2,
\]
where $\lambda \in [0, 1)$ is the pool-mediated coupling between sectors
$M_{d_1}$ and $M_{d_2}$.  The diagonal entries are the individual sector
costs; the off-diagonal $\lambda$ is the cross-term from $L_\Delta$
(pool-mediated defense cost); the $e_3$-diagonal entry is the direct
pool cost.

\begin{varlist}
  \item $\lambda = 0$: sectors uncoupled; $\Delta = 0$ (no cross-term);
  $G$ is diagonal; the interface is classical.
  \item $\lambda > 0$: sectors coupled through $\Pi$; $\Delta > 0$;
  $G$ has off-diagonal structure.  The pool mediates an admissibility
  interaction that is invisible to individual-sector accounting.
\end{varlist}

\noindent $G$ is symmetric and positive-definite for $|\lambda| < 1$
(eigenvalues $1 \pm \lambda$ and $1$, all positive).
$\varepsilon_{\mathrm{ext}}(v) = v^T G\, v$ is manifestly a quadratic
form.  The polarization identity gives $g(u,v) = u^T G\, v$, which is
the bilinear form.  All three conclusions of $T_{7B}$ are verified:
\begin{enumerate}[label=(\roman*),nosep]
  \item $\varepsilon_{\mathrm{ext}}$ is a quadratic form with kernel $G$.
  \item $G$ is positive-definite $\Rightarrow$ non-degenerate.
  \item the diagonal part $B$ of $G$ is unique up to one positive
  scalar \emph{per sector}, matching $T_{7B}$(iii); the off-diagonal
  $\lambda$ is fixed by the pool coupling, not by the normalization
  gauge.
\end{enumerate}

The connection to $\Delta$: the superadditive surplus is
$\Delta = 2\lambda \cdot |v_1 v_2|$, which vanishes iff $\lambda = 0$
(no pool coupling).  The relation between $G$ and $T_{7B}$'s objects
must be stated carefully, because sector-orthogonality is exactly the
vanishing of cross-sector terms.  The sector-orthogonal bilinear form
$B$ ($L_\omega$) is the \emph{diagonal part} of $G$: its cross-sector
values are zero by construction ($B(v_1, v_2) = 0$ for $v_1, v_2$ in
distinct sectors), its diagonal entries are $B$-norms of sector
projections ($T_{\mathrm{adj}}$), and K3 gives exact additivity over
disjoint supports.  The off-diagonal $\lambda$ is \emph{not} $B$'s
cross-sector value (which is zero); it is the $L_\Delta$ pool-mediated
coupling, a distinct object added to $B$: at $\lambda = 0$ the cost
matrix is $B$ itself (the Sep baseline), and $\lambda > 0$ is the IJC
deformation carried by $\Pi$.

\smallskip\noindent\textit{Illustrating $L_{\mathrm{irr,uniform}}$.}\quad
Set $\lambda = 1/2$ and $C_\Gamma = 3$.
Let $G = \{d_1\}$ (gauge sector) and $R_\Gamma = \{d_2\}$ (gravitational
record), with costs $\eps(G) = \eps(R_\Gamma) = 1$.
The superadditive surplus is $\Delta(G, R_\Gamma) = 2\lambda = 1$.
Interface saturation: $s = (1 + 1)/3 = 2/3$.

Check hypothesis~(b) of $L_{\mathrm{irr,uniform}}$:
$s = 2/3 > 1 - \Delta/(C_\Gamma - \eps(R_\Gamma))
= 1 - 1/(3 - 1) = 1/2$.  Satisfied.

After a gauge transition $G_1 \to G_2$, the reversal cost is:
$\eps(\text{reverse}) \ge \eps(G_1) + \Delta = 1 + 1 = 2$.
The remaining budget is $C_\Gamma - \eps(R_\Gamma) - \eps(G_2)
= 3 - 1 - 1 = 1 < 2$.  The reversal is inadmissible.

Now set $C_\Gamma = 5$ (loose budget, $s = 2/5$):
$s = 2/5 < 1/2$.  Hypothesis~(b) fails.
Remaining budget $= 5 - 1 - 1 = 3 > 2$.  The reversal succeeds.
The gauge sector is reversible at this low saturation, as the theorem
predicts.


% ------------------------------------------------------------------
\subsection{Consolidation: Theorem~R prerequisites closed}
\label{sec:TR_closed}

With $T_{7B}$, $L_{\mathrm{irr,uniform}}$, and $B1'$ proved in this
supplement, the three carrier requirements of Theorem~R
(Paper~2, §3.5) are now derived entirely within Paper~1, Paper~2, and
their respective supplements:

\smallskip
\begin{center}
\small
\renewcommand{\arraystretch}{1.15}
\resizebox{\textwidth}{!}{\begin{tabular}{@{}p{3.1cm}p{4.7cm}p{3.5cm}p{1.2cm}@{}}
\toprule
\textbf{Requirement} & \textbf{Derivation chain} & \textbf{Key new lemma} & \textbf{Status} \\
\midrule
R1 (complex ternary) & $L_{\mathrm{nc}} \to T_{\mathrm{conf}} \to T_M \to B1'$ & $B1'$ (\S\ref{sec:B1prime}) & [P] \\
R2 (pseudoreal binary) & $L_{\mathrm{irr,uniform}} \to T_M \to \text{chirality}$ & $L_{\mathrm{irr,uniform}}$ (\S\ref{sec:Lirr_uniform}) & [P] \\
R3 (abelian grading) & admissibility completeness (R-EC-inv $+$ I-typing) + A1 minimality & self-contained in Paper~2 & [P $\mid$ R-EC-inv $+$ I-typing] \\
\bottomrule
\end{tabular}}
\end{center}

\smallskip\noindent No result in this section depends on Papers~3--7.
The formerly external dependency on ``Paper~7 ($B1'$)'' cited in
Paper~2's abstract and §3.4 is now resolved.


\newpage


% ══════════════════════════════════════════════════════════════

\section{The gauge-cost program}
\label{sec:cost_program}

The composed operational statement $C(G) = \dim(G)$ --- lower bound a
theorem on the operational-encoding premise, upper bound the
disjoint-anchor realizability premise (\S\ref{sec:nc_chain}) --- is
this program's current state.  The named open is a cost functor: a
derivation of both premises from the admissibility semantics, which
would upgrade the equality from premise-conditioned to derived.  The
subsection below records what capacity arguments alone currently
deliver on the asymptotic-freedom side of the same cost story.

\subsection{Can the sign be established from capacity alone?}
\label{sec:AF_from_capacity}

The imported-coefficient check (\S2 of \TSI{}) verifies $b_0 > 0$
using the imported one-loop formula.  The reviewer asks whether any part of the sign can
be established purely from capacity or positivity arguments, without
importing the formula from perturbative QFT.  The answer is partially
affirmative: the \emph{dominance of the gauge self-interaction term
over the matter screening term} can be given a capacity argument,
even though the precise coefficient $11/3$ cannot.

\begin{proposition}[Capacity sign pressure with imported group weights]%
\label{prop:AF_capacity_bound}
Let $G$ be a non-abelian gauge group.  Capacity accounting fixes the sign
of the two competing contributions---adjoint self-correlation is
anti-screening, matter channels are screening---but it does not by itself
fix the numerical weights multiplying those contributions.  The comparison
used here imports the standard adjoint-Casimir and Dynkin-index weights
($C_A$ and $T_R$).  With those imported weights and the derived matter
content, the non-abelian factors have positive one-loop coefficient.
\end{proposition}

\begin{proof}[Argument]
\textit{Step 1 (Self-interaction as capacity self-cost).}\quad
Each gauge generator occupies one admissibility channel
(Lemma~\ref{lem:gen_channel}).  Non-abelian gauge fields
self-interact: the structure constants $f^{abc}$ couple all $p$
channels to each other.  In the admissibility language, this
self-interaction is the cross-channel cost of maintaining $p$
mutually correlated admissibility channels at a single interface.  The
superadditive surplus $\Delta > 0$ ($L_\Delta$) ensures that this
self-cost is \emph{positive}: co-located non-abelian channels cost
more than the sum of their individual costs.  This positive
self-cost makes the effective coupling \emph{decrease} at short
distances (where more channels are simultaneously active), because
the admissibility budget is finite and the self-cost consumes capacity
that would otherwise be available for matter interactions.

\textit{Step 2 (Matter as screening).}\quad
Charged matter fields screen the gauge interaction: a virtual
fermion--antifermion pair partially cancels the gauge field,
reducing the effective coupling at long distances.  In the
admissibility language, each charged matter DOF provides an
additional channel through which admissibility can operate, partially
offsetting the capacity consumed by the gauge self-interaction.  The
screening contribution is proportional to $N_{\mathrm{matter}}$.

\textit{Step 3 (Counting argument).}\quad
For the derived matter content:
$N_{\mathrm{matter}}^{\SU(3)} = 12$ charged Weyl DOF per generation
$\times$ 3 generations $= 36$, while $p = \dim(\mathfrak{su}(3)) = 8$.
Naively, $N_{\mathrm{matter}} > p$, so the counting argument alone
does not settle the sign.  However, the self-interaction term carries
a group-theoretic enhancement: each generator couples to all $p - 1$
others through the adjoint representation, giving an effective
self-interaction count of $p \times C_A = p^2$ (where $C_A$ is the
quadratic Casimir of the adjoint).  For $\SU(3)$:
$p \times C_A = 8 \times 3 = 24$.  The matter screening, by contrast,
is weighted by $T_R$ (the index of the matter representation, which is
$1/2$ for the fundamental), giving an effective screening count of
$T_R \times N_f = (1/2) \times 12 = 6$ per generation, or $18$ for
three generations.  Since $24 > 18$, the self-interaction dominates.
\end{proof}

\noindent\textit{What this establishes.}\quad
The argument above does not derive the one-loop formula or the
precise coefficient $11/3$.  What it establishes is the \emph{physical
mechanism}: non-abelian self-interaction produces antiscreening
(capacity self-cost from $L_\Delta$), while matter produces screening,
and the self-interaction dominates for the derived matter content
because the adjoint Casimir $C_A$ enhances the self-interaction
count.  The dominance can be stated as a \emph{capacity inequality}:
\begin{equation}\label{eq:AF_capacity_ineq}
  p \cdot C_A \;>\; T_R \cdot N_f^{\mathrm{total}}
  \qquad\Longrightarrow\qquad b_0 > 0.
\end{equation}
For $\SU(3)$: $8 \times 3 = 24 > 6 \times 3 = 18$.  For $\SU(2)$:
$3 \times 2 = 6 > (1/2) \times 12 = 6$; the inequality is
\emph{saturated}, and the precise coefficient $11/3$ (rather than
some other number) is needed to establish positivity.  This is
the boundary where the capacity argument alone becomes insufficient
and the perturbative formula is required.

\smallskip\noindent\textit{What remains for Paper~4A.}\quad
Paper~4A ($L_{\mathrm{beta,capacity}}$) derives the one-loop
$\beta$-function coefficient, including the precise $11/3$ and $4/3$
factors, from the capacity structure: the $11/3$ arises from the
number of independent polarization states of a spin-1 field in $d = 4$
(from $T_8$) combined with the self-interaction structure of the
adjoint representation; the $4/3$ arises from the spin-$1/2$ screening
of a Weyl fermion.  The capacity-based derivation replaces the
standard Feynman-diagram calculation with an admissibility-cost
calculation, but arrives at the same coefficients.  Here we can only
establish the mechanism (antiscreening from self-interaction) and the
dominance condition (\ref{eq:AF_capacity_ineq}); the coefficients are
imported.




\section{Open-lane register}
\label{sec:open_lanes}

The program's named opens, each stated as a question with its current
grade boundary.  Heading the register:

\smallskip\noindent\textbf{1.\ The fixed-point program.}\quad Is the
(template, matter content, single-abelian-grading) triple the unique
simultaneous solution of the full constraint set?  The classification
core proves minimality among the enumerated candidates given the
template; the closure theorems verify abelian consistency given the
matter content.  A fixed-point theorem would close the loop:
uniqueness of the triple under the joint constraints, with no
consumption order.  Nothing banked asserts it today.

\smallskip\noindent\textbf{2.\ The cost functor.}\quad Derive the
operational cost premise pair (operational encoding; disjoint-anchor
realizability) from the admissibility semantics, upgrading
$C(G) = \dim(G)$ from premise-conditioned to derived
(\S\ref{sec:cost_program}).

\smallskip\noindent\textbf{3.\ The EC reading.}\quad Derive R-EC-inv
from the axioms, or prove it underivable in a stronger sense than the
current counter-construction (which shows A1 $+$ MD $+$
$T_{\mathrm{sep}}$ $+$ $L_{\mathrm{loc}}$ do not entail EC-injectivity
over species at shape
level)\coderef{check\_L\_EC\_inventory\_reading}{ec\_inventory\_reading.py}.

\smallskip\noindent\textbf{4.\ A ground for F6.}\quad Electromagnetic
completeness is an independent declared assumption of the scan, and
its independence from EC is a computed
theorem\coderef{check\_L\_F6\_not\_from\_EC}{ec\_inventory\_reading.py}.
The candidate lane is a record-carrier route: whether the unbroken
electromagnetic charge's separating role can be grounded in the record
semantics rather than declared.

\smallskip\noindent\textbf{5.\ Cap completeness.}\quad The scan's
enumeration caps are declared inputs; the proof obligations that no
admissible sub-45-DOF template exists outside them are owned in the
classification core's reproducibility contract and remain open.

\smallskip\noindent\textbf{6.\ The full-invariance hypothesis.}\quad
The gauge identification's equality leg consumes it
(Theorem~\ref{thm:continuous_gauge_inner_v34}).  Derive it --- anchor
data entering at block granularity for gauge purposes --- or restate
the downstream chain so that only the unconditional inclusion is
consumed everywhere (the present text already consumes only the
inclusion plus the exact sequence).

\smallskip\noindent One conditioning is named here so it is not
mistaken for an open lane: the IJC \emph{occupancy} is empirical by
design --- whether a physical interface presents branch-(IJC) records
is read off its correlation record, exactly as coupling constants are
read off theirs.  It is a named conditioning of the chain, not a gap
the program owes a proof for.


\newpage

\appendix


\section{Comparison with reconstruction programs}
\label{app:literature_found}
% ══════════════════════════════════════════════════════════════

The comparisons kept in this document are the ones that bear on the
foundational chain --- operational reconstructions, algebraic QFT, and
assumption inventories; the ``SM from principles'' comparisons that
bear on the classification claim live in \TSI's comparison appendix.

\subsection{Operational and categorical reconstructions}
\label{sec:op_reconstructions}

Several programs reconstruct quantum theory from operational or
information-theoretic axioms.  We summarize each and note where the
APF diverges.

\smallskip\noindent\textbf{Hardy (2001)~\cite{Hardy2001}.}\quad
Five ``reasonable axioms'' (probabilities, simplicity, subspaces,
composite systems, continuous reversibility) derive finite-dimensional
quantum theory over $\mathbb{C}$.  The axioms are stated at the
kinematics level: they constrain the state space and measurement
structure but do not determine gauge content, matter spectrum, or
dynamics.  APF shares the ``finite capacity'' starting point
(Hardy's axiom 1 is a form of finite dimensionality) but pushes
through to gauge group and matter content via non-closure, which
Hardy does not address.

\smallskip\noindent\textbf{Chiribella--D'Ariano--Perinotti
(2011)~\cite{CDP2011}.}\quad
Six axioms (causality, perfect distinguishability, ideal compression,
local distinguishability, pure conditioning, purification) derive QT.
The purification axiom is the key differentiator from classical
probability.  The output is a general probabilistic theory isomorphic
to finite-dimensional quantum mechanics.  Like Hardy, CDP stops at
kinematics.  The APF's constitutive principle SP (substrate faithfulness)
plays a role analogous to CDP's local distinguishability, and K3
(forced additivity) parallels their ideal compression, but the APF
does not assume purification --- instead, it derives the Born rule
from the admissibility algebra via Busch's generalized Gleason theorem.


\begin{remark}[Complex structure and reconstruction axioms]
\label{rem:complex_reconstruction_v34}
The APF's $B1'$ result should be compared with local-tomography and
continuous-reversibility reconstructions of complex quantum theory.  Those
programs typically isolate $\mathbb C$ by assuming how composites behave
(e.g., local tomography, purification, or continuous reversible dynamics).  APF
uses a different but related composite axiom: oriented composite robustness.
Real carriers fail to preserve the needed oriented charged distinction, while
quaternionic carriers add an anti-linear/reality structure not forced by the
finite record ledger.  Thus $B1'$ is not a replacement proof of all complex
quantum theory; it is the APF-specific gate selecting complex charged modules
for the gauge/matter theorem.
\end{remark}

\smallskip\noindent\textbf{H\"ohn (2017)~\cite{Hohn2017}.}\quad
Reconstructs qubit quantum theory from ``questions'' (binary tests)
with finite information capacity.  The closest cousin to APF's
starting point: finite capacity $\to$ quantum structure.  H\"ohn's
framework is carefully delimited to qubits and does not attempt gauge
content.  APF can be viewed as an extension of H\"ohn's program that
asks: what happens when you push the ``finite information'' idea beyond
qubit reconstruction to the full field content?

\smallskip\noindent\textbf{Effectus theory (Westerbaan, van de Wetering).}\quad
Categorical framework using ``effectuses'' to axiomatize quantum-like
theories.  Sequential product axioms force operator-algebraic structure.
The APF's admissibility algebra (\S\ref{sec:nc_chain}) can be compared
to the sequential product in effectus theory: both derive
non-commutativity from the structure of sequential measurements.  The
effectus approach is more general (it classifies all effectuses, not
just the quantum one); the APF is more specific (it selects the
quantum case via A1 and derives gauge content).

\smallskip\noindent\textbf{Tull (2019); Spekkens (2007).}\quad
Tull derives operator algebras from categorical CPM (completely positive
maps) structure.  Spekkens constructs a toy model from epistemic
restrictions.  Both illuminate the boundary between classical and
quantum but do not derive field content or gauge structure.

\smallskip\noindent\textbf{Finite-algebra structures in quantum error
correction.}\quad A structurally close neighbor from a different
tradition: in operator-algebra quantum error correction and categorical
code frameworks, the correctable/logical structure of a finite code is
organized by a finite $*$-algebra, and the \emph{inner} automorphisms
of the code algebra are exactly the locally implementable
(``gauge-like'') transformations, while outer automorphisms act as
logical symmetries.  That is the same inner/outer split this
supplement derives operationally --- gauge $=$ inner $=$
implementable-at-no-observable-cost, kinematic/logical $=$ outer ---
arrived at from error-correction requirements rather than capacity
minimization.  The convergence is evidence that identifying the
physical gauge group with $\mathrm{Inn}(\mathcal{A})$ of a finite
operational algebra is natural in finite settings generally, not an
artifact of the admissibility axioms.  The APF adds what the QEC
tradition does not attempt: a selection principle (PLEC) over the
algebras themselves, from which the \emph{particular} finite algebra
$M_3 \oplus M_2 \oplus \mathbb{C}$ is argued to be the minimum-cost
choice.


\subsection{Locally covariant QFT and the General Boundary Formulation}
\label{sec:lcqft_oeckl}

Two programs are particularly relevant to the APF because they
share structural elements (local algebras, causality, boundaries)
while differing fundamentally on infinite-dimensionality.

\smallskip\noindent\textbf{Locally covariant QFT
(Brunetti--Fredenhagen--Verch, 2003).}\quad
LCQFT axiomatizes QFT as a functor from the category of globally
hyperbolic Lorentzian manifolds (with isometric embeddings as morphisms)
to the category of $*$-algebras (with injective $*$-homomorphisms).
The key axioms are:
\begin{itemize}[nosep]
  \item \textit{Locality}: spacelike-separated regions have commuting
  algebras.
  \item \textit{Covariance}: the assignment $\mathcal{O} \mapsto
  \mathcal{A}(\mathcal{O})$ is functorial under isometric embeddings.
  \item \textit{Time-slice}: if $\mathcal{O}_1 \hookrightarrow
  \mathcal{O}_2$ contains a Cauchy surface of $\mathcal{O}_2$, then
  $\mathcal{A}(\mathcal{O}_1) \cong \mathcal{A}(\mathcal{O}_2)$.
\end{itemize}
The LCQFT framework works with \emph{infinite-dimensional} algebras
(typically type~III factors for local regions) and
\emph{presupposes} a background Lorentzian manifold.

The APF differs from LCQFT on three structural points:
\begin{enumerate}[label=(\alph*),nosep]
  \item \textit{Finite vs.\ infinite dimensions.}\quad LCQFT's local
  algebras are infinite-dimensional by construction (they are obtained
  via the GNS representation of a state on a $C^*$-algebra of
  smeared fields).  The APF's local algebras are finite-dimensional
  (A1: $\dim(\mathcal{H}_\omega) < \infty$).  The reconciliation is
  via the controlled limit (Proposition~\ref{prop:typeIII_limit}):
  the APF's type~I algebras approximate LCQFT's type~III algebras
  in the $C/\varepsilon^* \to \infty$ limit, with the structural
  predictions (gauge group, matter content) stable under the limit
  (Theorem~\ref{thm:structural_stability}).
  \item \textit{Spacetime derived vs.\ presupposed.}\quad LCQFT takes
  a Lorentzian manifold as input; the APF derives it
  (exited spacetime bridge, Paper~6: $L_{\mathrm{irr}} \to \prec \to M \to
  g_{\mu\nu}$).  The APF can therefore be viewed as providing a
  \emph{foundation} for LCQFT: it derives the background manifold
  that LCQFT assumes, and the local algebras that LCQFT assigns to
  regions emerge as the admissibility algebras at interfaces.
  \item \textit{Causality.}\quad Both frameworks implement causality
  via commutativity of spacelike-separated algebras.  In LCQFT, this
  is an axiom; in the APF, it is derived
  (Lemma~\ref{lem:loc_commut}).  The LCQFT time-slice axiom has an
  APF analogue: if $\Gamma_1 \subseteq \Gamma_2$ and $\Gamma_1$
  ``contains a Cauchy surface'' (i.e., the admissibility structure at
  $\Gamma_1$ determines the admissibility structure at $\Gamma_2$),
  then $\mathcal{A}(\Gamma_1) \cong \mathcal{A}(\Gamma_2)$.  This
  is a consequence of $L_{\mathrm{loc}}$ and the deterministic
  evolution of the admissibility structure.
\end{enumerate}

\smallskip\noindent\textbf{Oeckl's Positive Formalism / General
Boundary Formulation (2003--2016).}\quad
Oeckl's program assigns quantum amplitudes to \emph{boundaries} of
spacetime regions, rather than to states on Cauchy surfaces.  The
formalism uses:
\begin{itemize}[nosep]
  \item \textit{Boundary Hilbert spaces}: each boundary $\Sigma$
  (not necessarily a Cauchy surface) carries a Hilbert space
  $\mathcal{H}_\Sigma$.
  \item \textit{Amplitude maps}: each spacetime region $M$ bounded by
  $\Sigma$ defines a linear map $\rho_M: \mathcal{H}_\Sigma \to
  \mathbb{C}$ (or more generally a positive map).
  \item \textit{Composition}: gluing regions along shared boundaries
  composes amplitude maps.
\end{itemize}
The General Boundary Formulation (GBF) generalizes this to
non-globally-hyperbolic spacetimes and is compatible with quantum
gravity scenarios where the background spacetime is not fixed.

The APF's interface structure has a striking structural parallel to
Oeckl's boundary formalism:
\begin{enumerate}[label=(\alph*),nosep]
  \item \textit{Interfaces as boundaries.}\quad An APF interface
  $\Gamma$ is the boundary between two admissibility regions.  The
  admissibility algebra $\mathcal{A}(\Gamma)$ plays the role of
  Oeckl's boundary Hilbert space $\mathcal{H}_\Sigma$ (in the
  finite-dimensional case, the algebra and its GNS Hilbert space
  are equivalent objects).
  \item \textit{Cost kernel as amplitude.}\quad The admissibility cost
  $\varepsilon_{\mathrm{ext}}(v) = B(v,v)$ ($T_{7B}$) is a positive
  functional on the interface sector space, analogous to Oeckl's
  positive amplitude $\rho_M$.  The positivity of $B$ corresponds to
  Oeckl's positivity condition, which ensures physical probabilities
  are non-negative.
  \item \textit{Composition via pool coupling.}\quad When two
  interfaces share a pool $\Pi$ ($L_\Delta$), the admissibility
  structure composes: the combined cost is superadditive
  ($\varepsilon(v_1 + v_2) \ge \varepsilon(v_1) + \varepsilon(v_2)$).
  This parallels Oeckl's composition of amplitude maps via gluing.
\end{enumerate}
The key differences:
\begin{itemize}[nosep]
  \item Oeckl works with infinite-dimensional Hilbert spaces and
  does not impose a capacity bound.  The APF's A1 restricts to finite
  dimensions.
  \item Oeckl's formalism is \emph{kinematic} (it assigns amplitudes
  but does not derive gauge content or matter spectrum).  The APF
  pushes through to field content via the admissibility algebra's
  Wedderburn decomposition.
  \item The GBF is compatible with indefinite causal structure
  (useful for quantum gravity); the APF derives a definite causal
  structure from $L_{\mathrm{irr}}$
  (exited causal-order construction; Paper~6, BA1).
\end{itemize}

\smallskip\noindent\textbf{Masanes--M\"uller (2011).}\quad
Derives finite-dimensional quantum theory from four postulates:
continuous reversibility, tomographic locality (joint states
determined by local measurements), existence of an information unit,
and the subspace axiom.  The output is quantum theory over
$\mathbb{R}$, $\mathbb{C}$, or $\mathbb{H}$ (with $\mathbb{C}$
selected by tomographic locality + bit symmetry).  The APF shares
the finite-dimensionality starting point and the field selection
($\mathbb{F} = \mathbb{C}$ from T2b), but arrives at $\mathbb{C}$
via a different mechanism (Skolem--Noether applied to the admissibility
algebra, not tomographic locality).  Like Hardy and CDP,
Masanes--M\"uller stops at kinematics and does not address gauge
content.


\subsection{Systematic comparison: infinite-dimensionality and causality}
\label{sec:systematic_comparison}

The reviewer asks for a systematic contrast on two axes ---
infinite-dimensionality and causality --- across the relevant
programs.  The following table provides this.

\smallskip
\begin{center}\footnotesize
\renewcommand{\arraystretch}{1.25}
\begin{tabular}{@{}p{2.4cm}p{3.2cm}p{3.2cm}p{3.5cm}@{}}
\toprule
\textbf{Program} & \textbf{Local algebra type} &
  \textbf{Causality treatment} & \textbf{APF relationship} \\
\midrule
AQFT (Haag--Kastler) &
  Type~III$_1$ factors (Reeh--Schlieder) &
  Axiom: $[\mathcal{A}(\mathcal{O}_1), \mathcal{A}(\mathcal{O}_2)] = 0$ for spacelike $\mathcal{O}_1, \mathcal{O}_2$ &
  APF recovers type~III in $n \to \infty$ limit; derives commutativity from $L_{\mathrm{loc}}$ \\[3pt]
LCQFT (Brunetti--Fredenhagen--Verch) &
  Infinite-dim.\ $*$-algebras; functorial on Lorentzian manifolds &
  Presupposes globally hyperbolic manifold; time-slice axiom &
  APF derives the manifold; provides a finite-dim.\ foundation \\[3pt]
Oeckl (GBF) &
  Infinite-dim.\ boundary Hilbert spaces &
  Compatible with indefinite causal structure &
  Structural parallel (interfaces $\sim$ boundaries); APF adds A1 finiteness and derives causal order \\[3pt]
Hardy (2001) &
  Finite-dim.\ (axiom 1) &
  Not addressed (kinematics only) &
  Shared starting point; APF extends to gauge content \\[3pt]
CDP (2011) &
  Finite-dim.\ (output) &
  Causality axiom (no signaling) &
  APF derives no-signaling from $L_{\mathrm{loc}}$; pushes beyond kinematics \\[3pt]
Masanes--M\"uller &
  Finite-dim.\ (postulate) &
  Not directly addressed &
  Shared $\mathbb{F} = \mathbb{C}$ result; different mechanism \\[3pt]
Connes NCG &
  Finite-dim.\ internal algebra; $\infty$-dim.\ external &
  Lorentzian structure from spectral triple &
  APF derives the algebra that NCG postulates \\[3pt]
DR reconstruction &
  Infinite-dim.\ (type~III in AQFT context) &
  Inherited from AQFT axioms &
  APF uses DR at finite $n$; structural parallel for abelian factors \\[3pt]
APF &
  Type~I$_n$ ($M_n(\mathbb{C})$); type~III recovered as $n \to \infty$ &
  Derived from $L_{\mathrm{irr}}$ (exited; Paper~6, BA1); HKM/Malament for signature &
  --- \\
\bottomrule
\end{tabular}
\end{center}

\smallskip\noindent Three patterns emerge from the comparison:

\textit{(i) Finite vs.\ infinite dimensions is a foundational choice,
not a technical detail.}\quad Programs that start with infinite
dimensions (AQFT, LCQFT, Oeckl, DR) must contend with type~III
algebras, ultraviolet divergences, and the absence of a faithful
trace.  Programs that start with finite dimensions (Hardy, CDP,
Masanes--M\"uller, APF) avoid these issues but must explain how the
infinite-dimensional continuum is recovered.  The APF's controlled
limit (Proposition~\ref{prop:typeIII_limit}) and structural stability
theorem (Theorem~\ref{thm:structural_stability}) provide this
explanation; the kinematic reconstruction programs do not address the
question (they stop at finite-dimensional QT and do not discuss the
continuum limit).

\textit{(ii) Causality is either axiomatized or derived, never
both.}\quad AQFT and CDP axiomatize causality (as commutativity of
spacelike algebras or no-signaling); the APF derives it from
$L_{\mathrm{irr}}$ (irreversibility $\to$ causal order) and
$L_{\mathrm{loc}}$ (independent interfaces $\to$ commutativity).
Oeckl's GBF is compatible with indefinite causality, making it the
most flexible --- but at the cost of not determining the causal
structure.  The APF's derivation of causality from a non-causal
starting point (finite admissibility capacity has no intrinsic temporal
or causal content) is the closest parallel to the causal-set program
(Sorkin, 1991), though the APF derives a smooth manifold rather than
working with a discrete causal set.

\textit{(iii) Only the APF and Connes NCG derive gauge content.}\quad
All other programs in the table stop at the kinematics of quantum
theory (state space, measurement structure, composition rules).  The
passage from kinematics to dynamics and gauge content requires
additional structure: in NCG, the spectral triple; in the APF, the
admissibility algebra's Wedderburn decomposition + non-closure +
anomaly cancellation.  The APF's advantage over NCG is that it
\emph{derives} the algebra, whereas NCG postulates it; NCG's
advantage is that it produces the full SM Lagrangian (including Yukawa
couplings and the Higgs potential), which the APF defers to Paper~4A.


\subsection{Assumption inventory comparison}
\label{sec:assumption_comparison}

\smallskip
\begin{center}
\small
\begin{tabular}{@{}lccl@{}}
\toprule
\textbf{Program} & \textbf{Inputs} & \textbf{Derives} & \textbf{Stops at} \\
\midrule
Hardy (2001) & 5 axioms & QT kinematics & State space \\
CDP (2011) & 6 axioms & QT kinematics & State space \\
Masanes--M\"uller (2011) & 4 postulates & QT over $\mathbb{C}$ & State space \\
H\"ohn (2017) & Finite info & Qubit QT & Qubit \\
LCQFT (BFV 2003) & Functorial axioms & QFT on curved spacetime & Assumes manifold \\
Oeckl (GBF 2003) & Boundary axioms & QT amplitudes & Kinematics \\
Connes NCG & Spectral triple & SM Lagrangian & Assumes algebra \\
APF & A1+MD+A2+BW+SP+K3 & QT + gauge chain at banked grades & Classification inputs (\TSI) \\
\bottomrule
\end{tabular}
\end{center}

\smallskip\noindent The APF's named assumption count (the four PLEC components A1, MD, A2, BW; plus constitutive semantics SP, K3, FD1--FD6) is numerically comparable to Hardy's 5 or CDP's 6,
though a direct comparison of ``assumption weight'' is imprecise: the
APF's constitutive definitions (FD1--FD6) under PLEC's four
constitutive features (A1+MD+A2+BW) carry more internal structure
than any single-axiom presentation suggests.  The significant
difference is scope: the APF's
downstream chain extends toward gauge content and the matter spectrum
--- as a graded program with named premises, closed by the
classification core of \TSI\ --- and, at bridge grade only, toward
cosmological readings; the other programs do not attempt this scope.


% ══════════════════════════════════════════════════════════════
\smallskip\noindent\textbf{DHR/DR categorical boundary.}\quad
The finite APF admissibility net is not asserted to be a type-III Haag--Kastler net.  The functorial comparison is: a finite interface category of admissible localized endomorphisms maps to a tensor subcategory of transportable sectors; the finite field algebra generated by those sectors has compact internal automorphism group $\mathrm{Inn}(\mathcal A)$; the continuum DHR/DR theorem is recovered only after adding the split/nuclearity/scaling assumptions stated in Proposition~\ref{prop:continuum_plan}.  Thus APF uses DR as a structural analogue and limiting target, not as a hidden premise inside the finite derivation.



\section{Self-contained proofs of imported results}
\label{app:imported}
% ══════════════════════════════════════════════════════════════

This appendix provides complete proofs of every result imported from the
Paper~1 Supplement that is used in this document.  With this appendix,
no external document is needed to verify any theorem in this supplement.

\smallskip\noindent\textit{Input assumptions.}\quad All proofs below
use only: A1 (finite admissibility capacity $0 < C(\Gamma) < \infty$),
MD (uniform per-test cost floor $\varepsilon \ge \mu^* > 0$),
K3 (additivity on disjoint supports), SP (no null directions in the
physical quotient), and the formal definitions FD1--FD6.  No result
from Papers~2--7 is used.


\subsection{$L_{\varepsilon^*}$: Uniform cost floor}
\label{app:Leps}

\begin{lemma_app}[$L_{\varepsilon^*}$]
There exists $\varepsilon^* > 0$ such that
$\varepsilon(d) \ge \varepsilon^*$ for all $d \in \mathcal{D}$.
\end{lemma_app}

\begin{proof}
By MD: $\varepsilon(d) \ge \mu^* \cdot n(d) \ge \mu^* > 0$ for every
$d$ (since $n(d) \ge 1$).  Set $\varepsilon^* := \mu^*$.
\end{proof}


\subsection{$T_{\mathrm{sep}}$: Sector decomposition}
\label{app:Tsep}

\begin{theorem_app}[$T_{\mathrm{sep}}$]
For distinctions $d_1, d_2 \in \mathcal{D}(\Gamma)$ with anchor sets
$M_{d_1}, M_{d_2} \subseteq S_\Gamma$:
\[
  M_{d_1} \cap M_{d_2} = \{0\}
  \quad\Longleftrightarrow\quad
  \varepsilon(\{d_1, d_2\}) = \varepsilon(d_1) + \varepsilon(d_2).
\]
The biconditional is stated at empty residual pool, $\Pi = \{0\}$; when
$\Pi \neq \{0\}$ the joint cost carries, in addition, the superadditive
surplus $\kappa_\Pi$ of $L_\Delta$ (Appendix~\ref{app:LDelta}).
For any maximal antichain $\mathcal{B}$ of pairwise independently
admissible distinctions, the substrate decomposes as
$S_\Gamma = \bigoplus_{d \in \mathcal{B}} M_d \oplus \Pi$,
where $\Pi$ is the residual pool.
\end{theorem_app}

\begin{proof}
($\Rightarrow$)\quad If $M_{d_1} \cap M_{d_2} = \{0\}$, perturbations
threatening $d_1$ and $d_2$ have disjoint supports (FD3: anchor-set
locality).  The joint defense decomposes into independent defenses on
$M_{d_1}$ and $M_{d_2}$.  By K3 (additivity on disjoint supports):
$\varepsilon(\{d_1,d_2\}) = \varepsilon(d_1) + \varepsilon(d_2)$.

($\Leftarrow$, contrapositive)\quad If $v \in M_{d_1} \cap M_{d_2}$
with $v \ne 0$, then defending $d_1$ and $d_2$ separately each commits
capacity to $\mathrm{span}\{v\}$; defending jointly commits once.
By SP and K2, $\kappa(\mathrm{span}\{v\}) > 0$.  Write
$M_{d_i} = \mathrm{span}\{v\} \oplus M_{d_i}'$.  By K3 on the three
disjoint subspaces:
\[
  \varepsilon(\{d_1,d_2\}) \le \kappa(\mathrm{span}\{v\})
  + \kappa(M_{d_1}') + \kappa(M_{d_2}')
  < \varepsilon(d_1) + \varepsilon(d_2).
\]
The direct-sum decomposition follows by iterating the biconditional
over a maximal antichain.
\end{proof}

\noindent\textit{Uses:} A1, K3, FD3, SP.  No inner product or linear
map structure.


\subsection{$L_\Delta$: Superadditivity}
\label{app:LDelta}

\begin{theorem_app}[$L_\Delta$]
For co-located distinctions $d_1, d_2$ at $\Gamma$ with
$M_{d_1} \cap M_{d_2} = \{0\}$ and $\Pi \ne \{0\}$:
$\Delta(d_1,d_2) := \varepsilon(\{d_1,d_2\}) - \varepsilon(d_1)
- \varepsilon(d_2) > 0$.
\end{theorem_app}

\begin{proof}
\textit{Step 1.}\quad The joint perturbation class
$\mathcal{P}(\{d_1,d_2\})$ strictly contains
$\mathcal{P}(d_1) \cup \mathcal{P}(d_2)$: perturbations
$p$ with $\mathrm{supp}(p) \subseteq \Pi$ are in
$\mathcal{P}(\{d_1,d_2\})$ but not in either individual class
(their anchors are disjoint from $\Pi$ by $T_{\mathrm{sep}}$).
These perturbations change physically meaningful DOF (SP) at positive
cost (K2).

\textit{Step 2.}\quad A defense confined to $M_{d_1} \oplus M_{d_2}$
acts as the identity on $\Pi$ and cannot resist $\Pi$-supported
perturbations.

\textit{Step 3.}\quad Any admissible joint defense must commit positive
capacity $\kappa_\Pi > 0$ in $\Pi$ (by K2: resisting a non-null
perturbation requires positive capacity).

\textit{Step 4.}\quad $M_{d_1}$, $M_{d_2}$, $\Pi$ are pairwise disjoint
($T_{\mathrm{sep}}$).  By K3:
$\varepsilon(\{d_1,d_2\}) = \varepsilon(d_1) + \varepsilon(d_2)
+ \kappa_\Pi$.

\textit{Step 5.}\quad $\Delta = \kappa_\Pi > 0$.
\end{proof}

\noindent\textit{Uses:} A1, K2, K3, $T_{\mathrm{sep}}$, SP.
Pre-algebraic.


\subsection{$L_\omega$: Sector-orthogonal bilinear form}
\label{app:Lomega}

\begin{lemma_app}[$L_\omega$]
Given $T_{\mathrm{sep}}$ and K3, there exists a positive-definite
symmetric bilinear form $B: V \times V \to \mathbb{R}$ with
$B(u_d, v_{d'}) = 0$ for $d \ne d'$ and $B(u_d, v_\Pi) = 0$ for
all $d$.  The inter-sector orthogonality is forced; the within-sector
forms are free (normalization gauge).
\end{lemma_app}

\begin{proof}
\textit{Construction.}\quad Choose any positive-definite form $B_d$ on
each sector $M_d$ and $B_\Pi$ on $\Pi$.  Define
$B(u,v) := \sum_d B_d(u_d, v_d) + B_\Pi(u_\Pi, v_\Pi)$
with zero cross-terms.

\textit{$B$ is well-defined, symmetric, positive-definite} by
construction (direct sum of positive-definite forms).

\textit{K3 forces cross-term vanishing.}\quad For $u_d \in M_d$ and
$v_{d'} \in M_{d'}$ with $d \ne d'$, K3 gives
$\|u_d + v_{d'}\|_B^2 = \|u_d\|_B^2 + \|v_{d'}\|_B^2$ (disjoint
supports).  Expanding: $B(u_d, u_d) + B(v_{d'}, v_{d'}) +
2B(u_d, v_{d'}) = B(u_d, u_d) + B(v_{d'}, v_{d'})$, so
$B(u_d, v_{d'}) = 0$.  The same argument gives
$B(u_d, v_\Pi) = 0$.
\end{proof}

\noindent\textit{Uses:} $T_{\mathrm{sep}}$, K3.


\subsection{$T_{\mathrm{adj}}$: Self-adjoint admissibility projections}
\label{app:Tadj}

\begin{theorem_app}[$T_{\mathrm{adj}}$]
With respect to $B$ from $L_\omega$, each sector map $E_d$ is a
$B$-orthogonal projection: $E_d^\dagger = E_d$, $E_d^2 = E_d$.
\end{theorem_app}

\begin{proof}
$\mathrm{im}(E_d) = M_d$ and
$\ker(E_d) = (\bigoplus_{d' \ne d} M_{d'}) \oplus \Pi$.
By $L_\omega$: $M_d \perp_B \ker(E_d)$.  Therefore $E_d$ is the
$B$-orthogonal projection onto $M_d$:
$B(E_d u, v) = B(u_d, v_d) = B(u, E_d v)$, so $E_d^\dagger = E_d$.
\end{proof}


\subsection{$L_{\mathrm{loc}}$: Locality (budget additivity)}
\label{app:Lloc}

\begin{theorem_app}[$L_{\mathrm{loc}}$]
For independent interfaces $\Gamma_1, \Gamma_2$ with disjoint
substrates $S_{\Gamma_1} \cap S_{\Gamma_2} = \emptyset$:
$C(\Gamma_1 \cup \Gamma_2) = C(\Gamma_1) + C(\Gamma_2)$.
\end{theorem_app}

\begin{proof}
Every distinction $d$ has its anchor entirely within $S_{\Gamma_1}$ or
$S_{\Gamma_2}$ (FD3: anchor-set locality).  The cost sum partitions:
$\sum_d \varepsilon(d) = \sum_{d \in \mathcal{D}(\Gamma_1)}
\varepsilon(d) + \sum_{d \in \mathcal{D}(\Gamma_2)} \varepsilon(d)
\le C(\Gamma_1) + C(\Gamma_2)$.  The bound is tight (each interface's
budget can be independently saturated).
\end{proof}

\noindent\textit{Uses:} A1, FD3.  Does not use MD, $L_{\varepsilon^*}$,
or any cost-functional structure.


\subsection{$L_{\mathrm{blk}}$: Block-diagonal defense failure}
\label{app:Lblk}

\begin{theorem_app}[$L_{\mathrm{blk}}$]
Let $\Delta(d_1,d_2) > 0$.  Then no idempotent with
$\mathrm{im}(P) \subseteq M_{d_1} \oplus M_{d_2}$ is an admissible
joint defender.
\end{theorem_app}

\begin{proof}
Such a $P$ satisfies $P|_\Pi = 0$ (by idempotency and the direct-sum
structure).  It therefore has zero defense against $\Pi$-supported
perturbations.  But $\Delta > 0$ requires positive $\Pi$-defense
($L_\Delta$, Step~3).  Contradiction.
\end{proof}


\subsection{$T_{\mathrm{alg}}$: Noncommutativity}
\label{app:Talg}

\begin{theorem_app}[$T_{\mathrm{alg}}$]
If $\Delta(d_1,d_2) > 0$, the full admissibility algebra
$\mathcal{A} := \mathrm{Alg}_\mathbb{R}(\{E_d\} \cup
\{E_{\{d_1,d_2\}}\})$ is noncommutative.
\end{theorem_app}

\begin{proof}
\textit{Step 1.}\quad The sector projections $\{E_d\}$ commute and
generate $\mathcal{A}_{\mathrm{diag}} \cong \mathbb{R}^{|\mathcal{B}|}$.

\textit{Step 2.}\quad By $L_{\mathrm{blk}}$, the minimum-cost joint
defender $E_{\{d_1,d_2\}} \notin \mathcal{A}_{\mathrm{diag}}$.
Define $F_\Pi := E_{\{d_1,d_2\}} - E_{d_1} - E_{d_2} \ne 0$.

\textit{Step 3.}\quad There exists $v \in M_{d_1}$ with
$E_{\{d_1,d_2\}} v = v + \pi_v$, $\pi_v \in \Pi \setminus \{0\}$
(otherwise $E_{\{d_1,d_2\}}$ is block-diagonal, contradicting Step~2).
Then:
$[E_{d_1}, F_\Pi]\,v = E_{d_1}(\pi_v) - F_\Pi(v) = 0 - \pi_v
= -\pi_v \ne 0$,
using $E_{d_1}|_\Pi = 0$ ($T_{\mathrm{sep}}$).

\textit{Step 4.}\quad $[E_{d_1}, F_\Pi] \ne 0$ in $\mathrm{End}(V)$.
Any faithful $*$-homomorphism to a commutative algebra would map this
to zero, contradicting faithfulness.  Therefore $\mathcal{A}$ is
noncommutative.
\end{proof}

\noindent\textit{Uses:} $T_{\mathrm{sep}}$, $L_\Delta$,
$L_{\mathrm{blk}}$, FD5a.  Pre-Hilbert.


\subsection{$L_{\mathrm{nc}}$: Non-closure (superadditivity at shared
interfaces)}
\label{app:Lnc}

\begin{lemma_app}[$L_{\mathrm{nc}}$]
At any shared interface $\Gamma$ with $\Pi \ne \{0\}$ and
co-located distinction sets $G, R_\Gamma$:
$\varepsilon(G \cup R_\Gamma) > \varepsilon(G)
+ \varepsilon(R_\Gamma)$.
\end{lemma_app}

\begin{proof}
This is $L_\Delta$ applied to the distinction sets $G$ and $R_\Gamma$:
the hypotheses of $L_\Delta$ (disjoint anchors, non-empty pool) are
satisfied by the co-location conditions of the shared interface.
\end{proof}


\subsection{$L_{\mathrm{irr}}$: Irreversibility}
\label{app:Lirr}

\begin{theorem_app}[$L_{\mathrm{irr}}$]
At any shared interface $\Gamma$ with $\Pi \ne \{0\}$, there exist
admissibility transitions that are locally irreversible: the capacity
committed in the transition cannot be recovered by any operation
confined to a single sector.
\end{theorem_app}

\begin{proof}
\textit{Step 1.}\quad By $L_\Delta$ (\S\ref{app:LDelta}), co-located
admissibility at $\Gamma$ commits a superadditive surplus
$\Delta = \kappa_\Pi > 0$ in the pool sector $\Pi$.

\textit{Step 2.}\quad By $L_{\mathrm{loc}}$ (\S\ref{app:Lloc}),
admissibility at separate interfaces has independent budgets.  A
single-sector observer (operating within $M_{d_1}$ alone) has access
only to the capacity committed in $M_{d_1}$, not to the pool
commitment $\kappa_\Pi$.

\textit{Step 3.}\quad The pool commitment $\kappa_\Pi$ is a
cross-sector correlation: it is committed in $\Pi$, which is shared
between the sectors.  No operation confined to $M_{d_1}$ can access
$\Pi$ (by the direct-sum structure of $T_{\mathrm{sep}}$:
$M_{d_1} \cap \Pi = \{0\}$).  Therefore the pool commitment is locally
unrecoverable.

\textit{Step 4.}\quad A transition that increases $\kappa_\Pi$ (e.g.,
by changing the gauge-sector distinctions at a shared interface) commits
capacity that is locally unrecoverable.  This is irreversibility: the
forward transition costs more to reverse than the budget allows, because
the reversal would require freeing pool capacity that no single-sector
operation can access.
\end{proof}

\noindent\textit{Uses:} $L_\Delta$, $L_{\mathrm{loc}}$,
$T_{\mathrm{sep}}$.  The argument is pre-algebraic: no Hilbert space
or GNS structure is invoked.

\medskip\noindent\textit{Summary.}\quad The ten results above
($L_{\varepsilon^*}$, $T_{\mathrm{sep}}$, $L_\Delta$, $L_\omega$,
$T_{\mathrm{adj}}$, $L_{\mathrm{loc}}$, $L_{\mathrm{blk}}$,
$T_{\mathrm{alg}}$, $L_{\mathrm{nc}}$, $L_{\mathrm{irr}}$) constitute
the complete set of results imported from the Paper~1 Supplement.  Each
is proved from A1, MD, K3, SP, and FD1--FD6 alone.  With these proofs
in hand, every theorem in this supplement can be verified end-to-end
without consulting any external document.




% ══════════════════════════════════════════════════════════════
\begin{thebibliography}{99}

\bibitem{BrookePaper1}
E.\,S.~Brooke, ``Paper~1: The Enforceability of Distinction,''
Admissibility Physics, 2025.
Submitted to \textit{Foundations of Physics}.

\bibitem{BrookeSupplement}
E.\,S.~Brooke, ``Technical Supplement: End-to-End Derivation from Finite
Enforcement Capacity to Quantum Structure,'' companion to Paper~1, 2025.

\bibitem{Humphreys1972}
J.\,E.~Humphreys, \textit{Introduction to Lie Algebras and Representation
Theory}, Springer, 1972.

\bibitem{FultonHarris1991}
W.~Fulton and J.~Harris, \textit{Representation Theory: A First Course},
Springer, 1991.

\bibitem{Witten1982}
E.~Witten, ``An SU(2) Anomaly,'' \textit{Phys.\ Lett.\ B} \textbf{117},
324--328 (1982).

\bibitem{PeskinSchroeder1995}
M.\,E.~Peskin and D.\,V.~Schroeder, \textit{An Introduction to Quantum
Field Theory}, Westview Press, 1995.

\bibitem{GrossWilczek1973}
D.\,J.~Gross and F.~Wilczek, ``Ultraviolet Behavior of Non-Abelian
Gauge Theories,'' \textit{Phys.\ Rev.\ Lett.} \textbf{30}, 1343 (1973).

\bibitem{Politzer1973}
H.\,D.~Politzer, ``Reliable Perturbative Results for Strong Interactions?''
\textit{Phys.\ Rev.\ Lett.} \textbf{30}, 1346 (1973).

\bibitem{Planck2020}
Planck Collaboration, ``Planck 2018 results.\ VI.\ Cosmological
parameters,'' \textit{Astron.\ Astrophys.} \textbf{641}, A6 (2020).

\bibitem{PDG2024}
Particle Data Group, ``Review of Particle Physics,''
\textit{Phys.\ Rev.\ D} \textbf{110}, 030001 (2024).

\bibitem{Hardy2001}
L.~Hardy, ``Quantum theory from five reasonable axioms,''
arXiv:quant-ph/0101012 (2001).

\bibitem{CDP2011}
G.~Chiribella, G.\,M.~D'Ariano, and P.~Perinotti,
``Informational derivation of quantum theory,''
\textit{Phys.\ Rev.\ A} \textbf{84}, 012311 (2011).

\bibitem{Hohn2017}
P.\,A.~H\"ohn, ``Quantum theory from questions,''
\textit{Phys.\ Rev.\ A} \textbf{95}, 012102 (2017).

\bibitem{Connes1996}
A.~Connes, ``Gravity coupled with matter and the foundation of
noncommutative geometry,''
\textit{Commun.\ Math.\ Phys.} \textbf{182}, 155--176 (1996).

\bibitem{Bousso1999}
R.~Bousso, ``A covariant entropy conjecture,''
\textit{JHEP} \textbf{07}, 004 (1999).

\bibitem{HKM1976}
S.\,W.~Hawking, A.\,R.~King, and P.\,J.~McCarthy,
``A new topology for curved space-time which incorporates the causal,
differential, and conformal structures,''
\textit{J.\ Math.\ Phys.} \textbf{17}, 174--181 (1976).

\bibitem{Malament1977}
D.\,B.~Malament, ``The class of continuous timelike curves determines
the topology of spacetime,''
\textit{J.\ Math.\ Phys.} \textbf{18}, 1399--1404 (1977).

\bibitem{Kolmogorov1933}
A.\,N.~Kolmogorov, \textit{Grundbegriffe der Wahrscheinlichkeitsrechnung},
Springer, 1933.

\bibitem{Lovelock1971}
D.~Lovelock, ``The Einstein tensor and its generalizations,''
\textit{J.\ Math.\ Phys.} \textbf{12}, 498--501 (1971).

\bibitem{BrookeFermionScan}
E.\,S.~Brooke, ``APF Fermion Template Scan---Standalone Verification
Script (v2),'' 2026.
Zenodo: \url{https://doi.org/10.5281/zenodo.19154197}.
GitHub: \url{https://github.com/Ethan-Brooke/APF-Paper-2-The-Structure-of-Admissible-Physics}.
Interactive walkthrough:
\url{https://colab.research.google.com/drive/1otO1Qwp_Lr3OsK6ykkeMRATanecPkKg6}.

\bibitem{DonnellyGiddings2016PRD93}
W.~Donnelly and S.\,B.~Giddings, ``Diffeomorphism-invariant observables
and their nonlocal algebra,'' \textit{Phys. Rev. D} \textbf{93}, 024030
(2016).

\bibitem{DonnellyGiddings2016PRD94}
W.~Donnelly and S.\,B.~Giddings, ``Observables, gravitational dressing,
and obstructions to locality and subsystems,'' \textit{Phys. Rev. D}
\textbf{94}, 104038 (2016).

\bibitem{SanchezQuevedo2023}
L.\,M.~S\'anchez and H.~Quevedo, ``Thermodynamics of the FLRW apparent
horizon,'' \textit{Phys. Lett. B} \textbf{839}, 137778 (2023).

\end{thebibliography}


\end{document}
